% Courbes représentatives — Mathématiques Tle A — Cours signé OnBuch+
% Programme officiel MINESEC. Compiler avec : tectonic MA04-courbes-representatives-a.tex
\newcommand{\DOCMATIERE}{Mathématiques}
\newcommand{\DOCNIVEAU}{Tle A}
\newcommand{\DOCMODULE}{Relations et opérations fondamentales dans l'ensemble des nombres réels}
\newcommand{\DOCLECON}{A04}
\newcommand{\DOCTITRE}{Courbes représentatives}
% OnBuch+ — préambule « cours signés » ENRICHI (compilé avec tectonic / XeLaTeX)
% Système visuel « Pop » d'OnBuch+ : fond crème (jamais blanc), bordures encre
% épaisses, ombres « dures » décalées, titres Archivo Black en capitales.
% Version MATHÉMATIQUES : graphes (pgfplots/tikz), boîtes pédagogiques
% (définition, propriété, méthode, attention, le savais-tu, exercice résolu,
% exercice, TP, bilan), page de garde et signature OnBuch+ ancrée en bas.
%
% Macros attendues dans main.tex AVANT \input{preamble} :
%   \DOCMATIERE  \DOCNIVEAU  \DOCMODULE  \DOCLECON (numéro)  \DOCTITRE
\documentclass[11pt]{article}

\usepackage{fontspec}
\usepackage[french]{babel}
\usepackage[a4paper,margin=2cm,top=3.4cm,bottom=2.8cm,headheight=1.4cm,footskip=1.3cm]{geometry}
\usepackage{amsmath,amssymb,amsfonts}
\usepackage{mathtools} % \xmapsto, \coloneqq... souvent utilisés par les modèles
\usepackage{cancel} % simplifications barrées
% \abs{z} (module, valeur absolue) et \norm{u} (norme), utilisés par les modèles
\providecommand{\abs}{}\let\abs\relax\DeclarePairedDelimiter{\abs}{\lvert}{\rvert}
\providecommand{\norm}{}\let\norm\relax\DeclarePairedDelimiter{\norm}{\lVert}{\rVert}
\usepackage[table]{xcolor} % [table] : fournit aussi \rowcolors (alternance de lignes)
\usepackage{pagecolor}
\usepackage{tikz}
\usetikzlibrary{arrows.meta,positioning,calc,shapes.geometric,shapes.misc,shapes.arrows,decorations.pathmorphing,decorations.markings,patterns,shadows,mindmap,trees,fit,backgrounds,matrix,intersections}
\usepackage{pgfplots}
\usepackage{pgf-pie} % diagrammes circulaires \pie (statistiques)
\pgfplotsset{compat=1.18}
\usepgfplotslibrary{fillbetween,groupplots} % aires entre courbes (calcul intégral)
\usepackage{enumitem}
\usepackage{fancyhdr}
% [most] charge toutes les bibliothèques tcolorbox (listings, minted, raster,
% documentation...) et gonfle inutilement la mémoire TeX sur un document long
% (~300 boîtes) : on ne charge que ce qui est réellement utilisé.
\usepackage[skins,breakable]{tcolorbox}
\usepackage{graphicx}
\usepackage{colortbl}
\usepackage{booktabs}
\usepackage{tabularx}
\usepackage{diagbox} % case d'en-tête diagonale des tableaux à double entrée
% Colonnes extensibles centrée / gauche / droite (C, L, R), souvent utilisées par les modèles
\newcolumntype{C}{>{\centering\arraybackslash}X}
\newcolumntype{L}{>{\raggedright\arraybackslash}X}
\newcolumntype{R}{>{\raggedleft\arraybackslash}X}
\usepackage{array}
\usepackage{multicol}
\usepackage{siunitx}
% parse-numbers=false : accepte aussi \SI{2,5 \times 10^{-3}}{...}
\sisetup{locale=FR,output-decimal-marker={,},group-minimum-digits=4,parse-numbers=false}
\DeclareSIUnit\ohm{\ensuremath{\symup{\Omega}}} % U+2126 absent de la police
\usepackage{qrcode}
% \degree : commande fréquemment utilisée par les modèles pour un angle ou une
% température en dehors de \SI{}{} (non fournie par les paquets chargés).
\providecommand{\degree}{^{\circ}}
% \qed (fin de démonstration), utilisé par les modèles sans amsthm
\providecommand{\qed}{\hfill$\square$}
% \barre : double barre des valeurs interdites dans un tableau de variations (tkz-tab)
\providecommand{\barre}{\ensuremath{\|}}
% environnement proof (amsthm non chargé) utilisé par les modèles
\ifdefined\proof\else\newenvironment{proof}[1][Démonstration]{\par\noindent\textit{#1.}\ }{\qed\par}\fi
\usepackage{lastpage}
\usepackage{hyperref}
\hypersetup{hidelinks}

% ============================================================
% Palette « Pop » OnBuch+ — mêmes hex que la classe P (plus_ui.dart)
% ============================================================
\definecolor{poporange}{HTML}{F59321}
\definecolor{popdark}{HTML}{E07A0C}
\definecolor{popcream}{HTML}{FFF6E8}
\definecolor{popink}{HTML}{1C1714}
\definecolor{poppurple}{HTML}{7A5AE0}
\definecolor{popgreen}{HTML}{2E9E6A}
\definecolor{poppink}{HTML}{E0457B}
\definecolor{popblue}{HTML}{2F7DE1}
\definecolor{popgold}{HTML}{C98A12}
\definecolor{popmuted}{HTML}{8A7F76}
\definecolor{popline}{HTML}{EADFD2}
\definecolor{popgray}{HTML}{8A7F76}
\colorlet{popgrey}{popgray}
% Alias tolérants (le modèle nomme parfois les couleurs différemment)
\colorlet{popwhite}{white}
\colorlet{popblack}{popink}
\colorlet{popred}{poppink}
\colorlet{popyellow}{popgold}
% Teintes claires pour fonds de tableaux / zones
\colorlet{popblueL}{popblue!10!white}
\colorlet{poporangeL}{poporange!14!white}
\colorlet{popgreenL}{popgreen!12!white}
\colorlet{poppurpleL}{poppurple!12!white}
\colorlet{poppinkL}{poppink!12!white}
\colorlet{popgoldL}{popgold!14!white}

\pagecolor{popcream}

% ============================================================
% Typographie
% ============================================================
\defaultfontfeatures{Ligatures=TeX}
\setmainfont{PlusJakartaSans-Medium.ttf}[Path=../../../fonts/,BoldFont=PlusJakartaSans-Bold.ttf]
% Maths en sans-serif (Fira Math) avec chiffres et lettres droites de Plus
% Jakarta, pour que les formules se fondent dans le texte.
\usepackage[math-style=ISO,bold-style=ISO]{unicode-math}
\setmathfont{FiraMath-Regular.otf}[Path=../../../fonts/]
\setmathfont{PlusJakartaSans-Medium.ttf}[Path=../../../fonts/,range={up/num,up/{Latin,latin},\mathup}]
\usepackage{icomma} % virgule décimale sans espace en mode math
% Caractères mathématiques Unicode absents des polices de texte (Plus Jakarta,
% Archivo Black) : rendus par leur équivalent mathématique, en texte comme en
% formule — sinon ils s'affichent en carré vide (ex. « Arithmétique dans ℤ »).
\usepackage{newunicodechar}
\newunicodechar{ℝ}{\ensuremath{\mathbb{R}}}
\newunicodechar{ℤ}{\ensuremath{\mathbb{Z}}}
\newunicodechar{ℂ}{\ensuremath{\mathbb{C}}}
\newunicodechar{ℕ}{\ensuremath{\mathbb{N}}}
\newunicodechar{ℚ}{\ensuremath{\mathbb{Q}}}
\newunicodechar{𝔻}{\ensuremath{\mathbb{D}}}
\newunicodechar{α}{\ensuremath{\alpha}}
\newunicodechar{β}{\ensuremath{\beta}}
\newunicodechar{γ}{\ensuremath{\gamma}}
\newunicodechar{δ}{\ensuremath{\delta}}
\newunicodechar{ε}{\ensuremath{\varepsilon}}
\newunicodechar{θ}{\ensuremath{\theta}}
\newunicodechar{λ}{\ensuremath{\lambda}}
\newunicodechar{μ}{\ensuremath{\mu}}
\newunicodechar{σ}{\ensuremath{\sigma}}
\newunicodechar{τ}{\ensuremath{\tau}}
\newunicodechar{φ}{\ensuremath{\varphi}}
\newunicodechar{ψ}{\ensuremath{\psi}}
\newunicodechar{ω}{\ensuremath{\omega}}
\newunicodechar{Σ}{\ensuremath{\Sigma}}
% majuscules produites par \MakeUppercase dans les titres (α-aminés -> Α)
\newunicodechar{Α}{\ensuremath{\alpha}}
\newunicodechar{Β}{\ensuremath{\beta}}
\newunicodechar{Γ}{\ensuremath{\Gamma}}
\newunicodechar{Δ}{\ensuremath{\Delta}}
\newunicodechar{Ε}{\ensuremath{\varepsilon}}
\newunicodechar{Θ}{\ensuremath{\Theta}}
\newunicodechar{Λ}{\ensuremath{\Lambda}}
\newunicodechar{Μ}{\ensuremath{\mu}}
\newunicodechar{Τ}{\ensuremath{\tau}}
\newunicodechar{Φ}{\ensuremath{\Phi}}
\newunicodechar{Ψ}{\ensuremath{\Psi}}
\newunicodechar{Ω}{\ensuremath{\Omega}}
\newunicodechar{↦}{\ensuremath{\mapsto}}
\newunicodechar{∈}{\ensuremath{\in}}
\newunicodechar{∉}{\ensuremath{\notin}}
\newunicodechar{∧}{\ensuremath{\wedge}}
\newunicodechar{⇔}{\ensuremath{\Leftrightarrow}}
\newunicodechar{⟺}{\ensuremath{\Longleftrightarrow}}
\newunicodechar{⇒}{\ensuremath{\Rightarrow}}
\newunicodechar{⟹}{\ensuremath{\Longrightarrow}}
\newunicodechar{∘}{\ensuremath{\circ}}
\newunicodechar{∪}{\ensuremath{\cup}}
\newunicodechar{∩}{\ensuremath{\cap}}
\newunicodechar{≡}{\ensuremath{\equiv}}
\newunicodechar{⊂}{\ensuremath{\subset}}
\newunicodechar{⊥}{\ensuremath{\perp}}
\newunicodechar{∃}{\ensuremath{\exists}}
\newunicodechar{∀}{\ensuremath{\forall}}
\newunicodechar{∛}{\ensuremath{\sqrt[3]{\,}}}
\newunicodechar{∜}{\ensuremath{\sqrt[4]{\,}}}
\newunicodechar{⃗}{\ensuremath{{}^{\to}}}
\newunicodechar{ⁿ}{\ifmmode^{n}\else\textsuperscript{\ensuremath{n}}\fi}
\newunicodechar{ˣ}{\ifmmode^{x}\else\textsuperscript{\ensuremath{x}}\fi}
\newunicodechar{ᵇ}{\ifmmode^{b}\else\textsuperscript{\ensuremath{b}}\fi}
\newunicodechar{ʳ}{\ifmmode^{r}\else\textsuperscript{\ensuremath{r}}\fi}
\newunicodechar{ᵘ}{\ifmmode^{u}\else\textsuperscript{\ensuremath{u}}\fi}
\newunicodechar{ʸ}{\ifmmode^{y}\else\textsuperscript{\ensuremath{y}}\fi}
\newunicodechar{ᵃ}{\ifmmode^{a}\else\textsuperscript{\ensuremath{a}}\fi}
\newunicodechar{ᵉ}{\ifmmode^{e}\else\textsuperscript{\ensuremath{e}}\fi}
\newunicodechar{ᵏ}{\ifmmode^{k}\else\textsuperscript{\ensuremath{k}}\fi}
\newunicodechar{ᵖ}{\ifmmode^{p}\else\textsuperscript{\ensuremath{p}}\fi}
\newunicodechar{ᴺ}{\ifmmode^{N}\else\textsuperscript{\ensuremath{N}}\fi}
\newunicodechar{ᶜ}{\ifmmode^{c}\else\textsuperscript{\ensuremath{c}}\fi}
\newunicodechar{ʰ}{\ifmmode^{h}\else\textsuperscript{\ensuremath{h}}\fi}
\newunicodechar{ᵠ}{\ifmmode^{q}\else\textsuperscript{\ensuremath{q}}\fi}
\newunicodechar{ₙ}{\ifmmode_{n}\else\textsubscript{\ensuremath{n}}\fi}
\newunicodechar{ₐ}{\ifmmode_{a}\else\textsubscript{\ensuremath{a}}\fi}
\newunicodechar{ᵢ}{\ifmmode_{i}\else\textsubscript{\ensuremath{i}}\fi}
\newunicodechar{ᵣ}{\ifmmode_{r}\else\textsubscript{\ensuremath{r}}\fi}
\newunicodechar{ₖ}{\ifmmode_{k}\else\textsubscript{\ensuremath{k}}\fi}
\newunicodechar{ᵦ}{\ifmmode_{\beta}\else\textsubscript{\ensuremath{\beta}}\fi}
\newunicodechar{ₓ}{\ifmmode_{x}\else\textsubscript{\ensuremath{x}}\fi}
\newfontfamily\popdisplay{ArchivoBlack-Regular.ttf}[Path=../../../fonts/]
\newfontfamily\poplogo{SpaceGrotesk-Bold.ttf}[Path=../../../fonts/]
\newfontfamily\popsemi{PlusJakartaSans-SemiBold.ttf}[Path=../../../fonts/]
\newfontfamily\popxbold{PlusJakartaSans-ExtraBold.ttf}[Path=../../../fonts/]
\newfontfamily\popxboldls{PlusJakartaSans-ExtraBold.ttf}[Path=../../../fonts/,LetterSpace=3.0]

\setlist[enumerate]{leftmargin=1.7em,itemsep=5pt,topsep=4pt}
\setlist[itemize]{leftmargin=1.7em,itemsep=4pt,topsep=4pt,label={\color{poporange}\textbullet}}
\setlength{\parindent}{0pt}
\setlength{\parskip}{5pt}
\renewcommand{\arraystretch}{1.25}

% pgfplots : style Pop par défaut
\pgfplotsset{
  every axis/.append style={axis line style={popink,line width=1pt},tick style={popink},
    grid style={popline,line width=0.5pt},label style={font=\small},tick label style={font=\footnotesize}},
  popcurve/.style={poporange,line width=1.8pt,no markers,smooth},
}
% Même style disponible dans un \draw TikZ simple (hors pgfplots)
\tikzset{popcurve/.style={poporange,line width=1.8pt}}

% ============================================================
% Pill / squiggle
% ============================================================
\newcommand{\poppill}[3]{%
  \tikz[baseline=-0.55ex]\node[draw=popink,line width=1.2pt,rounded corners=2.6pt,%
    fill=#2,inner xsep=6.5pt,inner ysep=3pt]{%
    \popxboldls\fontsize{7}{7}\selectfont\color{#3}\MakeUppercase{#1}};%
}
\newcommand{\popsquiggle}[1]{%
  \tikz[baseline=-0.4ex]{
    \draw[#1,line width=2.6pt,line cap=round]
      plot [smooth] coordinates {(0,0.09) (0.34,0.01) (0.68,0.21) (1.02,0.01) (1.36,0.10)};
  }%
}
% Mot-clé surligné (marqueur orange sous le texte)
\newcommand{\cle}[1]{{\popxbold\color{popdark}#1}}

% ============================================================
% En-tête / pied
% ============================================================
\pagestyle{fancy}
\fancyhf{}
\renewcommand{\headrulewidth}{0pt}
\renewcommand{\footrulewidth}{0pt}
\lhead{\raisebox{-0.15ex}{\poplogo\fontsize{13}{13}\selectfont\color{popink}OnBuch\color{poporange}+}}
\rhead{\poppill{\DOCMATIERE\ \textbullet\ \DOCNIVEAU\ \textbullet\ Leçon \DOCLECON}{white}{popink}}
\lfoot{\popxbold\fontsize{7.6}{7.6}\selectfont\color{popmuted}\MakeUppercase{Cours signé OnBuch+}\ \textbullet\ {\popsemi\DOCTITRE}}
\rfoot{\tikz[baseline=-0.6ex]\node[rounded corners=8.5pt,fill=popink,minimum height=17pt,minimum width=17pt,inner xsep=5pt,inner sep=0]{\popxbold\fontsize{8}{8}\selectfont\color{popcream}\thepage\,/\,\pageref*{LastPage}};}

% ============================================================
% Titres
% ============================================================
\newcommand{\chaptitre}[2]{%
  \begin{tcolorbox}[enhanced,colback=poporange,colframe=popink,boxrule=2.6pt,arc=11pt,
      drop shadow=popink,left=15pt,right=15pt,top=12pt,bottom=11pt,before skip=3pt,after skip=18pt]
    \poppill{#2}{popink}{popcream}\par\vspace{9pt}
    {\raggedright\hyphenpenalty=10000\exhyphenpenalty=10000\popdisplay\fontsize{22}{24}\selectfont\color{popcream}\MakeUppercase{#1}\par}
  \end{tcolorbox}
}
% Section numérotée : pastille orange avec le numéro + titre Archivo
\newcounter{popsec}
\newcommand{\coursec}[1]{%
  \stepcounter{popsec}\setcounter{popsub}{0}%
  \par\vspace{16pt}\noindent
  \begin{minipage}{\linewidth}
  \tikz[baseline=-0.7ex]\node[draw=popink,line width=1.6pt,fill=poporange,rounded corners=4pt,minimum size=22pt,inner sep=0]{\popdisplay\fontsize{12}{12}\selectfont\color{popcream}\thepopsec};%
  \hspace{8pt}{\popdisplay\fontsize{15}{17}\selectfont\color{popink}\MakeUppercase{#1}}\par
  \vspace{4pt}\noindent{\color{poporange}\rule{36pt}{3.6pt}}
  \end{minipage}\par\nopagebreak\vspace{8pt}
}
\newcounter{popsub}
\newcommand{\courssub}[1]{%
  \stepcounter{popsub}%
  \par\vspace{9pt}\noindent{\popxbold\fontsize{11.5}{13}\selectfont\color{popink}{\color{poporange}\thepopsec.\thepopsub}\hspace{6pt}#1}\par\nopagebreak\vspace{4pt}
}

% ============================================================
% Boîtes pédagogiques — même recette : fond blanc, bordure épaisse colorée,
% ombre dure encre, badge plein. Titre optionnel [#1].
% ============================================================
\tcbset{popbase/.style={breakable,enhanced,colback=white,boxrule=2.3pt,arc=9pt,
  shadow={3pt}{-3pt}{0pt}{popink},left=14pt,right=14pt,top=11pt,bottom=11pt,before skip=11pt,after skip=15pt,
  fontupper=\popsemi\fontsize{10.6}{14.5}\selectfont\color{popink},
  fontlower=\popsemi\fontsize{10.6}{14.5}\selectfont\color{popink}}}
% Coche dessinée (le glyphe ✓ n'existe pas dans les polices du document)
\newcommand{\popcheck}[1]{\tikz[baseline=-0.5ex]\draw[#1,line width=1.6pt,line cap=round,line join=round](-0.13,0.0)--(-0.04,-0.09)--(0.13,0.1);}
\newcommand{\popboxhead}[3]{\poppill{#1}{#2}{white}\ifx\relax#3\relax\else\hspace{6pt}{\popxbold\fontsize{10.6}{12}\selectfont\color{popink}#3}\fi\par\vspace{7pt}}

\newtcolorbox{aretenir}[1][]{popbase,colframe=popgreen,before upper={\popboxhead{À retenir}{popgreen}{#1}}}
\newtcolorbox{exemplebox}[1][]{popbase,colframe=poporange,before upper={\popboxhead{Exemple}{poporange}{#1}}}
\newtcolorbox{definition}[1][]{popbase,colframe=popblue,before upper={\popboxhead{Définition}{popblue}{#1}}}
\newtcolorbox{propriete}[1][]{popbase,colframe=poppurple,before upper={\popboxhead{Propriété}{poppurple}{#1}}}
\newtcolorbox{methode}[1][]{popbase,colframe=popgold,before upper={\popboxhead{Méthode}{popgold}{#1}}}
\newtcolorbox{attention}[1][]{popbase,colframe=poppink,before upper={\popboxhead{Attention}{poppink}{#1}}}
\newtcolorbox{experience}[1][]{popbase,colframe=popblue,colback=popblueL,before upper={\popboxhead{Expérience}{popblue}{#1}}}
% « Le savais-tu ? » : carte violette pleine (contexte, culture, Cameroun)
\newtcolorbox{savaistu}[1][]{popbase,colback=poppurple,colframe=popink,
  fontupper=\popsemi\fontsize{10.4}{14.2}\selectfont\color{white},
  before upper={\poppill{Le savais-tu\,?}{white}{poppurple}\ifx\relax#1\relax\else\hspace{6pt}{\popxbold\color{white}#1}\fi\par\vspace{7pt}}}
% Exercice résolu : énoncé en haut, \tcblower puis la solution sur fond crème
\newcounter{popexo}\newcounter{popexr}
\newtcolorbox{exoresolu}[1][]{popbase,colframe=poporange,colbacklower=popcream,
  segmentation style={popink,line width=1.2pt,dashed},
  before upper={\stepcounter{popexr}\popboxhead{Exercice résolu \thepopexr}{poporange}{#1}},
  before lower={\poppill{Solution}{popink}{popcream}\par\vspace{6pt}}}
% Exercice (non résolu en place) — la correction est dans la partie Corrigés
\newtcolorbox{exercice}[1][]{popbase,colframe=popink,
  before upper={\stepcounter{popexo}\popboxhead{Exercice \thepopexo}{popink}{#1}}}
% Corrigé
\newtcolorbox{corrige}[1][]{popbase,colframe=popgreen,colback=popgreenL,
  before upper={\popboxhead{Corrigé}{popgreen}{#1}}}
% Objectifs / prérequis / bilan
\newtcolorbox{objectifs}{popbase,colframe=popink,colback=poporangeL,
  before upper={\popboxhead{Objectifs de la leçon}{popink}{}}}
\newtcolorbox{prerequis}{popbase,colframe=popink,colback=popblueL,
  before upper={\popboxhead{Prérequis}{popblue}{}}}
\newtcolorbox{bilan}{popbase,colframe=popink,colback=white,boxrule=2.8pt,
  before upper={\popboxhead{Fiche bilan}{poporange}{L'essentiel à connaître pour l'examen}}}
% Figure encadrée avec légende
\newtcolorbox{popfigure}[1][]{enhanced,breakable=false,colback=white,colframe=popline,boxrule=1.4pt,arc=8pt,
  left=10pt,right=10pt,top=10pt,bottom=8pt,before skip=10pt,after skip=12pt,halign=center,
  lower separated=false,halign lower=center,
  fontlower=\popsemi\fontsize{9}{11.5}\selectfont\color{popmuted},
  #1}
\newcommand{\legende}[1]{\par\vspace{6pt}{\popsemi\fontsize{9}{11.5}\selectfont\color{popmuted}#1}}

% ============================================================
% Signature OnBuch+
% ============================================================
\newcommand{\poplogoinline}[1]{{\poplogo\fontsize{#1}{#1}\selectfont\color{popink}OnBuch\color{poporange}+}}
\newcommand{\signatureonbuch}{%
  \begin{tcolorbox}[enhanced,colback=popink,colframe=popink,boxrule=2.2pt,arc=10pt,
      drop shadow=poporange,left=14pt,right=14pt,top=10pt,bottom=10pt,before skip=0pt,after skip=0pt]
    \tikz[baseline=-0.7ex]\node[circle,fill=popgreen,draw=popcream,line width=1.3pt,minimum size=22pt,inner sep=0]{\popcheck{white}};%
    \hspace{9pt}\begin{minipage}[c]{0.62\linewidth}
      {\popxbold\fontsize{12}{13}\selectfont\color{popcream}Signé\ \textbullet\ {\poplogo OnBuch\color{poporange}+}}\par\vspace{2pt}
      {\raggedright\popsemi\fontsize{8}{10.5}\selectfont\color{popline}\MakeUppercase{Équipe pédagogique OnBuch+}\\\MakeUppercase{Conforme au programme officiel MINESEC}\par}
    \end{minipage}\hfill
    {\poplogo\fontsize{22}{22}\selectfont\color{popcream}OnBuch\color{poporange}+}
  \end{tcolorbox}%
}

% ============================================================
% Page de garde
% #1 titre · #2 matière · #3 niveau · #4 module · #5 n° leçon · #6 durée ·
% #7 description / objectifs (texte ou itemize)
% ============================================================
\newcommand{\pagegarde}[7]{%
  \thispagestyle{empty}
  \newgeometry{a4paper,margin=1.7cm,top=1.6cm,bottom=1.4cm}%
  \begin{tikzpicture}[remember picture,overlay]
    \fill[poporange!18] ([xshift=-3.2cm,yshift=-3.6cm]current page.north east) circle (4.6cm);
    \fill[poppurple!14] ([xshift=2.4cm,yshift=7.6cm]current page.south west) circle (3.8cm);
  \end{tikzpicture}%
  {\poplogo\fontsize{30}{30}\selectfont\color{popink}OnBuch\color{poporange}+}\hfill
  \raisebox{0.6ex}{\poppill{Cours signé \textbullet\ Premium}{popink}{popcream}}\par
  \vspace{1pt}\popsquiggle{poppurple}\par
  \vspace{8pt}
  \begin{tcolorbox}[enhanced,colback=poporange,colframe=popink,boxrule=2.8pt,arc=13pt,
      drop shadow=popink,left=18pt,right=18pt,top=12pt,bottom=13pt,after skip=10pt]
    \poppill{#2 \textbullet\ #3}{popink}{popcream}\hspace{5pt}\poppill{#4}{white}{popink}\par\vspace{8pt}
    {\popdisplay\fontsize{14}{14}\selectfont\color{popink}LEÇON #5}\par\vspace{4pt}
    {\raggedright\hyphenpenalty=10000\exhyphenpenalty=10000\popdisplay\fontsize{25}{27}\selectfont\color{popcream}\MakeUppercase{#1}\par}
    \vspace{8pt}
    {\popxbold\fontsize{9.5}{9.5}\selectfont\color{popink}\MakeUppercase{Durée conseillée : #6}}
  \end{tcolorbox}
  \begin{tcolorbox}[enhanced,colback=white,colframe=popink,boxrule=2.2pt,arc=10pt,
      drop shadow=popink,left=16pt,right=16pt,top=10pt,bottom=10pt,after skip=8pt]
    \poppill{Dans cette leçon}{poppurple}{white}\par\vspace{5pt}
    \popsemi\fontsize{9.3}{12.6}\selectfont\color{popink}#7
  \end{tcolorbox}
  \noindent\begin{minipage}[t]{0.487\linewidth}
    \begin{tcolorbox}[enhanced,colback=popgreen,colframe=popink,boxrule=2.2pt,arc=10pt,
        drop shadow=popink,left=11pt,right=11pt,top=10pt,bottom=10pt,height=3.1cm,valign=center]
      \tcbox[colback=white,colframe=popink,boxrule=1.3pt,arc=5pt,boxsep=0pt,nobeforeafter,
        left=3pt,right=3pt,top=3pt,bottom=3pt,valign=center]{\qrcode[height=1.9cm]{https://play.google.com/store/apps/details?id=com.luvvix.onbuch&hl=fr}}%
      \hspace{8pt}\begin{minipage}[c]{0.5\linewidth}
        {\popdisplay\fontsize{11}{12.5}\selectfont\color{white}\MakeUppercase{Télécharge OnBuch}}\par\vspace{4pt}
        {\raggedright\popsemi\fontsize{8.4}{11}\selectfont\color{white}Gratuit \textbullet\ Google Play\\Tuteur IA, annales, concours, résultats\par}
      \end{minipage}
    \end{tcolorbox}
  \end{minipage}\hfill
  \begin{minipage}[t]{0.487\linewidth}
    \begin{tcolorbox}[enhanced,colback=poppurple,colframe=popink,boxrule=2.2pt,arc=10pt,
        drop shadow=popink,left=11pt,right=11pt,top=10pt,bottom=10pt,height=3.1cm,valign=center]
      \tikz[baseline=-0.7ex]\node[circle,fill=white,draw=popink,line width=1.3pt,minimum size=1.6cm,inner sep=0]{\popdisplay\fontsize{24}{24}\selectfont\color{poppurple}+};%
      \hspace{8pt}\begin{minipage}[c]{0.6\linewidth}
        {\popdisplay\fontsize{11}{12.5}\selectfont\color{white}\MakeUppercase{Passe à OnBuch+}}\par\vspace{4pt}
        {\raggedright\popsemi\fontsize{8.4}{11}\selectfont\color{white}Tous les cours signés, Tuteur IA illimité, fiches PDF — onglet OnBuch+ de l'app\par}
      \end{minipage}
    \end{tcolorbox}
  \end{minipage}
  \vspace{8pt}
  \popcontenu
  \vfill
  \signatureonbuch
  \restoregeometry
  \newpage
}

% Rangée « Ce cours contient » (page de garde)
\newcommand{\poptile}[3]{% #1 couleur #2 gros texte #3 légende
  \begin{tcolorbox}[enhanced,colback=#1,colframe=popink,boxrule=2pt,arc=9pt,shadow={2.5pt}{-2.5pt}{0pt}{popink},
      left=6pt,right=6pt,top=9pt,bottom=9pt,halign=center,valign=center,height=1.75cm,width=\linewidth,nobeforeafter]
    {\popdisplay\fontsize{12.5}{13}\selectfont\color{white}\MakeUppercase{#2}}\par\vspace{3pt}
    {\centering\popsemi\fontsize{7.8}{9.5}\selectfont\color{white}#3\par}
  \end{tcolorbox}}
\newcommand{\popcontenu}{%
  \par\noindent{\popxboldls\fontsize{7.5}{7.5}\selectfont\color{popmuted}CE COURS SIGNÉ CONTIENT}\par\vspace{7pt}
  \noindent\begin{minipage}[t]{0.235\linewidth}\poptile{popblue}{Cours}{complet, illustré, pas à pas}\end{minipage}\hfill
  \begin{minipage}[t]{0.235\linewidth}\poptile{poporange}{Méthodes}{et exercices résolus}\end{minipage}\hfill
  \begin{minipage}[t]{0.235\linewidth}\poptile{poppink}{Exercices}{type examen + corrigés}\end{minipage}\hfill
  \begin{minipage}[t]{0.235\linewidth}\poptile{popgreen}{Fiche bilan}{l'essentiel à retenir}\end{minipage}\par}

% Sommaire
\newtcolorbox{sommaire}{enhanced,colback=white,colframe=popink,boxrule=2.2pt,arc=10pt,
  drop shadow=popink,left=18pt,right=18pt,top=13pt,bottom=13pt,before skip=6pt,after skip=20pt,
  before upper={\poppill{Sommaire}{poppurple}{white}\par\vspace{9pt}\popsemi\fontsize{10.6}{14.5}\selectfont\color{popink}}}

% Signature de fin de cours — poussée tout en bas de la dernière page
\newcommand{\findecours}{%
  \par\vfill\nopagebreak
  \begin{center}\popsquiggle{poporange}\end{center}\vspace{4pt}
  \signatureonbuch
}
\AtBeginDocument{\def\llbracket{\mathopen{[\mkern-3.5mu[}}\def\rrbracket{\mathclose{]\mkern-3.5mu]}}} % intervalles d'entiers (glyphe ⟦ absent de la police)
% Glyphes absents de Fira Math : redéfinis à partir de glyphes disponibles
\AtBeginDocument{%
  \def\setminus{\mathbin{\backslash}}\let\smallsetminus\setminus
  \def\vdots{\mathord{\vbox{\baselineskip3.2pt\lineskiplimit0pt\kern4pt\hbox{.}\hbox{.}\hbox{.}}}}%
  \def\ddots{\mathinner{\mkern1mu\raise6.5pt\hbox{.}\mkern2mu\raise3.6pt\hbox{.}\mkern2mu\raise0.7pt\hbox{.}\mkern1mu}}%
  \def\triangle{\mathord{\scalebox{0.9}{$\symup{\Delta}$}}}%
  \def\nabla{\mathord{\raisebox{\depth}{\scalebox{1}[-1]{$\symup{\Delta}$}}}}%
  \def\checkmark{\popcheck{popdark}}%
  \def\star{\mathbin{\ast}}\def\bigstar{\mathord{\ast}}%
  \def\diamond{\mathbin{\scalebox{0.6}{$\Diamond$}}}%
}

%% FIGFIX-BEGIN v5 — correctifs globaux des figures (docs/onbuch-generation/tools/figfix)
\usepackage{adjustbox}\usepackage{etoolbox}\tcbuselibrary{raster}
% toute figure TikZ/pgfplots plus large que la ligne est réduite à la largeur disponible
\BeforeBeginEnvironment{tikzpicture}{\begin{adjustbox}{max width=\linewidth}}
\AfterEndEnvironment{tikzpicture}{\end{adjustbox}}
% légendes pgfplots lisibles par-dessus les courbes
\pgfplotsset{every axis/.append style={legend style={fill=white,fill opacity=0.92,text opacity=1,draw=black!15,inner sep=3pt}}}
% étiquettes d'arêtes (syntaxe edge["..."]) avec fond blanc
\tikzset{every edge quotes/.append style={fill=white,inner sep=1.2pt}}
%% FIGFIX-END
\begin{document}

\pagegarde{Courbes représentatives}{Mathématiques}{Tle A}{Relations et opérations fondamentales dans l'ensemble des nombres réels}{A04}{5 h}{Cette leçon outille l'élève pour étudier les éléments remarquables d'une courbe représentative : axes et centres de symétrie, asymptotes verticales, horizontales et obliques justifiées par des calculs de limites. Elle développe aussi la capacité à lire une courbe, émettre des conjectures et les vérifier, compétence centrale à l'épreuve de mathématiques du Baccalauréat A.
\vspace{4pt}
\begin{multicols}{2}\small
\begin{itemize}
\item À la fin de cette leçon, je sais reconnaître graphiquement un axe ou un centre de symétrie d'une courbe
\item À la fin de cette leçon, je sais démontrer algébriquement que la droite x = a est axe de symétrie d'une courbe
\item À la fin de cette leçon, je sais démontrer algébriquement que le point (a ; b) est centre de symétrie d'une courbe
\item À la fin de cette leçon, je sais déterminer une asymptote verticale par un calcul de limite
\item À la fin de cette leçon, je sais déterminer une asymptote horizontale par un calcul de limite
\item À la fin de cette leçon, je sais déterminer une asymptote oblique en étudiant la limite de f(x) − (ax + b)
\end{itemize}\end{multicols}}

\begin{sommaire}
\begin{multicols}{2}
\begin{enumerate}
\item Lecture graphique et conjectures
\item Axe de symétrie d'une courbe
\item Centre de symétrie d'une courbe
\item Asymptotes verticales et horizontales
\item Asymptotes obliques
\item Synthèse : construire une courbe complète
\item Activité d'intégration
\item Méthodes et exercices résolus
\item Exercices
\item Corrigés des exercices
\item Fiche bilan
\end{enumerate}
\end{multicols}
\end{sommaire}

\begin{objectifs}
\begin{itemize}
\item À la fin de cette leçon, je sais reconnaître graphiquement un axe ou un centre de symétrie d'une courbe
\item À la fin de cette leçon, je sais démontrer algébriquement que la droite x = a est axe de symétrie d'une courbe
\item À la fin de cette leçon, je sais démontrer algébriquement que le point (a ; b) est centre de symétrie d'une courbe
\item À la fin de cette leçon, je sais déterminer une asymptote verticale par un calcul de limite
\item À la fin de cette leçon, je sais déterminer une asymptote horizontale par un calcul de limite
\item À la fin de cette leçon, je sais déterminer une asymptote oblique en étudiant la limite de f(x) − (ax + b)
\item À la fin de cette leçon, je sais tracer une courbe à partir de son tableau de variation et de ses asymptotes
\item À la fin de cette leçon, je sais émettre une conjecture à partir d'une représentation graphique et la vérifier par le calcul
\end{itemize}
\end{objectifs}

\begin{prerequis}
\begin{itemize}
\item Notion de limite d'une fonction en un point et en l'infini (début de Terminale)
\item Lecture d'un tableau de variation et dérivabilité (Première)
\item Symétrie centrale et symétrie axiale dans le plan (collège, rappel analytique)
\item Coordonnées du milieu d'un segment et transformations de points (Seconde/Première)
\item Fonctions de référence : inverse, rationnelles simples, racine carrée
\end{itemize}
\end{prerequis}

\newpage

\coursec{Lecture graphique et conjectures}

\emph{Situation d'accroche.} Mme Etoundi gère un cybercafé dans le quartier Nlongkak à Yaoundé. Chaque soir, elle note le nombre de clients connectés et le coût moyen de connexion par client (abonnement CAMTEL, électricité, amortissement du matériel, répartis entre les usagers). Elle constate un phénomène curieux : plus il y a de clients, plus le coût moyen par client diminue, mais il semble se rapprocher d'une valeur plancher — disons 100 FCFA — sans jamais l'atteindre. Son comptable lui trace la courbe du coût moyen en fonction du nombre de clients : la courbe semble \og coller \fg{} à une droite horizontale quand le nombre de clients devient grand, et \og exploser \fg{} près d'une droite verticale quand le nombre de clients approche zéro. De plus, la courbe des recettes hebdomadaires présente une étrange régularité autour d'un point : pliée mentalement autour de ce point, elle semble se superposer à elle-même.

\medskip
\textbf{Problématique.} Comment lire correctement une courbe ? Comment formuler précisément ce que l'on \emph{croit} voir (conjecture) ? Et surtout, comment \emph{prouver} ces observations par le calcul, car un dessin, aussi beau soit-il, ne constitue jamais une démonstration ? C'est tout l'objet de cette leçon : \cle{observer}, \cle{conjecturer}, puis \cle{démontrer} les propriétés d'une courbe représentative.

\courssub{Rappel : courbe représentative d'une fonction}

\begin{definition}[Courbe représentative]
Soit $f$ une fonction définie sur une partie $D_f$ de $\mathbb{R}$, et $(O ; \vec{i}, \vec{j})$ un repère du plan. La \cle{courbe représentative} de $f$, notée $\mathcal{C}_f$, est l'ensemble des points du plan de coordonnées $(x ; f(x))$ lorsque $x$ décrit $D_f$ :
\[ \mathcal{C}_f = \left\{ M\big(x ; f(x)\big) \; ; \; x \in D_f \right\}. \]
Le repère est dit \cle{orthogonal} lorsque les axes sont perpendiculaires, et \cle{orthonormal} (ou orthonormé) lorsque de plus les unités sur les deux axes sont égales ($\|\vec{i}\| = \|\vec{j}\|$).
\end{definition}

\begin{attention}[Repère et symétries]
Dans tout ce chapitre, on se placera dans un repère \textbf{orthonormal}. C'est essentiel : les propriétés de symétrie (axe vertical, centre de symétrie) ne se lisent correctement sur le dessin que si les deux axes portent la même unité. Une symétrie visible dans un repère orthonormal peut disparaître visuellement si l'on écrase l'un des axes.
\end{attention}

\courssub{Lire une courbe : les informations de base}

Face à une courbe donnée (au Baccalauréat, on te fournit souvent le graphique d'une fonction inconnue), il faut savoir en extraire méthodiquement quatre types d'informations.

\begin{enumerate}
\item \textbf{L'ensemble de définition.} Ce sont les abscisses $x$ pour lesquelles la courbe possède au moins un point. On le lit en \og projetant \fg{} la courbe sur l'axe des abscisses. Un trou dans la courbe (valeur interdite, asymptote verticale) se traduit par un trou dans l'ensemble de définition.
\item \textbf{Les variations.} On lit les intervalles où la courbe \og monte \fg{} (fonction croissante) et ceux où elle \og descend \fg{} (fonction décroissante), en parcourant la courbe de la gauche vers la droite.
\item \textbf{Les extrema.} Un \cle{maximum} local correspond à un point \og haut \fg{} de la courbe (sommet d'une colline) ; un \cle{minimum} local à un point \og bas \fg{} (creux). On donne toujours un extremum sous la forme : \og $f$ admet un maximum (resp. minimum) égal à $f(a)$, atteint en $x = a$ \fg{}.
\item \textbf{Le signe.} La courbe au-dessus de l'axe des abscisses correspond à $f(x) > 0$ ; en dessous, à $f(x) < 0$ ; les intersections avec l'axe donnent les solutions de $f(x) = 0$.
\end{enumerate}

\begin{methode}[Lire une courbe en trois temps]
\begin{enumerate}
\item \textbf{D'abord l'ensemble de définition} : quelles valeurs de $x$ sont concernées ? Y a-t-il des valeurs interdites, des \og trous \fg{}, des sauts ?
\item \textbf{Ensuite les variations et les extrema} : dresser mentalement le tableau de variations à partir de l'allure de la courbe.
\item \textbf{Enfin les comportements aux bornes} : que fait la courbe quand $x$ devient très grand (vers $+\infty$ ou $-\infty$) ? Que fait-elle près d'une valeur interdite ? Semble-t-elle se rapprocher d'une droite ?
\end{enumerate}
C'est ce troisième temps qui prépare la notion d'\cle{asymptote}, étudiée dans les sections suivantes.
\end{methode}

\courssub{Observer des régularités : vers la conjecture}

Une courbe bien tracée révèle souvent des \emph{régularités} que la simple lecture des variations ne suffit pas à décrire. Deux phénomènes doivent attirer ton œil :

\begin{itemize}
\item \textbf{La courbe semble se rapprocher d'une droite.} Par exemple, le coût moyen de Mme Etoundi se rapproche de la droite horizontale d'équation $y = 100$ sans jamais la toucher : on conjecturera l'existence d'une \cle{asymptote horizontale}. Près d'une valeur interdite $a$, la courbe peut \og exploser \fg{} vers $+\infty$ ou $-\infty$ le long de la droite verticale d'équation $x = a$ : on conjecturera une \cle{asymptote verticale}.
\item \textbf{La courbe semble symétrique.} Elle peut sembler invariante par pliage le long d'une droite verticale $x = a$ (\cle{axe de symétrie}), ou par demi-tour autour d'un point $(a ; b)$ (\cle{centre de symétrie}).
\end{itemize}

\begin{definition}[Conjecture]
Une \cle{conjecture} est une propriété que l'on \emph{pense} vraie à partir d'observations (graphiques, numériques, expérimentales), mais qui n'est pas encore démontrée. Une conjecture peut s'avérer vraie... ou fausse : seule une \cle{démonstration} tranche.
\end{definition}

\begin{attention}[Une observation graphique n'est pas une démonstration]
C'est l'erreur la plus fréquente — et la plus coûteuse — au Baccalauréat. Écrire \og on voit sur le graphique que la courbe est symétrique, donc... \fg{} ne rapporte \textbf{aucun point} de démonstration. Le graphique sert à \emph{conjecturer} ; le calcul sert à \emph{prouver}. Toute propriété observée (symétrie, asymptote, position relative) doit être vérifiée par un calcul rigoureux : calcul de $f(a+h)$ et $f(a-h)$ pour un axe, de $f(a+h) + f(a-h)$ pour un centre, ou calcul de limite pour une asymptote. Les sections suivantes te donneront tous ces outils.
\end{attention}

\begin{methode}[La démarche en trois temps : observer $\to$ conjecturer $\to$ vérifier]
\begin{enumerate}
\item \textbf{Observer} : tracer ou examiner la courbe avec soin, repérer les régularités (droite d'appui, symétrie apparente).
\item \textbf{Conjecturer} : formuler une phrase précise et mathématique : \og la droite d'équation $x = 2$ semble être axe de symétrie de $\mathcal{C}_f$ \fg{} ; \og la droite $y = 3$ semble être asymptote horizontale en $+\infty$ \fg{}.
\item \textbf{Vérifier par le calcul} : démontrer la propriété à l'aide des définitions (égalité entre expressions de $f$) ou des théorèmes sur les limites. Si le calcul confirme, la conjecture devient un \emph{théorème} ; s'il infirme, on abandonne la conjecture.
\end{enumerate}
\end{methode}

\courssub{Exemple guidé : la courbe de la fonction inverse}

Étudions ensemble la fonction $f : x \mapsto \dfrac{1}{x}$, la plus célèbre des fonctions de Terminale A, et menons la démarche complète.

\begin{popfigure}
\begin{center}
\begin{tikzpicture}
\begin{axis}[
  width=0.85\linewidth, height=7cm,
  axis lines=middle, xlabel=$x$, ylabel=$y$,
  xmin=-4.2, xmax=4.2, ymin=-4.2, ymax=4.2,
  xtick={-3,-2,-1,1,2,3}, ytick={-3,-2,-1,1,2,3},
  grid=both, grid style={popline!40},
  samples=200, clip=true]
\addplot[popcurve, domain=0.25:4] {1/x};
\addplot[popcurve, domain=-4:-0.25] {1/x};
\addplot[only marks, mark=*, poporange, mark size=2pt] coordinates {(0,0)};
\node[poporange, below right] at (axis cs:0.1,-0.4) {$O$};
\end{axis}
\end{tikzpicture}
\end{center}
\legende{La courbe de $f(x) = \dfrac{1}{x}$ (hyperbole) : deux branches symétriques par rapport à l'origine.}
\end{popfigure}

\textbf{Premier temps : observer.} La courbe possède deux \og branches \fg{} séparées : l'une pour $x > 0$, l'autre pour $x < 0$. Elle ne touche jamais les axes. Quand $x$ devient grand, la branche droite semble coller à l'axe des abscisses ; quand $x$ s'approche de $0$ par valeurs positives, elle monte très haut. Enfin, la courbe semble invariante par demi-tour autour de l'origine $O$.

\textbf{Deuxième temps : conjecturer.} On formule quatre conjectures précises :
\begin{enumerate}
\item[(C1)] L'ensemble de définition est $D_f = \mathbb{R}^* = ]-\infty ; 0[ \cup ]0 ; +\infty[$.
\item[(C2)] La droite d'équation $x = 0$ (l'axe des ordonnées) semble être asymptote verticale.
\item[(C3)] La droite d'équation $y = 0$ (l'axe des abscisses) semble être asymptote horizontale en $+\infty$ et en $-\infty$.
\item[(C4)] L'origine $O(0 ; 0)$ semble être centre de symétrie de la courbe.
\end{enumerate}

\textbf{Troisième temps : vérifier par le calcul.}
\begin{itemize}
\item (C1) : $f(x) = \dfrac{1}{x}$ existe si et seulement si $x \neq 0$. Conjecture confirmée.
\item (C4) : pour tout $x \neq 0$, l'opposé $-x$ est aussi dans $D_f$, et
\[ f(-x) = \frac{1}{-x} = -\frac{1}{x} = -f(x). \]
La fonction est \cle{impaire} : le point $(-x ; f(-x)) = (-x ; -f(x))$ est le symétrique de $(x ; f(x))$ par rapport à $O$. L'origine est bien centre de symétrie. Conjecture confirmée.
\item (C2) et (C3) : elles seront démontrées par des calculs de limites dans la section 4 — tu verras que $\lim\limits_{x \to 0^+} \dfrac{1}{x} = +\infty$ (asymptote verticale) et $\lim\limits_{x \to +\infty} \dfrac{1}{x} = 0$ (asymptote horizontale).
\end{itemize}

\courssub{Un deuxième exemple : conjecturer un axe de symétrie}

\begin{exemplebox}[La parabole de $f(x) = x^2 - 4x + 5$]
Traçons mentalement (ou à la calculatrice) la courbe de $f(x) = x^2 - 4x + 5$. Quelques valeurs :
\begin{center}
\begin{tabularx}{0.8\linewidth}{|l|X|X|X|X|X|}
\hline
\rowcolor{popblueL}
$x$ & $0$ & $1$ & $2$ & $3$ & $4$ \\
\hline
$f(x)$ & $5$ & $2$ & $1$ & $2$ & $5$ \\
\hline
\end{tabularx}
\end{center}
\textbf{Observation} : les valeurs se répètent de part et d'autre de $x = 2$ : $f(1) = f(3)$ et $f(0) = f(4)$. La courbe semble se replier sur elle-même le long de la droite verticale passant par $x = 2$.
\par\textbf{Conjecture} : la droite d'équation $x = 2$ semble être axe de symétrie de $\mathcal{C}_f$.
\par\textbf{Vérification} : un point d'abscisse $2 + h$ et un point d'abscisse $2 - h$ sont symétriques par rapport à la droite $x = 2$. Calculons, pour tout réel $h$ :
\begin{align*}
f(2 + h) &= (2+h)^2 - 4(2+h) + 5 = 4 + 4h + h^2 - 8 - 4h + 5 = h^2 + 1, \\
f(2 - h) &= (2-h)^2 - 4(2-h) + 5 = 4 - 4h + h^2 - 8 + 4h + 5 = h^2 + 1.
\end{align*}
On obtient $f(2+h) = f(2-h)$ pour tout réel $h$ : les deux points symétriques ont la même ordonnée. La conjecture est \textbf{démontrée} : la droite $x = 2$ est bien axe de symétrie de la parabole. La méthode générale (montrer que $f(a+h) = f(a-h)$) sera formalisée dans la section 2.
\end{exemplebox}

\courssub{Conjecturer un centre de symétrie sur une courbe cubique}

Voici la courbe d'une fonction polynôme de degré 3, $g(x) = x^3 - 3x^2 + 4x - 1$. Un point $\Omega$ y est marqué. Observe bien la régularité de la courbe autour de ce point.

\begin{popfigure}
\begin{center}
\begin{tikzpicture}
\begin{axis}[
  width=0.85\linewidth, height=7.5cm,
  axis lines=middle, xlabel=$x$, ylabel=$y$,
  xmin=-1.5, xmax=3.5, ymin=-3.5, ymax=4.5,
  xtick={-1,1,2,3}, ytick={-3,-2,-1,1,2,3,4},
  grid=both, grid style={popline!40},
  samples=200, clip=true]
\addplot[popcurve, domain=-1.1:3.1] {x^3 - 3*x^2 + 4*x - 1};
\addplot[only marks, mark=*, poporange, mark size=2.5pt] coordinates {(1,1)};
\addplot[only marks, mark=*, popgreen, mark size=2pt] coordinates {(0,-1) (2,3)};
\addplot[popgreen, dashed] coordinates {(0,-1) (2,3)};
\node[poporange, above left] at (axis cs:1,1) {$\Omega(1\,;\,1)$};
\node[popgreen, left] at (axis cs:0,-1) {$M$};
\node[popgreen, right] at (axis cs:2,3) {$M'$};
\end{axis}
\end{tikzpicture}
\end{center}
\legende{Courbe de $g(x) = x^3 - 3x^2 + 4x - 1$. Les points $M$ et $M'$ semblent symétriques par rapport à $\Omega$.}
\end{popfigure}

\textbf{Observation.} Le point $M(0 ; -1)$ et le point $M'(2 ; 3)$ sont sur la courbe, et $\Omega(1 ; 1)$ semble être le milieu de $[MM']$ : en effet, $\left(\dfrac{0+2}{2} ; \dfrac{-1+3}{2}\right) = (1 ; 1)$. Ce phénomène semble se reproduire pour tout couple de points symétriques.

\textbf{Conjecture.} Le point $\Omega(1 ; 1)$ semble être centre de symétrie de la courbe de $g$.

\textbf{Vérification (à titre d'entraînement).} Pour que $\Omega(1 ; 1)$ soit centre de symétrie, il faut et il suffit que pour tout réel $h$ : $g(1+h) + g(1-h) = 2 \times 1 = 2$. Calculons séparément, en développant avec soin :
\begin{align*}
g(1+h) &= (1+h)^3 - 3(1+h)^2 + 4(1+h) - 1 \\
&= \big(1 + 3h + 3h^2 + h^3\big) - 3\big(1 + 2h + h^2\big) + 4 + 4h - 1 \\
&= h^3 + h + 1, \\[4pt]
g(1-h) &= (1-h)^3 - 3(1-h)^2 + 4(1-h) - 1 \\
&= \big(1 - 3h + 3h^2 - h^3\big) - 3\big(1 - 2h + h^2\big) + 4 - 4h - 1 \\
&= -h^3 - h + 1.
\end{align*}
En additionnant : $g(1+h) + g(1-h) = \big(h^3 + h + 1\big) + \big(-h^3 - h + 1\big) = 2$ pour tout réel $h$. La conjecture est démontrée : $\Omega$ est bien centre de symétrie. La méthode générale sera présentée dans la section 3.

\begin{exoresolu}[Lecture complète d'une courbe]
On donne la courbe $\mathcal{C}$ d'une fonction $f$ définie sur $]-\infty ; 1[ \cup ]1 ; +\infty[$. On observe : la courbe monte sur $]-\infty ; -1]$, descend sur $[-1 ; 1[$ et sur $]1 ; 3]$, puis remonte sur $[3 ; +\infty[$ ; elle passe par un point haut en $(-1 ; 4)$ et par un point bas en $(3 ; 0)$ où elle touche l'axe des abscisses ; ailleurs, elle est strictement au-dessus de cet axe ; elle s'approche de la droite $y = 2$ quand $x$ devient très grand (en $+\infty$ comme en $-\infty$), et elle monte vers $+\infty$ des deux côtés de la droite $x = 1$.
\par \textbf{Questions.} 1) Dresser le tableau de variations conjecturé. 2) Formuler deux conjectures sur les asymptotes. 3) Que peut-on dire du signe de $f$ ?
\tcblower
\textbf{1) Tableau de variations conjecturé :} la valeur $1$ est interdite, on dédouble donc sa colonne et on la sépare par une double barre.
\begin{center}
\begin{tabular}{|c|ccccc||ccccc|}
\hline
\rowcolor{popblueL}
$x$ & $-\infty$ & & $-1$ & & $1$ & $1$ & & $3$ & & $+\infty$ \\
\hline
 & & & & & & & & & & \\
$f(x)$ & $2$ & $\nearrow$ & $4$ & $\searrow$ & $+\infty$ & $+\infty$ & $\searrow$ & $0$ & $\nearrow$ & $2$ \\
 & & & & & & & & & & \\
\hline
\end{tabular}
\end{center}
Autrement dit : $f$ est croissante sur $]-\infty ; -1]$, décroissante sur $[-1 ; 1[$, décroissante sur $]1 ; 3]$, croissante sur $[3 ; +\infty[$. Elle admet un maximum local égal à $4$, atteint en $x = -1$, et un minimum local égal à $0$, atteint en $x = 3$.
\par\textbf{2) Conjectures d'asymptotes :} la droite d'équation $x = 1$ semble être asymptote verticale (la courbe explose vers $+\infty$ de part et d'autre de cette droite) ; la droite d'équation $y = 2$ semble être asymptote horizontale en $+\infty$ et en $-\infty$ (la courbe s'en rapproche indéfiniment). Ces deux conjectures devront être prouvées par des calculs de limites : $\lim\limits_{x \to 1^-} f(x) = +\infty$, $\lim\limits_{x \to 1^+} f(x) = +\infty$ et $\lim\limits_{x \to +\infty} f(x) = 2$, $\lim\limits_{x \to -\infty} f(x) = 2$.
\par\textbf{3) Signe :} la courbe étant au-dessus de l'axe des abscisses, qu'elle touche seulement au point $(3 ; 0)$, on conjecture que $f(x) \geq 0$ pour tout $x$ de l'ensemble de définition, avec $f(x) = 0$ uniquement en $x = 3$. Le minimum $0$ atteint en $x = 3$ confirme cette lecture : $f(x) \geq 0$ partout sur $D_f$.
\end{exoresolu}

\begin{savaistu}[Conjecturer puis démontrer : des points faciles au Bac A]
Au Baccalauréat A, les problèmes d'analyse comportent presque toujours une question du type : \og À l'aide du graphique, conjecturer le sens de variation de $f$ / une asymptote / une symétrie, \emph{puis} démontrer votre conjecture \fg{}. La conjecture, correctement formulée en une phrase précise, rapporte souvent un point à elle seule ; la démonstration soignée en rapporte un ou deux autres. Les candidats qui rédigent proprement la démarche \og observer $\to$ conjecturer $\to$ vérifier \fg{} y gagnent des points faciles. À l'inverse, écrire \og d'après le graphique, c'est évident \fg{} sans aucun calcul fait perdre la totalité des points de la question. Retiens-le : au Bac, \textbf{le graphique suggère, le calcul décide}.
\end{savaistu}

\begin{exercice}[À toi de jouer]
Soit $h$ la fonction définie sur $\mathbb{R}$ par $h(x) = x^2 + 2x - 3$.
\begin{enumerate}
\item Compléter le tableau de valeurs pour $x \in \{-3 ; -2 ; -1 ; 0 ; 1\}$.
\item Quelle régularité observes-tu ? Formuler une conjecture sur un axe de symétrie de la courbe de $h$.
\item Démontrer cette conjecture en calculant $h(-1 + t)$ et $h(-1 - t)$ pour tout réel $t$ (on change de lettre pour éviter la confusion avec le nom de la fonction).
\end{enumerate}
\end{exercice}

\begin{corrige}[Exercice 1]
\textbf{1)} $h(-3) = 9 - 6 - 3 = 0$ ; $h(-2) = 4 - 4 - 3 = -3$ ; $h(-1) = 1 - 2 - 3 = -4$ ; $h(0) = -3$ ; $h(1) = 1 + 2 - 3 = 0$.
\par\textbf{2)} On observe $h(-2) = h(0) = -3$ et $h(-3) = h(1) = 0$ : les valeurs se répètent de part et d'autre de $x = -1$. \textbf{Conjecture :} la droite d'équation $x = -1$ semble être axe de symétrie de la courbe.
\par\textbf{3)} Pour tout réel $t$ :
\begin{align*}
h(-1 + t) &= (-1+t)^2 + 2(-1+t) - 3 = 1 - 2t + t^2 - 2 + 2t - 3 = t^2 - 4, \\
h(-1 - t) &= (-1-t)^2 + 2(-1-t) - 3 = 1 + 2t + t^2 - 2 - 2t - 3 = t^2 - 4.
\end{align*}
Donc $h(-1+t) = h(-1-t)$ pour tout réel $t$ : la conjecture est démontrée, la droite $x = -1$ est axe de symétrie de la parabole.
\end{corrige}

\begin{aretenir}[L'essentiel de cette section]
\begin{itemize}
\item La courbe représentative de $f$ est l'ensemble des points $(x ; f(x))$ pour $x \in D_f$ ; on travaille en repère orthonormal.
\item Lire une courbe se fait en trois temps : ensemble de définition, puis variations et extrema, puis comportements aux bornes.
\item Une régularité observée (droite d'appui, symétrie) donne une \cle{conjecture}, jamais une preuve.
\item La démarche complète est : \textbf{observer $\to$ conjecturer $\to$ vérifier par le calcul}.
\item Pour un axe de symétrie $x = a$, on teste $f(a+h) = f(a-h)$ ; pour un centre de symétrie $(a ; b)$, on teste $f(a+h) + f(a-h) = 2b$ ; pour une asymptote, on calcule une limite. Ces méthodes seront détaillées dans les sections suivantes.
\end{itemize}
\end{aretenir}

\coursec{Axe de symétrie d'une courbe}

Dans la section précédente, nous avons appris à \emph{lire} une courbe : à partir d'un tracé, tu as émis des conjectures sur les variations, les extremums, le signe d'une fonction. Or, en observant attentivement certaines courbes — la parabole de la fonction carré, par exemple — tu as sans doute remarqué une propriété frappante : la courbe semble se « plier » le long d'une droite verticale, chaque moitié étant le reflet exact de l'autre. Cette observation graphique n'est qu'une conjecture. Dans cette section, nous allons la transformer en certitude : nous allons définir rigoureusement ce qu'est un \cle{axe de symétrie} d'une courbe, puis apprendre à le \emph{démontrer} par le calcul, comme l'exige le Baccalauréat.

\courssub{Intuition et rappel géométrique : la symétrie orthogonale}

Imagine une feuille de papier sur laquelle tu as tracé la courbe d'une fonction. Plie la feuille le long d'une droite verticale $\Delta$. Si les deux morceaux de la courbe se superposent parfaitement, alors $\Delta$ est un axe de symétrie de la courbe. Ce pliage correspond, en géométrie, à la \cle{symétrie orthogonale} (ou réflexion) par rapport à la droite $\Delta$.

\begin{definition}[Symétrique d'un point par rapport à une droite verticale]
Soit $\Delta$ la droite d'équation $x = a$ dans un repère $(O\,;\vec{i},\vec{j})$. Le \cle{symétrique} d'un point $M(x\,;y)$ par rapport à $\Delta$ est le point $M'(x'\,;y')$ tel que $\Delta$ soit la médiatrice du segment $[MM']$. Comme $\Delta$ est verticale, le segment $[MM']$ est horizontal, et l'on obtient :
\[
x' = 2a - x \qquad \text{et} \qquad y' = y.
\]
\end{definition}

Vérifions cette formule sur un cas concret : le milieu de $[MM']$ a pour abscisse $\dfrac{x + (2a - x)}{2} = a$, donc il appartient bien à $\Delta$, et les deux points ont même ordonnée, donc $[MM']$ est bien perpendiculaire à $\Delta$. La formule est cohérente.

\begin{exemplebox}[Symétrique d'un point]
Par rapport à la droite d'équation $x = 3$, le symétrique du point $M(5\,; -2)$ est le point $M'(2 \times 3 - 5\,; -2) = M'(1\,; -2)$. Les points $M$ et $M'$ sont à égale distance de l'axe : $5 - 3 = 2$ et $3 - 1 = 2$.
\end{exemplebox}

\courssub{Définition de l'axe de symétrie d'une courbe}

Traduisons maintenant l'idée du pliage en langage fonctionnel. Dire que la courbe $\mathcal{C}_f$ est invariante par la symétrie d'axe $x = a$, c'est dire que pour tout point $M(x\,; f(x))$ de la courbe, son symétrique $M'(2a - x\,; f(x))$ appartient aussi à la courbe. Posons $x = a + h$ : le point $M$ est à « distance horizontale » $h$ à droite de l'axe, et son symétrique est le point d'abscisse $a - h$, à la même distance à gauche. L'appartenance de $M'$ à la courbe s'écrit alors $f(a - h) = f(a + h)$.

\begin{definition}[Axe de symétrie d'une courbe représentative]
Soit $f$ une fonction de domaine de définition $\mathcal{D}_f$ et $\mathcal{C}_f$ sa courbe représentative dans un repère. La droite $\Delta$ d'équation $x = a$ est \cle{axe de symétrie} de $\mathcal{C}_f$ si et seulement si :
\begin{enumerate}
\item le domaine $\mathcal{D}_f$ est \textbf{symétrique par rapport à $a$} : pour tout réel $h$, \; $a + h \in \mathcal{D}_f \iff a - h \in \mathcal{D}_f$ ;
\item pour tout réel $h$ tel que $a + h \in \mathcal{D}_f$ : 
\[ f(a - h) = f(a + h). \]
\end{enumerate}
\end{definition}

\begin{aretenir}[Lecture géométrique de la définition]
L'égalité $f(a - h) = f(a + h)$ signifie que deux points de la courbe situés à égale distance de part et d'autre de la droite $x = a$ ont \textbf{la même ordonnée}. La droite $x = a$ est donc la médiatrice du segment joignant ces deux points : c'est exactement la définition d'un axe de symétrie.
\end{aretenir}

\begin{propriete}[Caractérisation par le symétrique d'un point]
La droite d'équation $x = a$ est axe de symétrie de $\mathcal{C}_f$ si et seulement si, pour tout $x \in \mathcal{D}_f$, on a $2a - x \in \mathcal{D}_f$ et
\[ f(2a - x) = f(x). \]
\end{propriete}

\begin{experience}[Démonstration guidée : de la géométrie à la formule]
Nous allons démontrer que les deux formulations sont équivalentes, en partant de la définition du symétrique d'un point.
\begin{enumerate}
\item \textbf{Hypothèse.} Supposons que la droite $\Delta : x = a$ est axe de symétrie de $\mathcal{C}_f$ au sens géométrique : pour tout point $M$ de $\mathcal{C}_f$, son symétrique $M'$ par rapport à $\Delta$ appartient aussi à $\mathcal{C}_f$.
\item \textbf{Coordonnées.} Soit $M(x\,; f(x))$ un point de $\mathcal{C}_f$, avec $x \in \mathcal{D}_f$. D'après le rappel géométrique, son symétrique est $M'(2a - x\,; f(x))$.
\item \textbf{Traduction.} Dire que $M' \in \mathcal{C}_f$ signifie deux choses : d'une part $2a - x \in \mathcal{D}_f$ (le point existe sur la courbe), d'autre part l'ordonnée de $M'$ est l'image de son abscisse, c'est-à-dire $f(2a - x) = f(x)$.
\item \textbf{Lien avec la forme en $h$.} Posons $x = a + h$, c'est-à-dire $h = x - a$. Alors $2a - x = 2a - (a + h) = a - h$, et l'égalité $f(2a - x) = f(x)$ devient exactement $f(a - h) = f(a + h)$.
\item \textbf{Réciproque.} Inversement, si $f(a - h) = f(a + h)$ pour tout $h$ admissible, alors pour tout point $M(a + h\,; f(a+h))$ de la courbe, le point $M'(a - h\,; f(a+h))$ est aussi sur la courbe, et $\Delta$ est la médiatrice de $[MM']$ : la courbe est bien invariante par la symétrie d'axe $\Delta$. $\blacksquare$
\end{enumerate}
\end{experience}

\courssub{Cas particulier : l'axe des ordonnées et les fonctions paires}

Lorsque $a = 0$, la droite $x = 0$ est l'axe des ordonnées. La condition $f(-h) = f(h)$ pour tout $h \in \mathcal{D}_f$ (domaine symétrique par rapport à $0$) est exactement la définition d'une \cle{fonction paire}, que tu as rencontrée en classe de Seconde et de Première.

\begin{aretenir}[Fonction paire]
Une fonction $f$ est \textbf{paire} si et seulement si sa courbe représentative admet l'axe des ordonnées ($x = 0$) comme axe de symétrie. Exemples : $x \mapsto x^2$, $x \mapsto |x|$, $x \mapsto \cos x$, $x \mapsto \dfrac{1}{x^2}$.
\end{aretenir}

La notion d'axe de symétrie $x = a$ généralise donc la parité : une fonction dont la courbe admet l'axe $x = a$ est une fonction « paire après translation ». Cette remarque fournit d'ailleurs une stratégie de démonstration élégante, que nous détaillons dans la méthode suivante.

\courssub{Méthode de démonstration}

\begin{methode}[Prouver que la droite $x = a$ est axe de symétrie de $\mathcal{C}_f$]
\begin{enumerate}
\item \textbf{Vérifier le domaine.} Montrer que $\mathcal{D}_f$ est symétrique par rapport à $a$ : pour tout réel $h$, $a + h \in \mathcal{D}_f \iff a - h \in \mathcal{D}_f$. Cette étape est souvent négligée, elle est pourtant indispensable.
\item \textbf{Calculer $f(a + h)$.} Remplacer $x$ par $a + h$ dans l'expression de $f$ et simplifier.
\item \textbf{Calculer $f(a - h)$.} Remplacer $x$ par $a - h$ et simplifier.
\item \textbf{Comparer et conclure.} Si les deux expressions sont égales pour tout $h$ admissible, conclure par une phrase : « Pour tout $h$ tel que $a + h \in \mathcal{D}_f$, $f(a - h) = f(a + h)$, donc la droite d'équation $x = a$ est axe de symétrie de $\mathcal{C}_f$. »
\end{enumerate}
\textbf{Variante (changement de repère).} On peut aussi poser $X = x - a$ et étudier la fonction $g(X) = f(X + a)$ : la droite $x = a$ est axe de symétrie de $\mathcal{C}_f$ si et seulement si $g$ est \textbf{paire}. Cette approche est très efficace quand les calculs directs sont lourds.
\end{methode}

\begin{attention}[Le piège du domaine oublié]
L'erreur la plus fréquente au Baccalauréat consiste à vérifier l'égalité $f(a - h) = f(a + h)$ \emph{sans s'être assuré au préalable que le domaine est symétrique par rapport à $a$}. Contre-exemple : soit $f(x) = \sqrt{x}$ définie sur $\mathcal{D}_f = [0\,; +\infty[$. Pour $h = 1$, on aurait envie de comparer $f(1)$ et $f(-1)$, mais $f(-1)$ n'existe pas ! Un axe de symétrie $x = a$ impose que la courbe existe des deux côtés de l'axe, à distances égales. Rédige toujours la condition sur le domaine \textbf{avant} le calcul des images.
\end{attention}

\courssub{Exemples détaillés}

\begin{exoresolu}[Une parabole et son axe]
Soit $f$ la fonction définie sur $\mathbb{R}$ par $f(x) = x^2 - 2x + 3$. Démontrer que la droite $\Delta$ d'équation $x = 1$ est axe de symétrie de la courbe $\mathcal{C}_f$.
\tcblower
\textbf{Solution détaillée.}

\textbf{Étape 1 : le domaine.} $\mathcal{D}_f = \mathbb{R}$ est symétrique par rapport à tout réel, en particulier par rapport à $1$ : pour tout $h \in \mathbb{R}$, $1 + h \in \mathbb{R}$ et $1 - h \in \mathbb{R}$.

\textbf{Étape 2 : calcul de $f(1 + h)$.}
\begin{align*}
f(1 + h) &= (1 + h)^2 - 2(1 + h) + 3 \\
&= 1 + 2h + h^2 - 2 - 2h + 3 \\
&= h^2 + 2.
\end{align*}

\textbf{Étape 3 : calcul de $f(1 - h)$.}
\begin{align*}
f(1 - h) &= (1 - h)^2 - 2(1 - h) + 3 \\
&= 1 - 2h + h^2 - 2 + 2h + 3 \\
&= h^2 + 2.
\end{align*}

\textbf{Étape 4 : conclusion.} Pour tout réel $h$, $f(1 - h) = f(1 + h) = h^2 + 2$. Donc la droite d'équation $x = 1$ est axe de symétrie de $\mathcal{C}_f$.

\textbf{Remarque.} On retrouve ici la forme canonique : $f(x) = (x - 1)^2 + 2$. Le changement de variable $X = x - 1$ donne $g(X) = X^2 + 2$, qui est bien une fonction paire. Le sommet de la parabole, $S(1\,; 2)$, appartient à l'axe de symétrie : c'est toujours le cas pour une parabole.
\end{exoresolu}

\begin{popfigure}
\begin{center}
\begin{tikzpicture}
\begin{axis}[
  width=0.85\linewidth, height=7cm,
  axis lines=middle,
  xlabel={$x$}, ylabel={$y$},
  xmin=-2.2, xmax=4.2, ymin=-0.5, ymax=7.5,
  xtick={-2,-1,0,1,2,3,4},
  ytick={1,2,3,4,5,6,7},
  grid=both, grid style={popline!40},
  tick label style={font=\small},
]
\addplot[popcurve,domain=-1.7:3.7,samples=200]{x^2 - 2*x + 3};
\addplot[poporange,dashed,thick] coordinates {(1,-0.5) (1,7.5)};
\addplot[only marks,mark=*,popblue,mark size=2pt] coordinates {(-1,6) (3,6)};
\draw[popgreen,thick] (axis cs:-1,6) -- (axis cs:3,6);
\node[popblue,above left,font=\small] at (axis cs:-1,6) {$M(-1\,;6)$};
\node[popblue,above right,font=\small] at (axis cs:3,6) {$M'(3\,;6)$};
\node[poporange,font=\small,rotate=90,above] at (axis cs:1,6.6) {$x = 1$};
\addplot[only marks,mark=*,popink,mark size=2pt] coordinates {(1,2)};
\node[popink,below right,font=\small] at (axis cs:1,2) {$S(1\,;2)$};
\end{axis}
\end{tikzpicture}
\end{center}
\legende{La parabole d'équation $y = x^2 - 2x + 3$ et son axe de symétrie $x = 1$. Les points $M(-1\,;6)$ et $M'(3\,;6)$, situés à égale distance de l'axe, ont la même ordonnée : $f(1 - 2) = f(1 + 2) = 6$.}
\end{popfigure}

\begin{exoresolu}[Une fonction valeur absolue]
Soit $f$ la fonction définie sur $\mathbb{R}$ par $f(x) = |x - 3| + 1$. Montrer que la droite d'équation $x = 3$ est axe de symétrie de $\mathcal{C}_f$.
\tcblower
\textbf{Solution détaillée.}

\textbf{Étape 1.} $\mathcal{D}_f = \mathbb{R}$ est symétrique par rapport à $3$.

\textbf{Étape 2.} Pour tout réel $h$ :
\[ f(3 + h) = |(3 + h) - 3| + 1 = |h| + 1. \]

\textbf{Étape 3.} De même :
\[ f(3 - h) = |(3 - h) - 3| + 1 = |-h| + 1 = |h| + 1, \]
car $|-h| = |h|$ pour tout réel $h$.

\textbf{Étape 4.} Pour tout réel $h$, $f(3 - h) = f(3 + h)$, donc la droite d'équation $x = 3$ est axe de symétrie de $\mathcal{C}_f$. Géométriquement, la courbe de $f$ est un « V » dont la pointe, le point $(3\,; 1)$, repose sur l'axe.
\end{exoresolu}

\courssub{Le cas des fonctions rationnelles}

Les fonctions rationnelles fournissent des exemples plus riches, car le domaine de définition n'est plus $\mathbb{R}$ tout entier : la vérification de la symétrie du domaine devient une étape réellement significative.

\begin{exoresolu}[Une fonction rationnelle avec valeur interdite]
Soit $f$ la fonction définie par $f(x) = \dfrac{1}{(x - 2)^2}$.
\begin{enumerate}
\item Déterminer le domaine de définition $\mathcal{D}_f$ et vérifier qu'il est symétrique par rapport à $2$.
\item Démontrer que la droite d'équation $x = 2$ est axe de symétrie de $\mathcal{C}_f$.
\end{enumerate}
\tcblower
\textbf{Solution détaillée.}

\textbf{1. Domaine.} $f(x)$ existe si et seulement si $(x - 2)^2 \neq 0$, c'est-à-dire $x \neq 2$. Donc
\[ \mathcal{D}_f = \mathbb{R} \setminus \{2\} = \left]-\infty\,; 2\right[ \cup \left]2\,; +\infty\right[. \]
Soit $h$ un réel. On a la chaîne d'équivalences :
\[ 2 + h \in \mathcal{D}_f \iff 2 + h \neq 2 \iff h \neq 0 \iff 2 - h \neq 2 \iff 2 - h \in \mathcal{D}_f. \]
Le domaine est donc bien symétrique par rapport à $2$ : la valeur interdite $x = 2$ est précisément l'abscisse de l'axe de symétrie.

\textbf{2. Égalité des images.} Pour tout réel $h \neq 0$ :
\[ f(2 + h) = \frac{1}{\big((2 + h) - 2\big)^2} = \frac{1}{h^2} \qquad \text{et} \qquad f(2 - h) = \frac{1}{\big((2 - h) - 2\big)^2} = \frac{1}{(-h)^2} = \frac{1}{h^2}. \]
Donc $f(2 - h) = f(2 + h)$ pour tout $h \neq 0$ : la droite d'équation $x = 2$ est axe de symétrie de $\mathcal{C}_f$.

\textbf{Remarque.} La fonction $g(X) = f(X + 2) = \dfrac{1}{X^2}$ est paire : le changement de repère $X = x - 2$ confirme immédiatement le résultat.
\end{exoresolu}

\begin{popfigure}
\begin{center}
\begin{tikzpicture}
\begin{axis}[
  width=0.85\linewidth, height=7.5cm,
  axis lines=middle,
  xlabel={$x$}, ylabel={$y$},
  xmin=-2.5, xmax=6.5, ymin=-0.5, ymax=8,
  xtick={-2,-1,0,1,2,3,4,5,6},
  ytick={1,2,3,4,5,6,7},
  grid=both, grid style={popline!40},
  tick label style={font=\small},
]
\addplot[popcurve,domain=-2.3:1.6,samples=150]{1/(x-2)^2};
\addplot[popcurve,domain=2.4:6.3,samples=150]{1/(x-2)^2};
\addplot[poporange,dashed,thick] coordinates {(2,-0.5) (2,8)};
\node[poporange,font=\small,rotate=90,above] at (axis cs:2,7.2) {$x = 2$};
\addplot[only marks,mark=*,popblue,mark size=2pt] coordinates {(0,0.25) (4,0.25)};
\draw[popgreen,thick] (axis cs:0,0.25) -- (axis cs:4,0.25);
\node[popblue,below,font=\small] at (axis cs:0,0.25) {$M$};
\node[popblue,below,font=\small] at (axis cs:4,0.25) {$M'$};
\end{axis}
\end{tikzpicture}
\end{center}
\legende{Courbe de $f(x) = \dfrac{1}{(x-2)^2}$. La droite $x = 2$ est axe de symétrie : les points $M(0\,; 0{,}25)$ et $M'(4\,; 0{,}25)$ sont symétriques, car $f(2 - 2) = f(2 + 2) = \dfrac{1}{4}$. On observe aussi que la courbe « s'envole » près de l'axe : nous y reviendrons avec la notion d'asymptote verticale.}
\end{popfigure}

\begin{attention}[Deux pièges de rédaction à éviter]
\begin{itemize}
\item \textbf{Confondre les rôles.} L'égalité à établir est $f(a - h) = f(a + h)$, et non $f(a - x) = f(a + x)$ avec $x$ « quelconque » sans préciser le domaine de variation de la variable. Annonce toujours clairement : « Soit $h$ un réel tel que $a + h \in \mathcal{D}_f$ ».
\item \textbf{Conclure trop vite à partir de la calculatrice.} Observer sur l'écran que la courbe semble symétrique est une conjecture (section 1), pas une preuve. Seul le calcul algébrique de $f(a + h)$ et $f(a - h)$, valable pour \emph{tout} $h$ admissible, constitue une démonstration.
\end{itemize}
\end{attention}

\begin{savaistu}[La symétrie, un outil d'économie de travail]
Savoir qu'une courbe admet un axe de symétrie $x = a$ permet de diviser le travail d'étude par deux : il suffit d'étudier la fonction sur la moitié du domaine (par exemple pour $x \geq a$), puis de compléter le tracé par symétrie. Cette stratégie sera précieuse dans la section 6, lorsque nous construirons des courbes complètes. Les ingénieurs utilisent le même principe : dans la conception d'une antenne parabolique — comme celles que l'on voit sur les toits de Yaoundé ou de Douala pour capter les signaux satellite — le profil est symétrique par rapport à l'axe de la parabole, ce qui garantit que tous les rayons réfléchis convergent vers le foyer, où est placé le récepteur.
\end{savaistu}

\begin{exercice}[À toi de jouer]
\begin{enumerate}
\item Soit $f(x) = -2x^2 + 8x - 5$ définie sur $\mathbb{R}$. Démontrer que la droite d'équation $x = 2$ est axe de symétrie de $\mathcal{C}_f$.
\item Soit $g(x) = \dfrac{3}{|x + 1|}$. Déterminer $\mathcal{D}_g$, puis montrer que la droite d'équation $x = -1$ est axe de symétrie de $\mathcal{C}_g$.
\item La courbe de la fonction $h(x) = \dfrac{1}{x^2 - 4x + 5}$ admet-elle un axe de symétrie ? (Indication : écrire le dénominateur sous forme canonique.)
\end{enumerate}
\end{exercice}

\begin{corrige}[Exercice : À toi de jouer]
\textbf{1.} $\mathcal{D}_f = \mathbb{R}$, symétrique par rapport à $2$. Pour tout réel $h$ :
\begin{align*}
f(2 + h) &= -2(2 + h)^2 + 8(2 + h) - 5 = -2(4 + 4h + h^2) + 16 + 8h - 5 = -2h^2 + 3,\\
f(2 - h) &= -2(2 - h)^2 + 8(2 - h) - 5 = -2(4 - 4h + h^2) + 16 - 8h - 5 = -2h^2 + 3.
\end{align*}
Donc $f(2 - h) = f(2 + h)$ pour tout $h$ : la droite $x = 2$ est axe de symétrie de $\mathcal{C}_f$.

\textbf{2.} $g(x)$ existe ssi $x \neq -1$, donc $\mathcal{D}_g = \mathbb{R} \setminus \{-1\}$, symétrique par rapport à $-1$ (car $-1 + h \neq -1 \iff h \neq 0 \iff -1 - h \neq -1$). Pour $h \neq 0$ :
\[ g(-1 + h) = \frac{3}{|h|} \qquad \text{et} \qquad g(-1 - h) = \frac{3}{|-h|} = \frac{3}{|h|}. \]
Donc $g(-1 - h) = g(-1 + h)$ : la droite $x = -1$ est axe de symétrie de $\mathcal{C}_g$.

\textbf{3.} Forme canonique du dénominateur : $x^2 - 4x + 5 = (x - 2)^2 + 1$. Posons $X = x - 2$ : alors $h(x) = \dfrac{1}{X^2 + 1}$, fonction paire en $X$. Vérification directe : pour tout réel $h$ (le dénominateur ne s'annule jamais, $\mathcal{D}_h = \mathbb{R}$) :
\[ h(2 + h) = \frac{1}{h^2 + 1} = h(2 - h). \]
Donc la droite d'équation $x = 2$ est axe de symétrie de la courbe de $h$.
\end{corrige}

\begin{aretenir}[L'essentiel de la section]
\begin{itemize}
\item La droite $x = a$ est \cle{axe de symétrie} de $\mathcal{C}_f$ si le domaine $\mathcal{D}_f$ est symétrique par rapport à $a$ \textbf{et} si $f(a - h) = f(a + h)$ pour tout $h$ admissible.
\item Le cas $a = 0$ redonne les fonctions \cle{paires} (axe des ordonnées).
\item La démonstration type comporte trois temps : symétrie du domaine, calculs de $f(a + h)$ et $f(a - h)$, conclusion rédigée.
\item Le changement de variable $X = x - a$ ramène le problème à la parité de $g(X) = f(X + a)$ : une stratégie souvent plus rapide.
\end{itemize}
Dans la section suivante, nous étudierons l'autre grande symétrie des courbes : le \emph{centre} de symétrie, qui généralise quant à lui les fonctions impaires.
\end{aretenir}

\coursec{Centre de symétrie d'une courbe}

Dans la section précédente, nous avons vu comment une courbe peut être \emph{pliée} le long d'une droite verticale $x = a$ : chaque point $M$ de la courbe a alors un jumeau $M'$, son symétrique par rapport à cet axe. Il existe une autre manière remarquable pour une courbe de se reproduire elle-même : non plus par \emph{pliage}, mais par \emph{demi-tour} autour d'un point. C'est l'objet de cette section. Tu verras que la démarche est très proche de celle de l'axe de symétrie : une intuition géométrique, une traduction analytique, une méthode de démonstration, et enfin des applications concrètes.

\courssub{Intuition et rappels sur la symétrie centrale}

\textbf{Une expérience simple.} Prends une feuille, trace une courbe, pique un point $\Omega$ avec ton compas et fais tourner la feuille d'un demi-tour ($180\degres$) autour de ce point. Si la courbe retombe exactement sur elle-même, alors $\Omega$ est un \cle{centre de symétrie} de la courbe. C'est ce que tu observes sur la lettre S, sur la courbe de la fonction cube, ou encore sur certaines décorations des pagnes.

\textbf{Rappel de géométrie (Seconde).} La \cle{symétrie centrale} de centre $\Omega$ est la transformation qui, à tout point $M$, associe le point $M'$ tel que $\Omega$ soit le \emph{milieu} du segment $[MM']$. Autrement dit :
\[
\overrightarrow{\Omega M'} = -\overrightarrow{\Omega M} \qquad \text{ou encore} \qquad \overrightarrow{OM'} = 2\overrightarrow{O\Omega} - \overrightarrow{OM}.
\]

\begin{propriete}[Coordonnées du symétrique d'un point]
Soient $\Omega(a\,;\,b)$ et $M(x\,;\,y)$ deux points du plan muni d'un repère $(O\,;\vec{\imath},\vec{\jmath})$. Le symétrique de $M$ par rapport à $\Omega$ est le point $M'$ de coordonnées :
\[
M'\bigl(2a - x\,;\, 2b - y\bigr).
\]
\end{propriete}

\begin{experience}[Démonstration : pourquoi ces coordonnées ?]
Dire que $\Omega$ est le milieu de $[MM']$, c'est dire, en notant $M'(x'\,;\,y')$ :
\[
\frac{x + x'}{2} = a \qquad \text{et} \qquad \frac{y + y'}{2} = b.
\]
De la première égalité on tire $x + x' = 2a$, donc $x' = 2a - x$. De la seconde, $y' = 2b - y$. C'est exactement la formule annoncée. Retiens l'idée clé : \textbf{le centre de symétrie, c'est un milieu}.
\end{experience}

\courssub{Définition du centre de symétrie d'une courbe}

Plaçons-nous maintenant sur la courbe représentative $\mathcal{C}_f$ d'une fonction $f$. Un point $M$ de la courbe a pour coordonnées $(x\,;\,f(x))$. Son symétrique par rapport à $\Omega(a\,;\,b)$ est $M'\bigl(2a - x\,;\,2b - f(x)\bigr)$. Dire que la courbe est \emph{globalement invariante} par la symétrie de centre $\Omega$, c'est dire que $M'$ est \emph{aussi} un point de la courbe, c'est-à-dire que son ordonnée vaut $f(2a - x)$ :
\[
f(2a - x) = 2b - f(x).
\]

\begin{definition}[Centre de symétrie d'une courbe représentative]
Soit $f$ une fonction de domaine de définition $D_f$, de courbe représentative $\mathcal{C}_f$, et soit $\Omega(a\,;\,b)$ un point du plan. On dit que $\Omega$ est un \cle{centre de symétrie} de $\mathcal{C}_f$ lorsque les deux conditions suivantes sont réalisées :
\begin{enumerate}
\item le domaine $D_f$ est \textbf{symétrique par rapport à $a$} : pour tout réel $h$, si $a + h \in D_f$ alors $a - h \in D_f$ ;
\item pour tout $h$ tel que $a + h \in D_f$ :
\[
f(a - h) + f(a + h) = 2b.
\]
\end{enumerate}
\end{definition}

\begin{propriete}[Écritures équivalentes]
Les assertions suivantes sont équivalentes :
\begin{enumerate}
\item $\Omega(a\,;\,b)$ est centre de symétrie de $\mathcal{C}_f$ ;
\item pour tout $h$ tel que $a + h \in D_f$ : \quad $f(a + h) + f(a - h) = 2b$ ;
\item pour tout $h$ tel que $a + h \in D_f$ : \quad $f(a + h) - b = -\bigl(f(a - h) - b\bigr)$ ;
\item pour tout $x \in D_f$ : \quad $f(2a - x) = 2b - f(x)$.
\end{enumerate}
\end{propriete}

\begin{experience}[{Démonstration guidée : traduire « $\Omega$ milieu de $[MM']$ »}]
Montrons que l'assertion 2 traduit exactement le fait que $\Omega(a\,;\,b)$ est le milieu de $[MM']$, où $M$ et $M'$ sont deux points de la courbe échangés par la symétrie.

\textbf{Étape 1.} Prenons un point de la courbe dont l'abscisse s'écrit $a + h$ (c'est-à-dire situé à la distance horizontale $h$ à droite de $a$) : $M\bigl(a + h\,;\,f(a+h)\bigr)$.

\textbf{Étape 2.} Le point symétrique par rapport à $\Omega$ a pour abscisse $a - h$. Pour qu'il soit sur la courbe, son ordonnée doit être $f(a - h)$ : $M'\bigl(a - h\,;\,f(a-h)\bigr)$.

\textbf{Étape 3.} Le milieu de $[MM']$ a pour coordonnées :
\[
\left(\frac{(a+h)+(a-h)}{2}\,;\,\frac{f(a+h)+f(a-h)}{2}\right) = \left(a\,;\,\frac{f(a+h)+f(a-h)}{2}\right).
\]

\textbf{Étape 4.} Ce milieu est le point $\Omega(a\,;\,b)$ si et seulement si :
\[
\frac{f(a+h)+f(a-h)}{2} = b \qquad \Longleftrightarrow \qquad f(a+h)+f(a-h) = 2b.
\]

\textbf{Étape 5.} Enfin, $f(a+h)+f(a-h) = 2b$ équivaut, en soustrayant $2b$ et en répartissant, à $f(a+h) - b = -\bigl(f(a-h) - b\bigr)$ : les écarts verticaux par rapport à $b$ sont opposés de part et d'autre de $\Omega$. C'est exactement l'image du demi-tour. \hfill $\square$
\end{experience}

\begin{attention}[Ne confonds pas centre et axe de symétrie !]
L'erreur la plus fréquente au Baccalauréat est de mélanger les deux conditions :
\begin{itemize}
\item $f(a + h) = f(a - h)$ caractérise un \textbf{axe de symétrie} d'équation $x = a$ (les ordonnées sont \emph{égales}) ;
\item $f(a + h) + f(a - h) = 2b$ caractérise un \textbf{centre de symétrie} $\Omega(a\,;\,b)$ (la \emph{demi-somme} des ordonnées vaut $b$).
\end{itemize}
Avant de calculer, demande-toi toujours : « est-ce un pliage (axe) ou un demi-tour (centre) ? ». Autre piège : oublier de vérifier que le domaine $D_f$ est symétrique par rapport à $a$. Cette condition est indispensable et doit figurer dans ta rédaction.
\end{attention}

\courssub{Cas particulier : l'origine et les fonctions impaires}

Lorsque le centre de symétrie est l'origine $O(0\,;\,0)$ du repère, c'est-à-dire $a = 0$ et $b = 0$, la condition $f(a+h) + f(a-h) = 2b$ devient :
\[
f(h) + f(-h) = 0 \qquad \Longleftrightarrow \qquad f(-h) = -f(h).
\]
On retrouve exactement la définition d'une \cle{fonction impaire}, vue en classe de Première.

\begin{aretenir}[Fonctions impaires]
Une fonction $f$, définie sur un domaine $D_f$ symétrique par rapport à $0$, est \textbf{impaire} si et seulement si sa courbe représentative admet l'origine $O$ du repère pour centre de symétrie. Exemples : $x \mapsto x^3$, $x \mapsto \sin x$, $x \mapsto \dfrac{1}{x}$.
\end{aretenir}

\begin{exemplebox}[Vérifier que $O$ est centre de symétrie de la courbe de $f(x) = x^3$]
Soit $f(x) = x^3$, définie sur $D_f = \mathbb{R}$, qui est bien symétrique par rapport à $0$. Pour tout réel $h$ :
\[
f(h) + f(-h) = h^3 + (-h)^3 = h^3 - h^3 = 0 = 2 \times 0.
\]
La condition $f(0 + h) + f(0 - h) = 2 \times 0$ est vérifiée : l'origine $O$ est bien un centre de symétrie de la courbe de $f$. On dit que $f$ est impaire.
\end{exemplebox}

\courssub{Méthode de démonstration}

\begin{methode}[Prouver que $\Omega(a\,;\,b)$ est centre de symétrie de $\mathcal{C}_f$]
\begin{enumerate}
\item \textbf{Vérifier le domaine} : montrer que $D_f$ est symétrique par rapport à $a$ (si $a + h \in D_f$, alors $a - h \in D_f$).
\item \textbf{Calculer séparément} $f(a + h)$ et $f(a - h)$ en fonction de $h$.
\item \textbf{Former la demi-somme} : $\dfrac{f(a + h) + f(a - h)}{2}$.
\item \textbf{Conclure} : si cette demi-somme est égale à $b$ pour tout $h$ admissible, alors $\Omega(a\,;\,b)$ est centre de symétrie de $\mathcal{C}_f$. Sinon, ce n'en est pas un.
\end{enumerate}
\textbf{Astuce de rédaction} : la demi-somme est souvent plus lisible que la somme, car on « voit » apparaître directement le nombre $b$.
\end{methode}

\begin{exoresolu}[Montrer que $\Omega(1\,;\,2)$ est centre de symétrie de la courbe de $f(x) = \dfrac{2x-1}{x-1}$]
Énoncé. Soit $f$ la fonction définie par $f(x) = \dfrac{2x - 1}{x - 1}$ et $\mathcal{C}_f$ sa courbe représentative. Démontrer que le point $\Omega(1\,;\,2)$ est un centre de symétrie de $\mathcal{C}_f$.
\tcblower
Solution détaillée.

\textbf{Étape 1 : le domaine.} $f$ est définie pour $x - 1 \neq 0$, donc $D_f = \mathbb{R} \setminus \{1\}$. Si $1 + h \in D_f$, alors $h \neq 0$, donc $-h \neq 0$ et $1 - h \in D_f$ : le domaine est bien symétrique par rapport à $1$.

\textbf{Étape 2 : calcul de $f(1 + h)$ et $f(1 - h)$ pour $h \neq 0$.}
\begin{align*}
f(1 + h) &= \frac{2(1 + h) - 1}{(1 + h) - 1} = \frac{1 + 2h}{h} = \frac{1}{h} + 2,\\[4pt]
f(1 - h) &= \frac{2(1 - h) - 1}{(1 - h) - 1} = \frac{1 - 2h}{-h} = -\frac{1}{h} + 2.
\end{align*}

\textbf{Étape 3 : la demi-somme.}
\[
\frac{f(1 + h) + f(1 - h)}{2} = \frac{\left(\dfrac{1}{h} + 2\right) + \left(-\dfrac{1}{h} + 2\right)}{2} = \frac{4}{2} = 2.
\]

\textbf{Étape 4 : conclusion.} Pour tout $h \neq 0$, $\dfrac{f(1+h) + f(1-h)}{2} = 2 = b$. Donc le point $\Omega(1\,;\,2)$ est bien un centre de symétrie de $\mathcal{C}_f$.

\textbf{Remarque.} On pouvait le deviner en réécrivant $f(x) = \dfrac{2(x-1) + 1}{x - 1} = 2 + \dfrac{1}{x - 1}$ : la courbe de $f$ est l'image de l'hyperbole $y = \dfrac{1}{x}$ (de centre $O$) par la translation de vecteur $\vec{\imath} + 2\vec{\jmath}$, qui envoie $O$ sur $\Omega(1\,;\,2)$.
\end{exoresolu}

\begin{popfigure}
\begin{center}
\begin{tikzpicture}
\begin{axis}[
  width=0.92\linewidth, height=8cm,
  axis lines=middle, axis line style={->, popdark},
  xlabel={$x$}, ylabel={$y$},
  xlabel style={at={(ticklabel* cs:1)}, anchor=west},
  ylabel style={at={(ticklabel* cs:1)}, anchor=south},
  xmin=-3.5, xmax=5.5, ymin=-3.5, ymax=7.5,
  xtick={-3,-2,-1,1,2,3,4,5}, ytick={-3,-2,-1,1,2,3,4,5,6,7},
  tick label style={font=\small, color=popmuted},
  grid=both, grid style={popline!40, very thin},
  clip=true
]
  \addplot[popcurve, domain=-3.5:0.78, samples=120] {(2*x-1)/(x-1)};
  \addplot[popcurve, domain=1.22:5.5, samples=120] {(2*x-1)/(x-1)};
  \addplot[popline, dashed, domain=-3.5:5.5] {2};
  \draw[popline, dashed] (axis cs:1,-3.5) -- (axis cs:1,7.5);
  \addplot[only marks, mark=*, mark size=2.2pt, poporange] coordinates {(1,2)};
  \node[poporange, anchor=south west, font=\small] at (axis cs:1.1,2.1) {$\Omega(1\,;\,2)$};
  \addplot[only marks, mark=*, mark size=1.8pt, popgreen] coordinates {(3,2.5) (-1,1.5)};
  \node[popgreen, anchor=south west, font=\small] at (axis cs:3.05,2.5) {$M$};
  \node[popgreen, anchor=north east, font=\small] at (axis cs:-1.05,1.45) {$M'$};
  \draw[popgreen, dotted, thick] (axis cs:-1,1.5) -- (axis cs:3,2.5);
\end{axis}
\end{tikzpicture}
\end{center}
\legende{Courbe de $f(x) = \dfrac{2x-1}{x-1}$ : le point $\Omega(1\,;\,2)$ est le milieu de tout segment $[MM']$ joignant deux points symétriques de la courbe, par exemple $M(3\,;\,2{,}5)$ et $M'(-1\,;\,1{,}5)$.}
\end{popfigure}

Vérifie sur la figure : $f(3) = \dfrac{5}{2} = 2{,}5$ et $f(-1) = \dfrac{3}{2} = 1{,}5$. Le milieu de $[MM']$ est $\left(\dfrac{3 + (-1)}{2}\,;\,\dfrac{2{,}5 + 1{,}5}{2}\right) = (1\,;\,2) = \Omega$. La géométrie et le calcul se rejoignent parfaitement.

\courssub{Lien avec les fonctions cubiques}

Les fonctions polynômes de degré 3, dites \cle{cubiques}, offrent une famille remarquable d'exemples. Considérons $f(x) = x^3 - 3x$. Son domaine $\mathbb{R}$ est symétrique par rapport à $0$, et :
\[
f(-x) = (-x)^3 - 3(-x) = -x^3 + 3x = -\bigl(x^3 - 3x\bigr) = -f(x).
\]
La fonction est impaire : sa courbe admet l'origine $O$ pour centre de symétrie.

\begin{popfigure}
\begin{center}
\begin{tikzpicture}
\begin{axis}[
  width=0.92\linewidth, height=8cm,
  axis lines=middle, axis line style={->, popdark},
  xlabel={$x$}, ylabel={$y$},
  xlabel style={at={(ticklabel* cs:1)}, anchor=west},
  ylabel style={at={(ticklabel* cs:1)}, anchor=south},
  xmin=-3, xmax=3, ymin=-4.5, ymax=4.5,
  xtick={-3,-2,-1,1,2,3}, ytick={-4,-2,2,4},
  tick label style={font=\small, color=popmuted},
  grid=both, grid style={popline!40, very thin},
  clip=true
]
  \addplot[popcurve, domain=-2.35:2.35, samples=200] {x^3 - 3*x};
  \addplot[only marks, mark=*, mark size=2.2pt, poporange] coordinates {(0,0)};
  \node[poporange, anchor=south west, font=\small] at (axis cs:0.08,0.15) {$O$};
  \addplot[only marks, mark=*, mark size=1.8pt, popgreen] coordinates {(1.5,-1.125) (-1.5,1.125)};
  \node[popgreen, anchor=north west, font=\small] at (axis cs:1.55,-1.2) {$M$};
  \node[popgreen, anchor=south east, font=\small] at (axis cs:-1.55,1.2) {$M'$};
  \draw[popgreen, dotted, thick] (axis cs:-1.5,1.125) -- (axis cs:1.5,-1.125);
\end{axis}
\end{tikzpicture}
\end{center}
\legende{Courbe de $f(x) = x^3 - 3x$ : l'origine $O$ est centre de symétrie. Les points $M(1{,}5\,;\,-1{,}125)$ et $M'(-1{,}5\,;\,1{,}125)$ sont échangés par la symétrie de centre $O$.}
\end{popfigure}

\begin{propriete}[Centre de symétrie d'une courbe cubique (admis)]
Soit $f$ une fonction polynôme de degré 3, $f(x) = \alpha x^3 + \beta x^2 + \gamma x + \delta$ avec $\alpha \neq 0$. Sa courbe représentative admet toujours un centre de symétrie : c'est le point $\Omega$ d'abscisse $x_0 = -\dfrac{\beta}{3\alpha}$, appelé \cle{point d'inflexion} de la courbe (là où la courbe « traverse » sa tangente en changeant de sens de courbure). Ce résultat est admis à ce niveau ; il sera justifié en classe de Terminale C/D avec la dérivée seconde.
\end{propriete}

\begin{exemplebox}[Le point d'inflexion d'une cubique générale]
Pour $f(x) = x^3 - 3x^2 + 1$, on a $\alpha = 1$, $\beta = -3$, donc $x_0 = -\dfrac{-3}{3 \times 1} = 1$ et $f(1) = 1 - 3 + 1 = -1$. Le point $\Omega(1\,;\,-1)$ est centre de symétrie de la courbe. Vérifions-le avec notre méthode :
\begin{align*}
f(1 + h) &= (1+h)^3 - 3(1+h)^2 + 1 = h^3 - 3h - 1,\\
f(1 - h) &= (1-h)^3 - 3(1-h)^2 + 1 = -h^3 + 3h - 1.
\end{align*}
La demi-somme vaut $\dfrac{(h^3 - 3h - 1) + (-h^3 + 3h - 1)}{2} = \dfrac{-2}{2} = -1 = b$. La propriété est confirmée.
\end{exemplebox}

\begin{attention}[Un centre de symétrie n'est pas forcément sur les axes]
Le point $\Omega(1\,;\,-1)$ de l'exemple précédent n'est ni l'origine, ni sur un axe du repère. Ne cherche donc pas systématiquement $O$ comme centre de symétrie : lis l'énoncé, ou conjecture la position de $\Omega$ sur le graphique (c'est le point où la courbe semble « pivoter » sur elle-même), puis démontre.
\end{attention}

\begin{savaistu}[Pourquoi les mathématiciens aiment tant les symétries ?]
Démontrer qu'une courbe admet un centre de symétrie $\Omega(a\,;\,b)$ permet de \textbf{réduire de moitié le travail d'étude} : il suffit d'étudier la fonction sur la moitié du domaine, par exemple pour $x \geq a$, de dresser le tableau de variations et de tracer la demi-courbe correspondante, puis de \emph{compléter} par la symétrie de centre $\Omega$. C'est exactement le même gain que pour un axe de symétrie ou pour la parité. Ce principe — « la symétrie économise le travail » — est l'une des idées les plus profondes des mathématiques : le grand mathématicien allemand Felix Klein affirmait au XIX\textsuperscript{e} siècle que toute la géométrie pouvait se décrire en termes de symétries. En physique, un théorème célèbre d'Emmy Noether relie chaque symétrie de la nature à une loi de conservation (énergie, quantité de mouvement...). La prochaine fois que tu traceras une hyperbole ou une cubique, souviens-toi : tu manipules une idée qui fait tourner la science moderne !
\end{savaistu}

\begin{exercice}[À toi de jouer]
Soit $g$ la fonction définie par $g(x) = \dfrac{x^2 - 4x + 5}{x - 2}$.
\begin{enumerate}
\item Déterminer le domaine de définition $D_g$ et vérifier qu'il est symétrique par rapport à $2$.
\item Montrer que pour tout $x \in D_g$, $g(x) = x - 2 + \dfrac{1}{x - 2}$.
\item En déduire que le point $\Omega(2\,;\,0)$ est centre de symétrie de la courbe de $g$.
\end{enumerate}
\end{exercice}

\begin{corrige}[Exercice : centre de symétrie de la courbe de $g$]
\textbf{1.} $g$ est définie pour $x - 2 \neq 0$, donc $D_g = \mathbb{R} \setminus \{2\}$. Si $2 + h \in D_g$ alors $h \neq 0$, donc $2 - h \in D_g$ : le domaine est symétrique par rapport à $2$.

\textbf{2.} Pour $x \neq 2$ :
\[
x - 2 + \frac{1}{x - 2} = \frac{(x-2)^2 + 1}{x - 2} = \frac{x^2 - 4x + 4 + 1}{x - 2} = \frac{x^2 - 4x + 5}{x - 2} = g(x).
\]

\textbf{3.} Pour $h \neq 0$ :
\[
g(2 + h) = h + \frac{1}{h} \qquad \text{et} \qquad g(2 - h) = -h - \frac{1}{h}.
\]
Donc $\dfrac{g(2+h) + g(2-h)}{2} = \dfrac{0}{2} = 0 = b$. Le point $\Omega(2\,;\,0)$ est bien centre de symétrie de la courbe de $g$.
\end{corrige}

\begin{aretenir}[L'essentiel de la section]
\begin{itemize}
\item $\Omega(a\,;\,b)$ est \textbf{centre de symétrie} de $\mathcal{C}_f$ si $D_f$ est symétrique par rapport à $a$ et si, pour tout $h$ admissible : $f(a + h) + f(a - h) = 2b$, c'est-à-dire $\dfrac{f(a+h)+f(a-h)}{2} = b$.
\item Idée géométrique : $\Omega$ est le \textbf{milieu} de $[MM']$ pour tout couple de points symétriques de la courbe.
\item Cas $a = b = 0$ : on retrouve les \textbf{fonctions impaires}, $f(-x) = -f(x)$.
\item Toute courbe de fonction cubique admet un centre de symétrie : son \textbf{point d'inflexion}, d'abscisse $-\dfrac{\beta}{3\alpha}$ (admis).
\item Ne jamais confondre avec la condition d'axe de symétrie $f(a + h) = f(a - h)$.
\end{itemize}
\end{aretenir}

\coursec{Asymptotes verticales et horizontales}

Dans la section précédente, nous avons appris à repérer les \cle{symétries} d'une courbe : un point ou une droite autour desquels la courbe se « replie » sur elle-même. Nous allons maintenant étudier un autre type de comportement remarquable : celui de la courbe lorsqu'elle s'approche \emph{indéfiniment} d'une droite sans jamais la toucher (ou presque). Ces droites, appelées \cle{asymptotes}, sont des outils précieux : elles guident le tracé de la courbe aux « extrémités » du domaine de définition, là où le tableau de variations ne suffit plus à tout dire.

\courssub{Notion intuitive : la droite dont la courbe se rapproche indéfiniment}

Imagine que tu suives en moto-taxi la route Yaoundé--Douala, et qu'à l'horizon tu aperçoives une ligne droite : plus tu avances, plus ta trajectoire semble se coller à cette ligne, sans jamais la rejoindre complètement. C'est exactement l'idée d'une asymptote : une \cle{droite} dont la courbe se rapproche \emph{aussi près que l'on veut}, lorsque $x$ se rapproche d'une valeur interdite, ou lorsque $x$ devient très grand (vers $+\infty$ ou $-\infty$).

Deux situations se présentent naturellement :

\begin{itemize}
\item la courbe « explose » : ses ordonnées deviennent gigantesques (en valeur absolue) quand $x$ s'approche d'un nombre $a$. La courbe se plaque alors contre la droite \emph{verticale} d'équation $x = a$ ;
\item la courbe « se stabilise » : ses ordonnées se rapprochent d'un nombre fini $b$ quand $x$ tend vers $+\infty$ ou $-\infty$. La courbe se plaque alors contre la droite \emph{horizontale} d'équation $y = b$.
\end{itemize}

Toute la rigueur de la notion tient en un mot : \cle{limite}. Une asymptote n'est jamais décrétée « à l'œil » : elle est \textbf{justifiée par un calcul de limite}. C'est une exigence systématique au Baccalauréat.

\courssub{Asymptote verticale}

\begin{definition}[Asymptote verticale]
Soit $f$ une fonction définie sur un ensemble $D_f$ tel que $a$ est une valeur d'adhérence de $D_f$ (cas typique : $a$ est une borne finie d'un intervalle de $D_f$). On dit que la droite d'équation $x = a$ est une \cle{asymptote verticale} à $\mathcal{C}_f$ lorsque l'une au moins des limites suivantes est infinie :
\[
\lim_{\substack{x \to a \\ x < a}} f(x) = \pm\infty
\qquad\text{ou}\qquad
\lim_{\substack{x \to a \\ x > a}} f(x) = \pm\infty .
\]
Autrement dit : dès que $f(x)$ tend vers $+\infty$ ou $-\infty$ quand $x$ tend vers $a$ (à gauche, à droite, ou des deux côtés), la droite $x = a$ est asymptote verticale.
\end{definition}

\begin{propriete}[Critère de l'asymptote verticale — admis]
Si $\lim\limits_{x \to a^+} f(x) = \pm\infty$ ou $\lim\limits_{x \to a^-} f(x) = \pm\infty$, alors la droite d'équation $x = a$ est asymptote verticale à la courbe de $f$. Graphiquement, la distance entre le point $M(x\,;f(x))$ de la courbe et la droite $x = a$, qui vaut $|x - a|$, tend vers $0$, tandis que le point $M$ « s'éloigne à l'infini » verticalement.
\end{propriete}

\begin{exemplebox}[La fonction inverse décalée]
Étudions $f(x) = \dfrac{1}{x-2}$, définie sur $\mathbb{R}\setminus\{2\}$. La valeur $x = 2$ est interdite : calculons les limites de part et d'autre.
\begin{itemize}
\item Quand $x \to 2$ avec $x > 2$ : le dénominateur $x - 2$ tend vers $0$ en restant \textbf{positif}, donc
\[ \lim_{\substack{x \to 2 \\ x > 2}} \frac{1}{x-2} = +\infty . \]
\item Quand $x \to 2$ avec $x < 2$ : le dénominateur tend vers $0$ en restant \textbf{négatif}, donc
\[ \lim_{\substack{x \to 2 \\ x < 2}} \frac{1}{x-2} = -\infty . \]
\end{itemize}
\textbf{Conclusion :} la droite d'équation $x = 2$ est asymptote verticale à la courbe de $f$. La courbe « monte » vers $+\infty$ à droite de $2$ et « descend » vers $-\infty$ à gauche de $2$.
\end{exemplebox}

\begin{attention}[Une valeur interdite ne donne pas toujours une asymptote !]
C'est l'erreur la plus fréquente au Bac : conclure « $a$ est interdit, donc $x = a$ est asymptote verticale » \textbf{sans calculer la limite}. Faux en général ! Si la limite de $f$ en $a$ est \emph{finie}, il n'y a \textbf{pas} d'asymptote verticale : la courbe présente simplement un « trou » que l'on peut combler (on parle de \cle{prolongement par continuité}). Seule une limite \emph{infinie} donne une asymptote verticale.
\end{attention}

\begin{exemplebox}[Une valeur interdite sans asymptote]
Soit $g(x) = \dfrac{x^2 - 1}{x - 1}$, définie sur $\mathbb{R}\setminus\{1\}$. La valeur $x = 1$ est interdite. Mais pour $x \neq 1$ :
\[
g(x) = \frac{(x-1)(x+1)}{x-1} = x + 1 ,
\qquad\text{donc}\qquad
\lim_{x \to 1} g(x) = 2 \quad (\text{limite finie}).
\]
Il n'y a \textbf{pas} d'asymptote verticale en $x = 1$ : la courbe de $g$ est la droite $y = x + 1$ privée du seul point $(1\,;2)$. Moralité : \emph{on calcule toujours la limite avant de conclure}.
\end{exemplebox}

\courssub{Asymptote horizontale}

\begin{definition}[Asymptote horizontale]
Soit $f$ une fonction définie sur un intervalle du type $[A\,;+\infty[$ (ou $]-\infty\,;A]$). On dit que la droite d'équation $y = b$ (avec $b$ réel) est une \cle{asymptote horizontale} à $\mathcal{C}_f$ en $+\infty$ (respectivement en $-\infty$) lorsque :
\[
\lim_{x \to +\infty} f(x) = b
\qquad\text{(respectivement } \lim_{x \to -\infty} f(x) = b\text{)}.
\]
\end{definition}

\begin{propriete}[Critère de l'asymptote horizontale — admis]
Si $\lim\limits_{x \to +\infty} f(x) = b$ (nombre fini), alors la droite $y = b$ est asymptote horizontale à $\mathcal{C}_f$ en $+\infty$. En effet, la distance verticale entre le point $M(x\,;f(x))$ et la droite $y = b$ vaut $|f(x) - b|$, et cette quantité tend vers $0$ : la courbe se plaque contre la droite. Même conclusion en $-\infty$.
\end{propriete}

\begin{exemplebox}[Une fonction rationnelle]
Soit $f(x) = \dfrac{3x+1}{x-2}$, définie sur $\mathbb{R}\setminus\{2\}$. Pour $x \neq 0$, factorisons par $x$ au numérateur et au dénominateur :
\[
f(x) = \frac{x\left(3 + \dfrac{1}{x}\right)}{x\left(1 - \dfrac{2}{x}\right)} = \frac{3 + \dfrac{1}{x}}{1 - \dfrac{2}{x}} .
\]
Comme $\lim\limits_{x \to +\infty} \dfrac{1}{x} = 0$ et $\lim\limits_{x \to +\infty} \dfrac{2}{x} = 0$, on obtient :
\[
\lim_{x \to +\infty} f(x) = \frac{3}{1} = 3 .
\]
\textbf{Conclusion :} la droite d'équation $y = 3$ est asymptote horizontale à la courbe en $+\infty$ (le même calcul en $-\infty$ donne aussi $3$).
\end{exemplebox}

\begin{popfigure}
\begin{center}
\begin{tikzpicture}
\begin{axis}[
  width=0.92\linewidth, height=8cm,
  axis lines=middle, xlabel={$x$}, ylabel={$y$},
  xmin=-6, xmax=10, ymin=-6, ymax=10,
  xtick={-4,-2,2,4,6,8}, ytick={-4,-2,2,4,6,8},
  grid=none, samples=200, clip=true]
\addplot[popcurve,domain=-6:1.85]{(3*x+1)/(x-2)};
\addplot[popcurve,domain=2.15:10]{(3*x+1)/(x-2)};
\addplot[poporange,dashed,thick] coordinates {(2,-6) (2,10)};
\addplot[popgreen,dashed,thick,domain=-6:10]{3};
\node[poporange,anchor=west] at (axis cs:2.3,-4.5) {$x=2$};
\node[popgreen,anchor=south west] at (axis cs:6.2,3.15) {$y=3$};
\node[popblue,anchor=south] at (axis cs:-3,1.4) {$\mathcal{C}_f$};
\end{axis}
\end{tikzpicture}
\end{center}
\legende{Courbe de $f(x)=\dfrac{3x+1}{x-2}$ : asymptote verticale $x=2$ (orange) et asymptote horizontale $y=3$ (vert).}
\end{popfigure}

\courssub{Position relative de la courbe et de son asymptote horizontale}

Savoir que $y = b$ est asymptote ne suffit pas toujours : le correcteur du Bac apprécie que tu précises si la courbe est \emph{au-dessus} ou \emph{en dessous} de son asymptote. L'outil est simple : étudier le \cle{signe de la différence} $f(x) - b$.

\begin{methode}[Position de la courbe par rapport à l'asymptote $y = b$]
\begin{enumerate}
\item Calculer la différence $f(x) - b$ et la simplifier (réduire au même dénominateur).
\item Étudier le signe de cette différence sur l'intervalle considéré.
\item Conclure :
\begin{itemize}
\item si $f(x) - b > 0$, la courbe est \textbf{au-dessus} de l'asymptote ;
\item si $f(x) - b < 0$, la courbe est \textbf{en dessous} de l'asymptote.
\end{itemize}
\end{enumerate}
\end{methode}

\begin{exemplebox}[Position relative pour $f(x) = \dfrac{3x+1}{x-2}$]
Calculons :
\[
f(x) - 3 = \frac{3x+1}{x-2} - 3 = \frac{3x + 1 - 3(x-2)}{x-2} = \frac{7}{x-2}.
\]
Le numérateur $7$ est strictement positif. Donc :
\begin{itemize}
\item si $x > 2$, alors $f(x) - 3 > 0$ : la courbe est \textbf{au-dessus} de l'asymptote $y = 3$ ;
\item si $x < 2$, alors $f(x) - 3 < 0$ : la courbe est \textbf{en dessous} de l'asymptote $y = 3$.
\end{itemize}
C'est exactement ce que l'on observe sur la figure ci-dessus : à droite, la branche « descend » vers $y = 3$ par le haut ; à gauche, la branche « monte » vers $y = 3$ par le bas.
\end{exemplebox}

\courssub{Le cas des fonctions rationnelles}

Les fonctions rationnelles (quotients de polynômes) sont les grandes vedettes de cette leçon, car elles fournissent l'essentiel des exemples au Baccalauréat.

\begin{propriete}[Asymptotes d'une fonction rationnelle]
Soit $f(x) = \dfrac{P(x)}{Q(x)}$ une fonction rationnelle.
\begin{itemize}
\item \textbf{Asymptotes verticales :} si $a$ annule le dénominateur ($Q(a) = 0$) \textbf{sans} annuler le numérateur ($P(a) \neq 0$), alors $\lim\limits_{x \to a} f(x) = \pm\infty$ et la droite $x = a$ est asymptote verticale.
\item \textbf{Asymptotes horizontales :} en $+\infty$ et $-\infty$, la limite d'une fonction rationnelle est la limite du quotient des \cle{termes de plus haut degré}. Si le degré du numérateur est inférieur ou égal à celui du dénominateur, cette limite est finie et donne une asymptote horizontale.
\end{itemize}
\end{propriete}

\begin{attention}[Numérateur et dénominateur nuls en même temps]
Si $P(a) = 0$ \textbf{et} $Q(a) = 0$, on ne peut rien conclure directement : il faut \textbf{factoriser} par $(x - a)$ en haut et en bas, simplifier, puis recalculer la limite. C'est exactement ce qui s'est passé avec $g(x) = \dfrac{x^2-1}{x-1}$ : après simplification, la limite en $1$ valait $2$ (finie), donc pas d'asymptote. Ne jamais oublier cette étape de factorisation !
\end{attention}

\begin{methode}[Recherche systématique des asymptotes verticales et horizontales]
\begin{enumerate}
\item Déterminer l'ensemble de définition $D_f$ et repérer les \cle{valeurs interdites} (bornes finies des intervalles de $D_f$).
\item En chaque valeur interdite $a$, calculer les limites à gauche et à droite de $f$ (en factorisant si nécessaire). Limite infinie $\Rightarrow$ asymptote verticale $x = a$ ; limite finie $\Rightarrow$ pas d'asymptote.
\item Si $D_f$ contient un voisinage de $+\infty$ ou $-\infty$, calculer les limites en ces bornes. Limite finie $b$ $\Rightarrow$ asymptote horizontale $y = b$.
\item Préciser, si demandé, la position de la courbe par rapport à chaque asymptote horizontale (signe de $f(x) - b$).
\end{enumerate}
\end{methode}

\begin{popfigure}
\begin{center}
\begin{tikzpicture}[
  box/.style={draw=popblue, thick, rounded corners, fill=popblueL, text width=4.0cm, align=center, minimum height=1.5cm, font=\small},
  res/.style={draw=poporange, thick, rounded corners, fill=poporangeL, text width=4.0cm, align=center, minimum height=1.5cm, font=\small},
  arr/.style={-{Stealth[length=3mm]}, thick, popdark}]
\node[box, align=center] (a) at (0,2.2) {$x \to a$ (réel) \\ $\lim f(x) = \pm\infty$};
\node[box, align=center] (b) at (7.5,2.2) {$x \to a$ (réel) \\ $\lim f(x) = \ell$ finie};
\node[box, align=center] (c) at (0,-2.2) {$x \to \pm\infty$ \\ $\lim f(x) = b$ finie};
\node[box, align=center] (d) at (7.5,-2.2) {$x \to \pm\infty$ \\ $\lim f(x) = \pm\infty$};
\node[res, align=center] (ra) at (0,0) {Asymptote \textbf{verticale} \\ $x = a$};
\node[res, align=center] (rb) at (7.5,0) {\textbf{Pas} d'asymptote \\ (trou ou prolongement)};
\node[res, align=center] (rc) at (0,-4.4) {Asymptote \textbf{horizontale} \\ $y = b$};
\node[res, align=center] (rd) at (7.5,-4.4) {Pas d'asymptote horizontale \\ (voir asymptote oblique, section 5)};
\draw[arr] (a) -- (ra);
\draw[arr] (b) -- (rb);
\draw[arr] (c) -- (rc);
\draw[arr] (d) -- (rd);
\end{tikzpicture}
\end{center}
\legende{Les quatre situations de limites et la conclusion graphique correspondante.}
\end{popfigure}

\begin{exoresolu}[Étude complète des asymptotes de $f(x) = \dfrac{2x-3}{x+1}$]
\textbf{Énoncé.} Soit $f(x) = \dfrac{2x-3}{x+1}$.
\begin{enumerate}
\item Déterminer l'ensemble de définition de $f$.
\item Montrer que la courbe $\mathcal{C}_f$ admet une asymptote verticale et une asymptote horizontale dont on donnera les équations.
\item Préciser la position de $\mathcal{C}_f$ par rapport à son asymptote horizontale.
\end{enumerate}
\tcblower
\textbf{Solution détaillée.}
\begin{enumerate}
\item $f$ est définie lorsque $x + 1 \neq 0$, c'est-à-dire $x \neq -1$. Donc $D_f = \mathbb{R}\setminus\{-1\} = \left]-\infty\,;-1\right[ \cup \left]-1\,;+\infty\right[$.
\item \textbf{Asymptote verticale.} En $x = -1$, le numérateur vaut $2\times(-1) - 3 = -5 \neq 0$ et le dénominateur tend vers $0$.
\begin{itemize}
\item Si $x \to -1$ avec $x > -1$, alors $x + 1 \to 0^{+}$, donc $\lim\limits_{\substack{x \to -1 \\ x > -1}} f(x) = \dfrac{-5}{0^{+}} = -\infty$.
\item Si $x \to -1$ avec $x < -1$, alors $x + 1 \to 0^{-}$, donc $\lim\limits_{\substack{x \to -1 \\ x < -1}} f(x) = \dfrac{-5}{0^{-}} = +\infty$.
\end{itemize}
Les limites étant infinies, la droite d'équation $\boxed{x = -1}$ est asymptote verticale à $\mathcal{C}_f$.

\textbf{Asymptote horizontale.} Pour $x \neq 0$ :
\[
f(x) = \frac{x\left(2 - \dfrac{3}{x}\right)}{x\left(1 + \dfrac{1}{x}\right)} = \frac{2 - \dfrac{3}{x}}{1 + \dfrac{1}{x}}
\ \xrightarrow[x \to \pm\infty]{}\ \frac{2}{1} = 2 .
\]
La limite est finie, donc la droite d'équation $\boxed{y = 2}$ est asymptote horizontale à $\mathcal{C}_f$ en $+\infty$ et en $-\infty$.
\item Étudions le signe de la différence :
\[
f(x) - 2 = \frac{2x - 3}{x+1} - 2 = \frac{2x - 3 - 2(x+1)}{x+1} = \frac{-5}{x+1}.
\]
Le numérateur $-5$ est strictement négatif. Donc :
\begin{itemize}
\item sur $]-1\,;+\infty[$, $x + 1 > 0$, donc $f(x) - 2 < 0$ : la courbe est \textbf{en dessous} de l'asymptote $y = 2$ ;
\item sur $]-\infty\,;-1[$, $x + 1 < 0$, donc $f(x) - 2 > 0$ : la courbe est \textbf{au-dessus} de l'asymptote $y = 2$.
\end{itemize}
\end{enumerate}
\end{exoresolu}

\begin{savaistu}[Les asymptotes dans la vie réelle]
Les asymptotes horizontales modélisent des phénomènes de \emph{saturation} : la charge d'un téléphone branché au secteur se rapproche de $100\,\%$ sans jamais dépasser ce plafond ; le nombre de clients d'un réseau mobile comme MTN ou Orange dans une ville tend vers une limite quand le marché est saturé. Les asymptotes verticales, elles, modélisent des « explosions » : l'intensité du courant dans un circuit dont la résistance tend vers zéro, ou le temps d'attente dans un guichet CAMTEL quand le nombre d'employés tend vers zéro. Les mathématiques décrivent ces comportements extrêmes avec une précision remarquable !
\end{savaistu}

\begin{aretenir}[L'essentiel de la section]
\begin{itemize}
\item $\lim\limits_{x \to a} f(x) = \pm\infty$ (à gauche ou à droite) $\Longleftrightarrow$ la droite $x = a$ est \cle{asymptote verticale}.
\item $\lim\limits_{x \to \pm\infty} f(x) = b$ (fini) $\Longleftrightarrow$ la droite $y = b$ est \cle{asymptote horizontale}.
\item Une asymptote se justifie \textbf{toujours} par un calcul de limite écrit proprement.
\item Une valeur interdite ne donne une asymptote verticale que si la limite est \emph{infinie} ; si elle est finie, on factorise et on conclut à un prolongement par continuité.
\item La position de la courbe par rapport à $y = b$ s'obtient par le signe de $f(x) - b$.
\end{itemize}
\end{aretenir}

\begin{exercice}[Pour t'entraîner]
Soit $h(x) = \dfrac{x + 4}{2x - 6}$.
\begin{enumerate}
\item Déterminer l'ensemble de définition de $h$.
\item Justifier l'existence d'une asymptote verticale et d'une asymptote horizontale, et donner leurs équations.
\item Étudier la position de la courbe de $h$ par rapport à son asymptote horizontale.
\end{enumerate}
\end{exercice}

\begin{corrige}[Exercice : pour t'entraîner]
\begin{enumerate}
\item $2x - 6 = 0 \iff x = 3$, donc $D_h = \mathbb{R}\setminus\{3\}$.
\item En $x = 3$ : le numérateur vaut $3 + 4 = 7 \neq 0$. Si $x \to 3^{+}$, $2x - 6 \to 0^{+}$ donc $h(x) \to +\infty$ ; si $x \to 3^{-}$, $2x - 6 \to 0^{-}$ donc $h(x) \to -\infty$. La droite $x = 3$ est asymptote verticale. En $\pm\infty$ : $h(x) = \dfrac{1 + \frac{4}{x}}{2 - \frac{6}{x}} \to \dfrac{1}{2}$. La droite $y = \dfrac{1}{2}$ est asymptote horizontale en $+\infty$ et en $-\infty$.
\item $h(x) - \dfrac{1}{2} = \dfrac{x + 4}{2x - 6} - \dfrac{1}{2} = \dfrac{2(x+4) - (2x-6)}{2(2x-6)} = \dfrac{14}{4x - 12} = \dfrac{7}{2x - 6}$. Sur $]3\,;+\infty[$, $h(x) - \frac{1}{2} > 0$ : courbe au-dessus. Sur $]-\infty\,;3[$, $h(x) - \frac{1}{2} < 0$ : courbe en dessous.
\end{enumerate}
\end{corrige}

\coursec{Asymptotes obliques}

Dans la section précédente, nous avons appris à repérer et à justifier les \cle{asymptotes verticales} (quand $f(x)$ explose vers $\pm\infty$ près d'une valeur interdite $a$) et les \cle{asymptotes horizontales} (quand $f(x)$ se stabilise vers un nombre $b$ quand $x$ tend vers $\pm\infty$). Mais il arrive qu'en $+\infty$ ou en $-\infty$, la courbe ne se plaque ni contre une droite verticale, ni contre une droite horizontale : elle suit une droite \emph{inclinée}, comme une route qui longe une voie ferrée en pente. C'est le phénomène d'\cle{asymptote oblique}, que nous étudions maintenant.

\courssub{Motivation : une courbe qui suit une droite inclinée}

Considérons la fonction définie sur $\mathbb{R}^*$ par
\[ f(x) = \frac{x^2 + x + 1}{x}. \]
Quand $x$ devient très grand, le terme $x^2$ domine le numérateur : $f(x)$ tend vers $+\infty$. Il n'y a donc \emph{pas} d'asymptote horizontale en $+\infty$. Pourtant, si l'on trace la courbe sur un grand intervalle, on observe qu'elle se rapproche visiblement d'une droite montante. Laquelle ? En écrivant
\[ f(x) = \frac{x^2}{x} + \frac{x}{x} + \frac{1}{x} = x + 1 + \frac{1}{x}, \]
on voit apparaître la droite d'équation $y = x + 1$, et un \emph{reste} $\dfrac{1}{x}$ qui devient négligeable quand $x$ grandit. La courbe « colle » de plus en plus à la droite $y = x + 1$ : c'est exactement le comportement d'une asymptote oblique.

Remarquons au passage que $\lim\limits_{x \to 0^+} f(x) = +\infty$ et $\lim\limits_{x \to 0^-} f(x) = -\infty$ : la droite $x = 0$ (l'axe des ordonnées) est aussi asymptote \emph{verticale} à la courbe. Une même courbe peut donc posséder plusieurs asymptotes de natures différentes.

\begin{popfigure}
\begin{center}
\begin{tikzpicture}
\begin{axis}[
    width=0.92\linewidth, height=8cm,
    axis lines=middle,
    xlabel={$x$}, ylabel={$y$},
    xmin=-6, xmax=6, ymin=-8, ymax=10,
    xtick={-5,-4,...,5}, ytick={-6,-4,...,8},
    grid=both, grid style={popline!40},
    legend style={at={(0.02,0.98)}, anchor=north west, font=\small}
]
\addplot[popcurve,domain=-6:-0.35,samples=150]{(x^2+x+1)/x};
\addplot[popcurve,domain=0.35:6,samples=150]{(x^2+x+1)/x};
\addplot[poporange,thick,dashed,domain=-5.5:5.5,samples=2]{x+1};
\addplot[popgreen,thick,dashed] coordinates {(0,-8) (0,10)};
\legend{$\mathcal{C}_f$,\; $y=x+1$ (asymptote oblique),\; $x=0$ (asymptote verticale)}
\end{axis}
\end{tikzpicture}
\end{center}
\legende{La courbe de $f(x)=\dfrac{x^2+x+1}{x}$ se rapproche de la droite $y=x+1$ quand $x\to\pm\infty$, et de la droite $x=0$ près de $0$.}
\end{popfigure}

\courssub{Définition et caractérisation}

\begin{definition}[Asymptote oblique]
Soit $f$ une fonction définie au voisinage de $+\infty$ (respectivement $-\infty$) et $(d)$ la droite d'équation $y = ax + b$ avec $a \neq 0$.
On dit que $(d)$ est \cle{asymptote oblique} à la courbe $\mathcal{C}_f$ en $+\infty$ (respectivement en $-\infty$) lorsque
\[ \lim_{x \to +\infty} \bigl[ f(x) - (ax + b) \bigr] = 0
\quad \text{(respectivement } \lim_{x \to -\infty} \bigl[ f(x) - (ax + b) \bigr] = 0 \text{)}. \]
\end{definition}

\begin{attention}[Pourquoi $a \neq 0$ ?]
La condition $a \neq 0$ distingue l'asymptote \emph{oblique} de l'asymptote \emph{horizontale} : si $a = 0$, la droite $y = b$ est horizontale et l'on retrouve la définition déjà vue. Le vocabulaire compte au Bac : annoncer « asymptote oblique $y = 3$ » est une erreur, car $y = 3$ est une asymptote \emph{horizontale}.
\end{attention}

\begin{propriete}[Caractérisation par l'écart]
La droite $(d) : y = ax + b$ ($a \neq 0$) est asymptote oblique à $\mathcal{C}_f$ en $+\infty$ si et seulement si l'écart vertical $f(x) - (ax+b)$ tend vers $0$ quand $x \to +\infty$. De même en $-\infty$.
\end{propriete}

Cette propriété est admise, mais elle est très intuitive : pour une abscisse $x$ donnée, le point $M(x\,;\,f(x))$ de la courbe et le point $N(x\,;\,ax+b)$ de la droite ont la même abscisse. Leur écart vertical est $MN = \bigl| f(x) - (ax+b) \bigr|$. Dire que cet écart tend vers $0$, c'est dire que la distance entre la courbe et la droite s'écrase : la courbe épouse la droite. C'est exactement la même idée que pour l'asymptote horizontale, sauf que la droite de référence est inclinée.

\begin{popfigure}
\begin{center}
\begin{tikzpicture}
\begin{axis}[
    width=0.92\linewidth, height=7cm,
    axis lines=middle,
    xlabel={$x$}, ylabel={$y$},
    xmin=0, xmax=12, ymin=0, ymax=14,
    xtick={1,2,...,11}, ytick={2,4,...,12},
    grid=both, grid style={popline!40},
    legend style={at={(0.02,0.98)}, anchor=north west, font=\small}
]
\addplot[popcurve,domain=0.4:12,samples=200]{x+1+1/x};
\addplot[poporange,thick,dashed,domain=0:12,samples=2]{x+1};
\legend{$\mathcal{C}_f$,\; $y=x+1$}
\addplot[popgreen,thick,<->] coordinates {(2,3.5) (2,3)};
\addplot[popgreen,thick,<->] coordinates {(8,9.125) (8,9)};
\node[popgreen,font=\small] at (axis cs:2.7,3.25) {$\frac12$};
\node[popgreen,font=\small] at (axis cs:9,9.06) {$\frac18$};
\end{axis}
\end{tikzpicture}
\end{center}
\legende{Zoom sur l'écart $f(x)-(x+1)=\dfrac{1}{x}$ : il vaut $\tfrac12$ en $x=2$, $\tfrac18$ en $x=8$, et tend vers $0$.}
\end{popfigure}

\courssub{Méthode 1 : la division euclidienne (fractions rationnelles)}

\begin{methode}[Faire apparaître $ax + b + g(x)$]
Pour montrer qu'une fraction rationnelle $f(x) = \dfrac{P(x)}{D(x)}$ admet une asymptote oblique :
\begin{enumerate}
\item Vérifier que le degré du numérateur est exactement le degré du dénominateur \emph{plus un}.
\item Effectuer la \cle{division euclidienne} de $P(x)$ par $D(x)$ (ou réorganiser le numérateur) pour écrire
\[ f(x) = ax + b + \frac{R(x)}{D(x)}, \quad \text{avec } \deg R < \deg D. \]
\item Poser $g(x) = \dfrac{R(x)}{D(x)}$ et montrer que $\lim\limits_{x \to \pm\infty} g(x) = 0$ (une fraction dont le numérateur est de degré strictement inférieur à celui du dénominateur tend toujours vers $0$ en $\pm\infty$).
\item Conclure : la droite $y = ax + b$ est asymptote oblique à $\mathcal{C}_f$.
\end{enumerate}
\end{methode}

\begin{exemplebox}[Exemple de référence]
Soit $f(x) = \dfrac{x^2 + x + 1}{x}$ sur $\mathbb{R}^*$. Le numérateur est de degré $2$, le dénominateur de degré $1$ : on est bien dans le cas « degré $+1$ ». En divisant chaque terme du numérateur par $x$ :
\[ f(x) = x + 1 + \frac{1}{x}. \]
Ici $g(x) = \dfrac{1}{x}$ et $\lim\limits_{x \to +\infty} \dfrac{1}{x} = 0$, de même en $-\infty$. Donc la droite $y = x + 1$ est asymptote oblique à $\mathcal{C}_f$ en $+\infty$ \emph{et} en $-\infty$.
\end{exemplebox}

\courssub{Méthode 2 : conjecture graphique puis vérification}

Au Baccalauréat, on te donne souvent la courbe ou on te demande de conjecturer l'asymptote. La démarche est alors :

\begin{enumerate}
\item Sur le graphique, repérer la droite inclinée que la courbe semble suivre en $\pm\infty$ ; lire son coefficient directeur $a$ et son ordonnée à l'origine $b$.
\item \emph{Vérifier par le calcul} que $\lim\limits_{x \to +\infty} \bigl[ f(x) - (ax+b) \bigr] = 0$.
\end{enumerate}

\begin{attention}[Conjecturer ne suffit pas !]
L'erreur la plus fréquente est de conclure « la droite $y = ax+b$ est asymptote » uniquement parce que la courbe en a l'air sur le graphique, ou parce que $f(x)$ et $ax+b$ tendent tous deux vers $+\infty$. Deux quantités qui tendent vers $+\infty$ peuvent s'écarter indéfiniment : par exemple $f(x) = x + \sqrt{x}$ et la droite $y = x$ s'écartent de $\sqrt{x}$, qui tend vers $+\infty$. \textbf{Il faut toujours calculer la limite de l'écart} $f(x) - (ax+b)$ et vérifier qu'elle vaut $0$.
\end{attention}

\courssub{Position relative de la courbe et de l'asymptote}

Une fois l'asymptote oblique $y = ax + b$ établie, le signe de l'écart $f(x) - (ax+b)$ donne la position de la courbe :

\begin{itemize}
\item si $f(x) - (ax+b) > 0$ sur un intervalle, la courbe est \cle{au-dessus} de l'asymptote ;
\item si $f(x) - (ax+b) < 0$, la courbe est \cle{en dessous} ;
\item si l'écart change de signe en s'annulant, la courbe \emph{croise} son asymptote (c'est autorisé : une asymptote n'est pas une barrière !).
\end{itemize}

Pour $f(x) = \dfrac{x^2+x+1}{x}$, l'écart est $\dfrac{1}{x}$ : positif pour $x > 0$ (courbe au-dessus de $y = x+1$), négatif pour $x < 0$ (courbe en dessous). C'est exactement ce qu'on observe sur la première figure.

\begin{exoresolu}[Asymptote oblique et position relative]
Soit $g$ la fonction définie sur $\mathbb{R} \setminus \{-1\}$ par
\[ g(x) = 2x - 1 + \frac{3}{x+1}. \]
\begin{enumerate}
\item Montrer que la droite $(\Delta) : y = 2x - 1$ est asymptote oblique à $\mathcal{C}_g$ en $+\infty$ et en $-\infty$.
\item Étudier la position de $\mathcal{C}_g$ par rapport à $(\Delta)$.
\end{enumerate}
\tcblower
\textbf{1.} Calculons l'écart entre la courbe et la droite :
\[ g(x) - (2x - 1) = \frac{3}{x+1}. \]
Or $\lim\limits_{x \to +\infty} (x+1) = +\infty$, donc $\lim\limits_{x \to +\infty} \dfrac{3}{x+1} = 0$. De même, $\lim\limits_{x \to -\infty} (x+1) = -\infty$, donc $\lim\limits_{x \to -\infty} \dfrac{3}{x+1} = 0$. Ainsi
\[ \lim_{x \to +\infty} \bigl[ g(x) - (2x-1) \bigr] = 0
\quad \text{et} \quad
\lim_{x \to -\infty} \bigl[ g(x) - (2x-1) \bigr] = 0. \]
La droite $(\Delta) : y = 2x - 1$ est donc asymptote oblique à $\mathcal{C}_g$ en $+\infty$ et en $-\infty$.

\textbf{2.} Le signe de l'écart $\dfrac{3}{x+1}$ est celui de $x+1$ :
\begin{itemize}
\item si $x > -1$, alors $g(x) - (2x-1) > 0$ : la courbe est \textbf{au-dessus} de $(\Delta)$ ;
\item si $x < -1$, alors $g(x) - (2x-1) < 0$ : la courbe est \textbf{en dessous} de $(\Delta)$.
\end{itemize}
\end{exoresolu}

\courssub{Complément utile au Bac : trouver $a$ et $b$ par le calcul}

\begin{savaistu}[Hors strict programme, mais précieux]
Quand la forme $ax + b + g(x)$ n'est pas évidente (par exemple avec des racines carrées), on peut déterminer l'asymptote oblique en deux étapes :
\begin{enumerate}
\item le coefficient directeur : $\displaystyle a = \lim_{x \to +\infty} \frac{f(x)}{x}$ (cette limite doit être un réel \emph{fini non nul}) ;
\item l'ordonnée à l'origine : $\displaystyle b = \lim_{x \to +\infty} \bigl[ f(x) - ax \bigr]$ (cette limite doit être un réel \emph{fini}).
\end{enumerate}
Si l'une des deux limites est infinie ou n'existe pas, il n'y a pas d'asymptote oblique. Par exemple, pour $f(x) = \dfrac{x^2+x+1}{x}$ : $\dfrac{f(x)}{x} = 1 + \dfrac{1}{x} + \dfrac{1}{x^2} \to 1$, donc $a = 1$ ; puis $f(x) - x = 1 + \dfrac{1}{x} \to 1$, donc $b = 1$. On retrouve $y = x + 1$.
\end{savaistu}

\begin{aretenir}[L'essentiel sur les asymptotes obliques]
\begin{itemize}
\item $(d) : y = ax + b$ ($a \neq 0$) est asymptote oblique à $\mathcal{C}_f$ en $+\infty$ si et seulement si $\lim\limits_{x \to +\infty} \bigl[ f(x) - (ax+b) \bigr] = 0$ (idem en $-\infty$).
\item Pour une fraction rationnelle, on obtient $a$ et $b$ par \textbf{division euclidienne} : $f(x) = ax + b + g(x)$ avec $g(x) \to 0$.
\item Une conjecture graphique doit toujours être \textbf{confirmée par le calcul de la limite de l'écart}.
\item Le \textbf{signe} de $f(x) - (ax+b)$ donne la position de la courbe par rapport à l'asymptote ; la courbe peut la croiser.
\end{itemize}
\end{aretenir}

\begin{exercice}[Pour t'entraîner]
Soit $h$ définie sur $\mathbb{R} \setminus \{2\}$ par $h(x) = \dfrac{x^2 - 3x + 5}{x - 2}$.
\begin{enumerate}
\item Vérifier que $h(x) = x - 1 + \dfrac{3}{x-2}$.
\item En déduire que $\mathcal{C}_h$ admet une asymptote oblique dont on donnera l'équation.
\item Préciser la position de $\mathcal{C}_h$ par rapport à cette asymptote.
\end{enumerate}
\end{exercice}

\begin{corrige}[Exercice : asymptote oblique de $h$]
\textbf{1.} On réduit au même dénominateur :
\[ x - 1 + \frac{3}{x-2} = \frac{(x-1)(x-2) + 3}{x-2} = \frac{x^2 - 3x + 2 + 3}{x-2} = \frac{x^2 - 3x + 5}{x-2} = h(x). \]

\textbf{2.} L'écart vaut $h(x) - (x-1) = \dfrac{3}{x-2}$, et $\lim\limits_{x \to \pm\infty} \dfrac{3}{x-2} = 0$. Donc la droite $y = x - 1$ est asymptote oblique à $\mathcal{C}_h$ en $+\infty$ et en $-\infty$.

\textbf{3.} L'écart $\dfrac{3}{x-2}$ est positif si $x > 2$ (courbe au-dessus) et négatif si $x < 2$ (courbe en dessous).
\end{corrige}

Tu sais désormais détecter et justifier les trois types d'asymptotes : verticales, horizontales et obliques. Dans la prochaine section, nous assemblerons tous ces outils — tableau de variations, symétries, asymptotes — pour construire une courbe complète et soignée, comme on te le demandera au Baccalauréat.

\coursec{Synthèse : construire une courbe complète}

Dans la section précédente, tu as appris à détecter les asymptotes obliques : quand $f(x) - (ax+b)$ tend vers $0$, la droite d'équation $y = ax + b$ « aspire » la courbe à l'infini. Tu disposes maintenant de toutes les pièces du puzzle : limites, asymptotes, dérivée, variations, symétries. Il reste à les assembler dans le bon ordre. C'est exactement ce que demande l'épreuve du Baccalauréat A : à partir d'une fonction donnée, produire une \cle{courbe complète}, cohérente et soignée. Cette section te donne la méthode, puis la met en œuvre sur un exemple fil rouge étudié de bout en bout.

\courssub{Le plan d'étude complet d'une fonction}

Imaginons que tu doives décrire un paysage à quelqu'un qui ne le verra jamais. Tu commencerais par les limites du terrain (le domaine), puis les grandes directions (les asymptotes), puis le relief (les variations), puis les points remarquables (sommets, creux), avant de dessiner la carte. L'étude d'une fonction suit exactement cette logique.

\begin{methode}[Les 6 étapes de l'étude complète d'une fonction au Bac A]
\begin{enumerate}
\item \textbf{Domaine de définition.} Déterminer $D_f$ en excluant les valeurs qui annulent un dénominateur ou rendent négatif un radicande. Repérer immédiatement une éventuelle \cle{parité} ou \cle{symétrie} : cela peut diviser le travail par deux.
\item \textbf{Limites aux bornes et asymptotes.} Calculer les limites de $f$ aux bornes ouvertes de $D_f$ (y compris en $+\infty$ et $-\infty$). Chaque limite infinie en un point fini $a$ donne une \cle{asymptote verticale} $x = a$ ; chaque limite finie $b$ en l'infini donne une \cle{asymptote horizontale} $y = b$ ; sinon, chercher une \cle{asymptote oblique} $y = ax+b$ en étudiant $f(x) - (ax+b)$.
\item \textbf{Dérivée et signe de la dérivée.} Calculer $f'$, factoriser si possible, et résoudre $f'(x) > 0$ et $f'(x) < 0$.
\item \textbf{Tableau de variation enrichi.} Y faire figurer : les valeurs interdites (double barre), les limites aux bornes, les extremums avec leurs valeurs exactes.
\item \textbf{Éléments graphiques.} Placer les asymptotes en pointillés, les points remarquables (extremums, intersections avec les axes), les tangentes horizontales aux points où $f'$ s'annule.
\item \textbf{Tracé et vérification.} Tracer la courbe branche par branche en suivant les asymptotes, puis vérifier la cohérence globale : la courbe respecte-t-elle le tableau de variation ? rejoint-elle bien ses asymptotes ?
\end{enumerate}
\end{methode}

\begin{attention}[La courbe ne traverse jamais son asymptote verticale]
Deux erreurs reviennent chaque année dans les copies :
\begin{itemize}
\item Tracer une branche qui \textbf{traverse} l'asymptote verticale $x = a$. C'est impossible : $f$ n'est pas définie en $a$, la courbe ne peut pas exister des deux côtés « collée » à la droite en un même point. Chaque branche doit \emph{fuir} vers $+\infty$ ou $-\infty$ en se rapprochant de la droite.
\item Tracer une branche qui \textbf{s'éloigne} de son asymptote (horizontale ou oblique) quand $x$ grandit. C'est le contraire de la définition : l'écart $f(x) - (ax+b)$ doit \emph{diminuer} vers $0$. La courbe doit « coller » de plus en plus à l'asymptote.
\end{itemize}
En revanche, une courbe \emph{peut} traverser son asymptote horizontale ou oblique en un point fini : l'asymptote ne décrit que le comportement à l'infini.
\end{attention}

\courssub{Réduire le domaine d'étude grâce aux symétries}

Voici un réflexe qui fait gagner un temps précieux le jour de l'épreuve. Si la courbe admet un centre ou un axe de symétrie, il est inutile d'étudier la fonction sur tout son domaine : la moitié suffit.

\begin{propriete}[Réduction du domaine d'étude par symétrie]
Soit $f$ une fonction de domaine $D_f$ et $\Omega(a\,;\,b)$ un point tel que, pour tout $x$ avec $a+x \in D_f$ et $a-x \in D_f$ :
\[ f(a+x) + f(a-x) = 2b. \]
Alors $\Omega$ est \cle{centre de symétrie} de la courbe $\mathcal{C}_f$, et il suffit d'étudier $f$ sur la moitié du domaine (par exemple $x \geq a$) : on obtient l'autre moitié de la courbe par la symétrie de centre $\Omega$.
\end{propriete}

\begin{experience}[Pourquoi la moitié du domaine suffit]
Rappelons la démonstration vue dans la section 3. Soit $M(x\,;\,f(x))$ un point de $\mathcal{C}_f$ avec $x = a + t$, $t \geq 0$. Son symétrique par rapport à $\Omega(a\,;\,b)$ est le point $M'(2a - x\,;\,2b - f(x))$, c'est-à-dire $M'(a - t\,;\,2b - f(a+t))$. Or la relation $f(a+t) + f(a-t) = 2b$ donne $f(a-t) = 2b - f(a+t)$. Donc $M' = (a-t\,;\,f(a-t))$ est bien un point de $\mathcal{C}_f$. La courbe est donc globalement invariante par la symétrie de centre $\Omega$ : toute l'information graphique pour $x \leq a$ est contenue dans celle pour $x \geq a$. \hfill$\square$
\end{experience}

Le cas le plus fréquent au Bac A est celui d'une fonction \cle{impaire} ($f(-x) = -f(x)$, domaine symétrique par rapport à $0$) : l'origine est centre de symétrie, on étudie sur $[0\,;\,+\infty[$ et on complète par symétrie. De même, si $f(a+x) = f(a-x)$, la droite $x = a$ est axe de symétrie et on étudie seulement pour $x \geq a$.

\courssub{Exemple chiffré : étude complète de $f(x) = x + \dfrac{1}{x}$}

\begin{exemplebox}[Étude guidée d'une fonction impaire]
Soit $f(x) = x + \dfrac{1}{x}$. Appliquons les 6 étapes.

\textbf{1. Domaine.} Le dénominateur s'annule en $0$ : $D_f = \mathbb{R}^* = ]-\infty\,;\,0[ \cup ]0\,;\,+\infty[$. De plus $f(-x) = -x - \dfrac{1}{x} = -f(x)$ : $f$ est \cle{impaire}. On étudie sur $]0\,;\,+\infty[$ et on complétera par symétrie centrale de centre $O$.

\textbf{2. Limites et asymptotes.}
\[ \lim_{x \to 0^+} f(x) = 0 + (+\infty) = +\infty \quad\Longrightarrow\quad \text{asymptote verticale } x = 0. \]
\[ \lim_{x \to +\infty} \left[f(x) - x\right] = \lim_{x \to +\infty} \frac{1}{x} = 0 \quad\Longrightarrow\quad \text{asymptote oblique } y = x. \]

\textbf{3. Dérivée.} $f'(x) = 1 - \dfrac{1}{x^2} = \dfrac{x^2 - 1}{x^2}$. Sur $]0\,;\,+\infty[$ : $f'(x) < 0$ sur $]0\,;\,1[$ et $f'(x) > 0$ sur $]1\,;\,+\infty[$.

\textbf{4. Variations.} $f$ décroît de $+\infty$ à $f(1) = 2$ sur $]0\,;\,1]$, puis croît de $2$ vers $+\infty$ sur $[1\,;\,+\infty[$. Minimum local : $(1\,;\,2)$, avec tangente horizontale.

\textbf{5 et 6. Tracé.} On place les pointillés $x = 0$ et $y = x$, le point $(1\,;\,2)$, on trace la branche sur $]0\,;\,+\infty[$, puis on complète par symétrie de centre $O$ : le point $(-1\,;\,-2)$ est maximum local de la branche de gauche.
\end{exemplebox}

\begin{popfigure}
\begin{center}
\begin{tikzpicture}[scale=0.85]
\begin{axis}[
  axis lines=middle, xlabel={$x$}, ylabel={$y$},
  xmin=-4.5, xmax=4.5, ymin=-5.5, ymax=5.5,
  xtick={-3,-2,-1,1,2,3}, ytick={-4,-2,2,4},
  width=0.8\linewidth, height=7cm, clip=true]
\addplot[dashed,popmuted,domain=-4.5:4.5]{x};
\draw[dashed,popmuted] (axis cs:0,-5.5) -- (axis cs:0,5.5);
\addplot[popcurve,domain=0.22:4.5,samples=150]{x + 1/x};
\addplot[popcurve,domain=-4.5:-0.22,samples=150]{x + 1/x};
\addplot[mark=*,poporange,only marks] coordinates {(1,2) (-1,-2)};
\node[popmuted,anchor=west] at (axis cs:2.6,2.2) {\small $y=x$};
\node[popink,anchor=south west] at (axis cs:1.05,2.05) {\small $(1\,;\,2)$};
\end{axis}
\end{tikzpicture}
\end{center}
\legende{Courbe de $f(x) = x + \dfrac{1}{x}$ : asymptotes $x = 0$ et $y = x$ en pointillés, symétrie centrale de centre $O$.}
\end{popfigure}

\begin{savaistu}[Une hyperbole déguisée]
La courbe de $f(x) = x + \dfrac{1}{x}$ est une \cle{hyperbole}, comme celle de $x \mapsto \dfrac{1}{x}$, mais « tournée » : ses asymptotes ne sont plus les axes du repère, mais la droite $x = 0$ et la droite $y = x$. Les asymptotes sont les \emph{directions} vers lesquelles la courbe fuit à l'infini : les géomètres grecs (Ménæchme, Apollonius, III\textsuperscript{ème} siècle av. J.-C.) étudiaient déjà ces courbes obtenues en coupant un cône par un plan. Aujourd'hui, les hyperboles décrivent des phénomènes concrets : la loi d'Ohm de la puissance, les courbes offre-demande en économie, ou encore la trajectoire de certaines comètes qui frôlent le Soleil puis repartent pour ne jamais revenir.
\end{savaistu}

\courssub{Exemple fil rouge : étude et tracé de $f(x) = \dfrac{x^2 + x + 1}{x}$}

Traitons maintenant un exemple complet, du type de celui qui tombe au Baccalauréat, en suivant scrupuleusement la méthode et en construisant la courbe \emph{progressivement}, comme tu dois le faire sur ta copie.

\textbf{Étape 1 — Domaine.} Le dénominateur s'annule en $0$ : $D_f = \mathbb{R}^*$. Astuce utile pour la suite : en divisant chaque terme, on obtient la forme réduite
\[ f(x) = x + 1 + \frac{1}{x}. \]
C'est exactement la fonction précédente translatée de $1$ vers le haut ! Ni paire ni impaire, mais on devine déjà que le point $\Omega(0\,;\,1)$ sera centre de symétrie : en effet $f(x) + f(-x) = \left(x + 1 + \tfrac{1}{x}\right) + \left(-x + 1 - \tfrac{1}{x}\right) = 2 = 2 \times 1$.

\textbf{Étape 2 — Limites et asymptotes.}
\[ \lim_{x \to 0^+} f(x) = 1 + (+\infty) = +\infty, \qquad \lim_{x \to 0^-} f(x) = 1 + (-\infty) = -\infty. \]
La droite $x = 0$ (l'axe des ordonnées) est \cle{asymptote verticale}. Ensuite :
\[ \lim_{x \to \pm\infty} \left[f(x) - (x+1)\right] = \lim_{x \to \pm\infty} \frac{1}{x} = 0, \]
donc la droite d'équation $y = x + 1$ est \cle{asymptote oblique} en $-\infty$ et en $+\infty$. Mieux : le signe de $f(x) - (x+1) = \dfrac{1}{x}$ donne la \cle{position relative} : la courbe est \emph{au-dessus} de son asymptote pour $x > 0$ et \emph{en dessous} pour $x < 0$.

\textbf{Étape 3 — Dérivée.}
\[ f'(x) = 1 - \frac{1}{x^2} = \frac{x^2 - 1}{x^2} = \frac{(x-1)(x+1)}{x^2}. \]
Le dénominateur est strictement positif sur $\mathbb{R}^*$ ; le signe de $f'$ est celui de $(x-1)(x+1)$ : positif sur $]-\infty\,;\,-1[$ et $]1\,;\,+\infty[$, négatif sur $]-1\,;\,0[$ et $]0\,;\,1[$.

\textbf{Étape 4 — Tableau de variation enrichi.} On y fait figurer les limites aux bornes, la valeur interdite $0$ (double barre) et les extremums $f(-1) = -1$ et $f(1) = 3$.

\begin{center}
\begin{tabular}{|c|ccccccccc|}
\hline
\rowcolor{popblueL}
$x$ & $-\infty$ & & $-1$ & & $0$ & & $1$ & & $+\infty$ \\
\hline
Signe de $f'(x)$ & & $+$ & $0$ & $-$ & $||$ & $-$ & $0$ & $+$ & \\
\hline
 & & & $-1$ & & $||$ & & & & $+\infty$ \\
Variations de $f$ & & $\nearrow$ & & $\searrow$ & $||$ & $\searrow$ & & $\nearrow$ & \\
 & $-\infty$ & & & & $||$ & & $3$ & & \\
 & & & & $-\infty$ & $||$ & $+\infty$ & & & \\
\hline
\end{tabular}
\end{center}

(Lecture : sur $]-\infty\,;\,-1]$, $f$ croît de $-\infty$ à $-1$ ; sur $[-1\,;\,0[$, elle décroît de $-1$ vers $-\infty$ ; sur $]0\,;\,1]$, elle décroît de $+\infty$ à $3$ ; sur $[1\,;\,+\infty[$, elle croît de $3$ vers $+\infty$.)

\textbf{Étape 5 — Éléments graphiques.} On place en pointillés l'asymptote verticale $x = 0$ et l'asymptote oblique $y = x+1$, puis les points remarquables $A(-1\,;\,-1)$ (maximum local) et $B(1\,;\,3)$ (minimum local), avec leurs tangentes horizontales.

\textbf{Étape 6 — Tracé progressif.} Voici la construction en trois temps, telle qu'on la réalise sur une copie.

\begin{popfigure}
\begin{center}
\begin{tikzpicture}[scale=0.62]
\begin{axis}[
  axis lines=middle, xlabel={$x$}, ylabel={$y$},
  xmin=-5, xmax=5, ymin=-6, ymax=7,
  width=0.48\linewidth, height=6.5cm, clip=true,
  title={\small Étape A : les asymptotes}]
\addplot[dashed,popmuted,domain=-5:4.8]{x+1};
\draw[dashed,popmuted] (axis cs:0,-6) -- (axis cs:0,7);
\node[popmuted,anchor=west,font=\small] at (axis cs:2.2,2.6) {$y=x+1$};
\end{axis}
\end{tikzpicture}\hfill
\begin{tikzpicture}[scale=0.62]
\begin{axis}[
  axis lines=middle, xlabel={$x$}, ylabel={$y$},
  xmin=-5, xmax=5, ymin=-6, ymax=7,
  width=0.48\linewidth, height=6.5cm, clip=true,
  title={\small Étape B : points et variations}]
\addplot[dashed,popmuted,domain=-5:4.8]{x+1};
\draw[dashed,popmuted] (axis cs:0,-6) -- (axis cs:0,7);
\addplot[mark=*,poporange,only marks] coordinates {(-1,-1) (1,3)};
\draw[popgreen,thick] (axis cs:-1.7,-1) -- (axis cs:-0.3,-1);
\draw[popgreen,thick] (axis cs:0.3,3) -- (axis cs:1.7,3);
\draw[->,popblue,thick] (axis cs:-3.5,-4) -- (axis cs:-1.3,-1.3);
\draw[->,popblue,thick] (axis cs:-0.8,-1.4) -- (axis cs:-0.25,-4.5);
\draw[->,popblue,thick] (axis cs:0.25,5.5) -- (axis cs:0.8,3.4);
\draw[->,popblue,thick] (axis cs:1.3,3.3) -- (axis cs:3.5,5.6);
\end{axis}
\end{tikzpicture}
\end{center}
\legende{Étapes A et B : on place d'abord les asymptotes en pointillés, puis les extremums $A(-1\,;\,-1)$ et $B(1\,;\,3)$ avec leurs tangentes horizontales et le sens de variation.}
\end{popfigure}

\begin{popfigure}
\begin{center}
\begin{tikzpicture}[scale=0.85]
\begin{axis}[
  axis lines=middle, xlabel={$x$}, ylabel={$y$},
  xmin=-5, xmax=5, ymin=-6, ymax=7,
  xtick={-4,-2,2,4}, ytick={-4,-2,2,4,6},
  width=0.85\linewidth, height=8cm, clip=true,
  title={\small Étape C : la courbe complète}]
\addplot[dashed,popmuted,domain=-5:4.8]{x+1};
\draw[dashed,popmuted] (axis cs:0,-6) -- (axis cs:0,7);
\addplot[popcurve,domain=0.16:5,samples=150]{x + 1 + 1/x};
\addplot[popcurve,domain=-5:-0.16,samples=150]{x + 1 + 1/x};
\addplot[mark=*,poporange,only marks] coordinates {(-1,-1) (1,3)};
\node[popink,anchor=south west,font=\small] at (axis cs:1.05,3.05) {$B(1\,;\,3)$};
\node[popink,anchor=north east,font=\small] at (axis cs:-1.05,-1.05) {$A(-1\,;\,-1)$};
\node[popmuted,anchor=west,font=\small] at (axis cs:2.4,2.7) {$y=x+1$};
\end{axis}
\end{tikzpicture}
\end{center}
\legende{Courbe complète de $f(x) = \dfrac{x^2+x+1}{x}$ : chaque branche suit ses asymptotes et passe par les extremums en respectant le tableau de variation.}
\end{popfigure}

\courssub{Vérification finale : le contrôle de cohérence}

Avant de rendre ta copie, prends trente secondes pour vérifier que ta courbe « raconte la même histoire » que tes calculs. C'est ce réflexe qui distingue une copie excellente d'une copie correcte.

\begin{methode}[La check-list de cohérence]
\begin{enumerate}
\item \textbf{Domaine} : la courbe n'existe-t-elle que là où $f$ est définie ? Aucun point sur la droite $x = 0$ ici.
\item \textbf{Asymptote verticale} : chaque branche fuit-elle vers le bon infini ? Ici $+\infty$ à droite de $0$, $-\infty$ à gauche — conforme aux limites calculées.
\item \textbf{Asymptote oblique} : la courbe se rapproche-t-elle de $y = x+1$ aux deux extrémités, du bon côté (au-dessus à droite, en dessous à gauche, signe de $\tfrac{1}{x}$) ?
\item \textbf{Variations} : en parcourant la courbe de gauche à droite, lit-on bien croissant, décroissant, décroissant, croissant, avec les sommets aux bonnes ordonnées $-1$ et $3$ ?
\item \textbf{Symétrie} : la courbe est-elle invariante par la symétrie de centre $\Omega(0\,;\,1)$ ? Les points $A(-1\,;\,-1)$ et $B(1\,;\,3)$ sont-ils symétriques par rapport à $\Omega$ ? Oui : le milieu de $[AB]$ est $(0\,;\,1)$.
\end{enumerate}
\end{methode}

\begin{exoresolu}[Bac blanc : étude complète et tracé]
Soit $g$ la fonction définie par $g(x) = \dfrac{x^2 - 2x + 2}{x - 1}$.
\begin{enumerate}
\item Déterminer le domaine de $g$ et montrer que $g(x) = x - 1 + \dfrac{1}{x-1}$.
\item Déterminer les asymptotes de la courbe $\mathcal{C}_g$.
\item Montrer que le point $\Omega(1\,;\,0)$ est centre de symétrie de $\mathcal{C}_g$.
\item Dresser le tableau de variation de $g$ et décrire le tracé de $\mathcal{C}_g$.
\end{enumerate}
\tcblower
\textbf{1.} Le dénominateur s'annule en $x = 1$ : $D_g = \mathbb{R} \setminus \{1\}$. Effectuons la division : $x^2 - 2x + 2 = (x-1)(x-1) + 1$, car $(x-1)^2 = x^2 - 2x + 1$. Donc
\[ g(x) = \frac{(x-1)^2 + 1}{x-1} = x - 1 + \frac{1}{x-1}. \]
\textbf{2.} Limite en $1$ : le numérateur tend vers $1 > 0$ et $x - 1 \to 0^\pm$, donc $\lim_{x \to 1^+} g(x) = +\infty$ et $\lim_{x \to 1^-} g(x) = -\infty$ : la droite $x = 1$ est asymptote verticale. Ensuite $g(x) - (x-1) = \dfrac{1}{x-1} \xrightarrow[x \to \pm\infty]{} 0$ : la droite $y = x - 1$ est asymptote oblique.
\textbf{3.} Posons $x = 1 + t$. Alors
\[ g(1+t) + g(1-t) = \left(t + \frac{1}{t}\right) + \left(-t - \frac{1}{t}\right) = 0 = 2 \times 0, \]
donc $\Omega(1\,;\,0)$ est centre de symétrie de $\mathcal{C}_g$. (On reconnaît d'ailleurs la courbe de $x + \tfrac{1}{x}$ translatée du vecteur $\vec{\imath}$.)
\textbf{4.} $g'(x) = 1 - \dfrac{1}{(x-1)^2} = \dfrac{(x-1)^2 - 1}{(x-1)^2} = \dfrac{x(x-2)}{(x-1)^2}$. Le signe est celui de $x(x-2)$ : $g$ est croissante sur $]-\infty\,;\,0]$ et $[2\,;\,+\infty[$, décroissante sur $[0\,;\,1[$ et $]1\,;\,2]$. Extremums : $g(0) = -2$ (maximum local) et $g(2) = 2$ (minimum local), symétriques par rapport à $\Omega$. Pour le tracé : pointillés $x = 1$ et $y = x-1$, points $(0\,;\,-2)$ et $(2\,;\,2)$ à tangentes horizontales, branches qui fuient le long des asymptotes, et contrôle final par la symétrie de centre $\Omega$.
\end{exoresolu}

\begin{exercice}[À toi de jouer]
Soit $h(x) = \dfrac{x^2 + 3}{x}$. Étudier $h$ complètement (domaine, parité, limites, asymptotes, dérivée, tableau de variation enrichi) puis tracer sa courbe dans un repère orthonormé, en suivant les trois étapes de construction progressive : asymptotes, points remarquables, courbe. Vérifier enfin la cohérence du dessin avec la check-list.
\end{exercice}

\begin{corrige}[Exercice 1]
$D_h = \mathbb{R}^*$ ; $h(x) = x + \dfrac{3}{x}$ et $h(-x) = -h(x)$ : $h$ est impaire, on étudie sur $]0\,;\,+\infty[$. Limites : $\lim_{0^+} h = +\infty$ (asymptote verticale $x = 0$) ; $h(x) - x = \dfrac{3}{x} \to 0$ en $\pm\infty$ (asymptote oblique $y = x$, courbe au-dessus pour $x > 0$). Dérivée : $h'(x) = 1 - \dfrac{3}{x^2} = \dfrac{x^2 - 3}{x^2}$, négative sur $]0\,;\,\sqrt{3}[$, positive sur $]\sqrt{3}\,;\,+\infty[$. Minimum en $\sqrt{3}$ : $h(\sqrt{3}) = \sqrt{3} + \dfrac{3}{\sqrt{3}} = 2\sqrt{3} \approx 3,46$. Sur $]0\,;\,+\infty[$, $h$ décroît de $+\infty$ à $2\sqrt{3}$ puis croît vers $+\infty$ ; on complète par symétrie de centre $O$ (maximum local $(-\sqrt{3}\,;\,-2\sqrt{3})$). Tracé : pointillés $x = 0$ et $y = x$, points $(\sqrt{3}\,;\,2\sqrt{3})$ et $(-\sqrt{3}\,;\,-2\sqrt{3})$ à tangentes horizontales, branches collées aux asymptotes. Cohérence : les deux extremums sont bien symétriques par rapport à $O$, et la courbe reste au-dessus de $y = x$ à droite, en dessous à gauche.
\end{corrige}

\begin{aretenir}[L'essentiel de la synthèse]
\begin{itemize}
\item L'étude complète d'une fonction suit \textbf{6 étapes} : domaine, limites et asymptotes, dérivée, tableau de variation enrichi, éléments graphiques, tracé vérifié.
\item Une \textbf{symétrie} (centre ou axe) permet de n'étudier que la moitié du domaine, puis de compléter la courbe par la transformation géométrique.
\item Les \textbf{asymptotes} sont les rails du tracé : elles guident les branches infinies et doivent toujours figurer en pointillés.
\item Le \textbf{tableau de variation enrichi} (limites, valeurs interdites, extremums) est le plan de construction de la courbe : tout ce qui est dans le tableau doit se lire sur le dessin, et réciproquement.
\end{itemize}
\end{aretenir}

\newpage

\coursec{Activité d'intégration}

\coursec{Courbes représentatives}

\courssub{Objectifs de l'activité}

À l'issue de cette activité d'intégration, tu seras capable de :
\begin{itemize}
\item lire et interpréter une courbe représentative dans un contexte économique concret (coût moyen d'un forage, recettes d'une commune) ;
\item déterminer les équations des asymptotes verticales, horizontales et obliques d'une courbe à l'aide de calculs de limites, et donner un sens concret à ces asymptotes ;
\item dresser un tableau de variation et construire une courbe complète en utilisant ses asymptotes et ses éventuels centres ou axes de symétrie ;
\item démontrer qu'un point est centre de symétrie d'une courbe par la méthode du changement de repère (ou par la relation $f(a-h)+f(a+h)=2b$) ;
\item émettre une conjecture à partir d'une représentation graphique, puis la vérifier par le calcul algébrique.
\end{itemize}

\courssub{Situation}

La commune de Bangangté, dans la région de l'Ouest, a décidé de financer le forage d'un puits afin d'alimenter en eau potable le quartier de Bamendou. L'entreprise de forage retenue présente le devis suivant :
\begin{itemize}
\item des \textbf{frais fixes} (transport du matériel, installation de la foreuse, études géologiques) s'élevant à \num{2000000} FCFA ;
\item un \textbf{coût variable} de \num{1000} FCFA par mètre foré (carburant, tubage, main-d'œuvre).
\end{itemize}

Si la profondeur totale forée est de $x$ mètres (avec $x > 0$), le coût total est donc, en FCFA :
\[ \num{2000000} + \num{1000}\,x. \]
Le \textbf{coût moyen par mètre foré}, noté $C(x)$, est le quotient du coût total par le nombre de mètres forés :
\[ C(x) = \frac{\num{2000000} + \num{1000}\,x}{x} = \num{1000} + \frac{\num{2000000}}{x}, \qquad x \in \mathopen]0\,;\,+\infty\mathclose[. \]

Le comité de gestion de l'eau souhaite comprendre comment ce coût moyen évolue lorsque la profondeur augmente, afin de convaincre les contribuables qu'il est plus rentable de forer profond (pour atteindre une nappe pérenne) plutôt que de s'arrêter à une faible profondeur.

Par ailleurs, la commune revend une partie de l'eau à une société de distribution. Les recettes cumulées, exprimées en \textbf{millions de FCFA}, sont modélisées en fonction de la quantité d'eau vendue $x$ (en milliers de mètres cubes) par :
\[ R(x) = \frac{2x+3}{x+1}, \qquad x \geq 0. \]
Le trésorier communal affirme, après avoir tracé la courbe de $R$ sur son tableur, que celle-ci possède « une belle symétrie autour du point $\Omega(-1\,;\,2)$ » et qu'elle « se rapproche de certaines droites sans jamais les toucher ». Il te demande de vérifier rigoureusement ses affirmations, puis de déterminer graphiquement et par le calcul à partir de quelle profondeur le coût moyen du mètre foré devient inférieur à \num{3000} FCFA.

\courssub{Consignes}

\textbf{Partie A : le coût moyen du forage}
\begin{enumerate}
\item Justifier que $C(x) = \num{1000} + \dfrac{\num{2000000}}{x}$ pour tout $x > 0$.
\item Calculer $\lim\limits_{x \to 0^+} C(x)$. En déduire que la courbe $\mathcal{C}$ de $C$ admet une asymptote verticale dont on donnera l'équation. Interpréter ce résultat en langage économique.
\item Calculer $\lim\limits_{x \to +\infty} C(x)$. En déduire que $\mathcal{C}$ admet une asymptote horizontale dont on donnera l'équation. Pourquoi le trésorier appelle-t-il \num{1000} FCFA le « coût plancher » ?
\item Étudier le sens de variation de $C$ sur $\mathopen]0\,;\,+\infty\mathclose[$ et dresser son tableau de variation complet (avec les limites).
\item Tracer soigneusement $\mathcal{C}$ et ses deux asymptotes dans un repère orthogonal (unités : 1 cm pour 50 m en abscisse, 1 cm pour \num{1000} FCFA en ordonnée, en commençant l'axe des ordonnées à \num{1000}).
\end{enumerate}

\textbf{Partie B : la recette de la commune}
\begin{enumerate}
\setcounter{enumi}{5}
\item Déterminer l'ensemble de définition de $R$, puis calculer les limites de $R$ en $-1$ (à gauche et à droite) et en $+\infty$. En déduire toutes les asymptotes de la courbe $\mathcal{C}_R$.
\item Démontrer que le point $\Omega(-1\,;\,2)$ est centre de symétrie de $\mathcal{C}_R$ (on pourra poser $X = x+1$ et $Y = y - 2$, ou vérifier que $R(-1-h) + R(-1+h) = 4$ pour tout $h \neq 0$).
\item Tracer $\mathcal{C}_R$ et ses asymptotes dans un repère orthonormé, en exploitant la symétrie démontrée.
\end{enumerate}

\textbf{Partie C : conjecture et vérification}
\begin{enumerate}
\setcounter{enumi}{8}
\item À l'aide du graphique de la partie A, conjecturer à partir de quelle profondeur $x_0$ le coût moyen par mètre devient inférieur à \num{3000} FCFA.
\item Vérifier cette conjecture par le calcul, en résolvant l'inéquation $C(x) < \num{3000}$.
\end{enumerate}

\courssub{Ressources à mobiliser}

\begin{methode}[Recherche des asymptotes d'une courbe]
\begin{enumerate}
\item \textbf{Asymptotes verticales} : Chercher les valeurs $x = a$ non appartenantes à l'ensemble de définition (zéros du dénominateur pour une fonction rationnelle) et calculer $\lim\limits_{x \to a^-} f(x)$ et $\lim\limits_{x \to a^+} f(x)$. Si l'une de ces limites est $\pm\infty$, la droite $x = a$ est asymptote verticale.
\item \textbf{Asymptotes horizontales} : Calculer $\lim\limits_{x \to +\infty} f(x)$ et $\lim\limits_{x \to -\infty} f(x)$. Si la limite est un réel $b$, la droite $y = b$ est asymptote horizontale dans la direction correspondante.
\item \textbf{Asymptotes obliques} (si pas d'asymptote horizontale en $+\infty$ ou $-\infty$) : Calculer $a = \lim\limits_{x \to \pm\infty} \frac{f(x)}{x}$. Si $a \in \mathbb{R}^*$, calculer $b = \lim\limits_{x \to \pm\infty} \bigl(f(x) - ax\bigr)$. La droite $y = ax + b$ est l'asymptote oblique.
\end{enumerate}
\end{methode}

\begin{methode}[Preuve d'un centre de symétrie $\Omega(a\,;\,b)$]
Deux méthodes équivalentes au choix :
\begin{enumerate}
\item \textbf{Relation caractéristique} : Vérifier que pour tout $h$ tel que $a-h$ et $a+h$ appartiennent à $D_f$,
      \[ f(a-h) + f(a+h) = 2b. \]
\item \textbf{Changement de repère} : Poser $X = x - a$ et $Y = y - b$ (ou $g(X) = f(a+X) - b$). Montrer que la fonction $g$ est \cle{impaire} (c.-à-d. $g(-X) = -g(X)$ pour tout $X$ dans son ensemble de définition).
\end{enumerate}
\end{methode}

\begin{aretenir}[Boîte à outils de la leçon]
\begin{itemize}
\item \textbf{Axe de symétrie (vertical)} : La droite $x = a$ est axe de symétrie de $\mathcal{C}_f$ si, pour tout $h$ tel que $a-h$ et $a+h$ appartiennent à $D_f$,
      \[ f(a-h) = f(a+h). \]
      Équivalent : la fonction $g : h \mapsto f(a+h)$ est \cle{paire}.
\item \textbf{Centre de symétrie} : Le point $\Omega(a\,;\,b)$ est centre de symétrie de $\mathcal{C}_f$ si, pour tout $h$ admissible,
      \[ f(a-h) + f(a+h) = 2b. \]
      Équivalent : la courbe de $g : X \mapsto f(a+X) - b$ est celle d'une fonction \cle{impaire}.
\item \textbf{Asymptote verticale} : Si $\lim\limits_{x \to a} f(x) = \pm\infty$, alors la droite d'équation $x = a$ est asymptote verticale à $\mathcal{C}_f$.
\item \textbf{Asymptote horizontale} : Si $\lim\limits_{x \to \pm\infty} f(x) = b \in \mathbb{R}$, alors la droite d'équation $y = b$ est asymptote horizontale à $\mathcal{C}_f$ dans la direction correspondante.
\item \textbf{Asymptote oblique} : Si $\lim\limits_{x \to \pm\infty} \frac{f(x)}{x} = a \in \mathbb{R}^*$ et $\lim\limits_{x \to \pm\infty} \bigl(f(x) - ax\bigr) = b \in \mathbb{R}$, alors la droite $y = ax + b$ est asymptote oblique.
\item \textbf{Lecture graphique} : Une conjecture issue d'un graphique doit toujours être confirmée par un calcul (résolution d'équation ou d'inéquation).
\end{itemize}
\end{aretenir}

\courssub{Démarche et corrigé commenté}

\begin{exoresolu}[Le forage de la commune : corrigé complet]
\textbf{Rappel des questions (Parties A, B, C).}
\tcblower

\textbf{Partie A.}
\begin{enumerate}
\item Pour tout $x > 0$ :
      \[
      \frac{\num{2000000} + \num{1000}x}{x}
      = \frac{\num{2000000}}{x} + \frac{\num{1000}x}{x}
      = \num{1000} + \frac{\num{2000000}}{x}.
      \]
\item Quand $x \to 0^+$ : le numérateur $\num{2000000} + \num{1000}x \to \num{2000000} > 0$ et le dénominateur $x \to 0^+$.
      Donc $\lim\limits_{x \to 0^+} C(x) = +\infty$.
      La droite d'équation $x = 0$ (l'axe des ordonnées) est une \textbf{asymptote verticale} à la courbe $\mathcal{C}$.
      \emph{Interprétation économique} : si l'on fore très peu de mètres, les frais fixes de \num{2000000} FCFA sont répartis sur un tout petit nombre de mètres ; le coût moyen par mètre « explose » (tend vers $+\infty$).
\item $\lim\limits_{x \to +\infty} \dfrac{\num{2000000}}{x} = 0$, donc $\lim\limits_{x \to +\infty} C(x) = \num{1000}$.
      La droite d'équation $y = \num{1000}$ est une \textbf{asymptote horizontale} à $\mathcal{C}$ en $+\infty$.
      \emph{Interprétation} : plus on fore profond, plus les frais fixes se « diluent » dans le total. Le coût moyen se rapproche de \num{1000} FCFA/m sans jamais descendre en dessous : c'est le « coût plancher », égal au coût variable unitaire.
\item $C$ est dérivable sur $\mathopen]0\,;\,+\infty\mathclose[$ et
      \[
      C'(x) = -\frac{\num{2000000}}{x^2} < 0 \quad \text{pour tout } x > 0.
      \]
      $C$ est donc \textbf{strictement décroissante} sur $\mathopen]0\,;\,+\infty\mathclose[$.
      \begin{center}
      \begin{tabular}{|c|ccccc|}
      \hline
      $x$ & $0$ & & & & $+\infty$ \\
      \hline
      $C'(x)$ & & $-$ & $-$ & $-$ & \\
      \hline
      $C(x)$ & $+\infty$ & $\searrow$ & & $\searrow$ & $\num{1000}$ \\
      \hline
      \end{tabular}
      \end{center}
\item \textbf{Tracé de la courbe $\mathcal{C}$ et de ses asymptotes.}
\end{enumerate}

\begin{popfigure}
\begin{tikzpicture}
\begin{axis}[
    width=\linewidth,
    height=0.55\linewidth,
    axis lines=middle,
    xlabel={$x$ (m)},
    ylabel={$C(x)$ (FCFA)},
    xmin=0, xmax=5000,
    ymin=1000, ymax=22000,
    xtick={0,1000,2000,3000,4000,5000},
    ytick={1000,5000,10000,15000,20000},
    yticklabels={1000,5000,10000,15000,20000},
    x tick label style={/pgf/number format/1000 sep={\,}, font=\small},
    y tick label style={/pgf/number format/1000 sep={\,}, font=\small},
    domain=50:5000,
    samples=200,
    clip=false,
    axis line style={-Stealth, popdark},
    tick style={thin, popmuted},
    grid=major,
    grid style={popline, opacity=0.4},
    enlargelimits={abs=0.05},
]
% Asymptote verticale x=0 (axe des ordonnées) - déjà là par axis lines=middle
% On la met en évidence
\addplot[poporange, dashed, thick, domain=1000:22000] ({0}, x) 
    node[above left, pos=0.95, font=\small, poporange] {$x=0$};
% Asymptote horizontale y=1000
\addplot[poporange, dashed, thick, domain=0:5000] (x, {1000}) 
    node[above right, pos=0.95, font=\small, poporange] {$y=1000$};
% Courbe C(x) = 1000 + 2000000/x
\addplot[popcurve, thick, popblue, domain=50:5000] {1000 + 2000000/x};
% Points remarquables
\addplot[only marks, mark=*, popgreen, mark size=2.5] coordinates {
    (100, 21000)
    (500, 5000)
    (1000, 3000)
    (2000, 2000)
    (4000, 1500)
};
% Droite y=3000 pour la conjecture
\addplot[poppink, dashed, thin, domain=0:5000] (x, {3000}) 
    node[right, pos=0.9, font=\small, poppink] {$y=3000$};
\end{axis}
\end{tikzpicture}
\legende{Courbe $\mathcal{C}$ du coût moyen $C(x) = 1000 + \frac{2\,000\,000}{x}$, ses asymptotes $x=0$ (orange) et $y=1000$ (orange), points remarquables (verts) et droite $y=3000$ (rose) pour la conjecture.}
\end{popfigure}

\textbf{Partie B.}
\begin{enumerate}
\setcounter{enumi}{5}
\item $R(x) = \dfrac{2x+3}{x+1}$. Le dénominateur s'annule pour $x = -1$.
      Ensemble de définition : $D_R = \mathbb{R} \setminus \{-1\}$.
      \begin{itemize}
      \item \textbf{En $x = -1$} : numérateur $2(-1)+3 = 1 > 0$.
            \begin{align*}
            \lim_{x \to -1^-} (x+1) = 0^- &\implies \lim_{x \to -1^-} R(x) = \frac{1}{0^-} = -\infty, \\
            \lim_{x \to -1^+} (x+1) = 0^+ &\implies \lim_{x \to -1^+} R(x) = \frac{1}{0^+} = +\infty.
            \end{align*}
            La droite $x = -1$ est \textbf{asymptote verticale}.
      \item \textbf{En $+\infty$} :
            \[
            R(x) = \frac{x(2 + \frac{3}{x})}{x(1 + \frac{1}{x})} = \frac{2 + \frac{3}{x}}{1 + \frac{1}{x}} \xrightarrow[x \to +\infty]{} 2.
            \]
            La droite $y = 2$ est \textbf{asymptote horizontale} en $+\infty$.
            (De même en $-\infty$, limite $2$, même asymptote).
      \end{itemize}
\item \textbf{Démonstration que $\Omega(-1\,;\,2)$ est centre de symétrie.}
      \textit{Méthode 1 (relation caractéristique) :} Pour tout $h \neq 0$,
      \begin{align*}
      R(-1-h) + R(-1+h)
      &= \frac{2(-1-h)+3}{-1-h+1} + \frac{2(-1+h)+3}{-1+h+1} \\
      &= \frac{1 - 2h}{-h} + \frac{1 + 2h}{h} \\
      &= -\frac{1-2h}{h} + \frac{1+2h}{h} = \frac{-1+2h+1+2h}{h} = \frac{4h}{h} = 4 = 2 \times 2.
      \end{align*}
      Donc $R(-1-h) + R(-1+h) = 2 \times 2$ pour tout $h \neq 0$ : $\Omega(-1\,;\,2)$ est bien centre de symétrie.
      
      \textit{Méthode 2 (changement de repère) :} On écrit $R(x) = 2 + \dfrac{1}{x+1}$.
      Posons $X = x + 1$ et $Y = y - 2$. Alors $y = R(x) \iff Y + 2 = 2 + \dfrac{1}{X} \iff Y = \dfrac{1}{X}$.
      La courbe de $Y = \frac{1}{X}$ est l'hyperbole équilatère, courbe d'une fonction \cle{impaire}, qui admet l'origine $(0,0)$ comme centre de symétrie.
      Par translation de vecteur $\overrightarrow{(-1,\,2)}$, l'origine devient $\Omega(-1\,;\,2)$, centre de symétrie de $\mathcal{C}_R$.
\item \textbf{Tracé de $\mathcal{C}_R$ et de ses asymptotes.}
      On trace la branche pour $x > -1$ (domaine du modèle $x \ge 0$ inclus) : décroissante, de $+\infty$ (en $-1^+$) vers $2$ (en $+\infty$).
      La branche pour $x < -1$ s'en déduit par symétrie de centre $\Omega$ : décroissante, de $2$ (en $-\infty$) vers $-\infty$ (en $-1^-$).
\end{enumerate}

\begin{popfigure}
\begin{tikzpicture}
\begin{axis}[
    width=\linewidth,
    height=0.6\linewidth,
    axis lines=middle,
    xlabel={$x$ (milliers de m$^3$)},
    ylabel={$R(x)$ (millions FCFA)},
    xmin=-5, xmax=5,
    ymin=-5, ymax=5,
    xtick={-5,-4,-3,-2,-1,0,1,2,3,4,5},
    ytick={-5,-4,-3,-2,-1,0,1,2,3,4,5},
    tick label style={font=\small},
    domain=-5:5,
    samples=200,
    restrict y to domain=-10:10,
    clip=false,
    axis line style={-Stealth, popdark},
    tick style={thin, popmuted},
    grid=major,
    grid style={popline, opacity=0.4},
    unit vector ratio=1 1, % Repère orthonormé
]
% Asymptotes
\addplot[poporange, dashed, thick, domain=-5:5] ({-1}, x) 
    node[above, pos=0.9, font=\small, poporange, rotate=90, anchor=south] {$x=-1$};
\addplot[poporange, dashed, thick, domain=-5:5] (x, {2}) 
    node[above right, pos=0.95, font=\small, poporange] {$y=2$};
% Branche droite (x > -1)
\addplot[popcurve, thick, popblue, domain=-0.9:5] {2 + 1/(x+1)};
% Branche gauche (x < -1)
\addplot[popcurve, thick, popblue, domain=-5:-1.1] {2 + 1/(x+1)};
% Centre de symétrie
\addplot[only marks, mark=*, poppurple, mark size=3] coordinates {(-1, 2)};
\node[poppurple, font=\small, anchor=south east] at (-1, 2) {$\Omega(-1;2)$};
% Flèches de symétrie (optionnel, pour illustration)
\draw[poppurple, <->, thick] (-2, 1) -- (-1, 2) -- (0, 3);
\draw[poppurple, <->, thick] (-3, 0.5) -- (-1, 2) -- (1, 3.5);
\end{axis}
\end{tikzpicture}
\legende{Courbe $\mathcal{C}_R$ de $R(x) = \frac{2x+3}{x+1}$, ses asymptotes $x=-1$ et $y=2$ (orange), son centre de symétrie $\Omega(-1;2)$ (violet) et illustration de la symétrie centrale.}
\end{popfigure}

\textbf{Partie C.}
\begin{enumerate}
\setcounter{enumi}{8}
\item Sur le graphique de la Partie A (Figure 1), la courbe $\mathcal{C}$ passe sous la droite $y = \num{3000}$ aux alentours de $x = \num{1000}$ m.
      \emph{Conjecture} : le coût moyen devient inférieur à \num{3000} FCFA à partir d'environ \num{1000} mètres.
\item \textbf{Vérification par le calcul} : pour $x > 0$,
      \begin{align*}
      C(x) < \num{3000}
      &\iff \num{1000} + \frac{\num{2000000}}{x} < \num{3000} \\
      &\iff \frac{\num{2000000}}{x} < \num{2000} \\
      &\iff x > \frac{\num{2000000}}{\num{2000}} = \num{1000} \quad (\text{car } x > 0).
      \end{align*}
      La conjecture est confirmée : le coût moyen devient \textbf{strictement inférieur à \num{3000} FCFA} dès que la profondeur dépasse \textbf{\num{1000} mètres}.
\end{enumerate}
\end{exoresolu}

\courssub{Prolongement : Symétrie axiale et asymptotes obliques}

Le programme exige aussi la maîtrise de la \textbf{symétrie axiale} et des \textbf{asymptotes obliques}. Voici deux exemples types pour compléter ta boîte à outils.

\begin{exemplebox}[Symétrie axiale : $f(x) = \frac{x^2}{x^2+1}$]
Soit $f(x) = \frac{x^2}{x^2+1}$, $D_f = \mathbb{R}$.
\begin{enumerate}
\item $f(-x) = \frac{(-x)^2}{(-x)^2+1} = \frac{x^2}{x^2+1} = f(x)$ : $f$ est \cle{paire}.
\item La courbe $\mathcal{C}_f$ admet donc l'axe des ordonnées ($x=0$) comme \textbf{axe de symétrie}.
\item Asymptote horizontale : $\lim\limits_{x \to \pm\infty} f(x) = 1$, donc $y=1$ est asymptote horizontale (la courbe est toujours en dessous car $x^2 < x^2+1$).
\item Pas d'asymptote verticale (dénominateur jamais nul).
\end{enumerate}
\textbf{Tableau de variation} (à faire : dériver, signe de $f'(x) = \frac{2x}{(x^2+1)^2}$) : décroissante sur $]-\infty,0]$, croissante sur $[0,+\infty[$, minimum $f(0)=0$.
\end{exemplebox}

\begin{exemplebox}[Asymptote oblique : $g(x) = \frac{x^2+1}{x}$]
Soit $g(x) = x + \frac{1}{x}$, $D_g = \mathbb{R}^*$.
\begin{enumerate}
\item \textbf{Asymptote verticale} : $x=0$ car $\lim\limits_{x \to 0^\pm} g(x) = \pm\infty$.
\item \textbf{Asymptote oblique} : Pas d'asymptote horizontale ($\lim\limits_{x \to \pm\infty} g(x) = \pm\infty$).
      \begin{align*}
      a &= \lim_{x \to \pm\infty} \frac{g(x)}{x} = \lim_{x \to \pm\infty} \frac{x + 1/x}{x} = 1, \\
      b &= \lim_{x \to \pm\infty} \bigl(g(x) - x\bigr) = \lim_{x \to \pm\infty} \frac{1}{x} = 0.
      \end{align*}
      La droite $y = x$ est \textbf{asymptote oblique} (commune aux deux branches).
\item \textbf{Centre de symétrie} : $g(-x) = -g(x)$, $g$ est impaire $\implies$ l'origine $O(0,0)$ est centre de symétrie.
\end{enumerate}
\end{exemplebox}

\begin{exercice}[Entraînement : Asymptote oblique et symétrie]
Soit $h(x) = \frac{x^2 - 2x + 3}{x - 1}$, $D_h = \mathbb{R} \setminus \{1\}$.
\begin{enumerate}
\item Effectuer la division euclidienne du numérateur par le dénominateur pour écrire $h(x)$ sous la forme $ax + b + \frac{c}{x-1}$.
\item En déduire les asymptotes de $\mathcal{C}_h$ (verticale et oblique).
\item Montrer que $\mathcal{C}_h$ admet un centre de symétrie et donner ses coordonnées.
\item Dresser le tableau de variation de $h$ et tracer $\mathcal{C}_h$ avec ses asymptotes.
\end{enumerate}
\end{exercice}

\begin{corrige}[Exercice : Asymptote oblique et symétrie]
\begin{enumerate}
\item Division : $(x^2 - 2x + 3) \div (x-1) = x - 1$ reste $2$.
      Donc $h(x) = x - 1 + \frac{2}{x-1}$.
\item Asymptote verticale : $x = 1$ (pôle).
      Asymptote oblique : $y = x - 1$ (car $\frac{2}{x-1} \to 0$ en $\pm\infty$).
\item Posons $X = x-1$, $Y = y + 1$. Alors $Y = \frac{2}{X}$, courbe d'une fonction impaire.
      Centre de symétrie : $\Omega(1\,;\,-1)$.
\item $h'(x) = 1 - \frac{2}{(x-1)^2} = \frac{(x-1)^2 - 2}{(x-1)^2} = \frac{(x-1-\sqrt{2})(x-1+\sqrt{2})}{(x-1)^2}$.
      Variation : $\nearrow$ sur $]-\infty, 1-\sqrt{2}]$, $\searrow$ sur $[1-\sqrt{2}, 1[$, $\searrow$ sur $]1, 1+\sqrt{2}]$, $\nearrow$ sur $[1+\sqrt{2}, +\infty[$.
      Max local en $1-\sqrt{2}$, min local en $1+\sqrt{2}$.
\end{enumerate}
\end{corrige}

\courssub{Bilan de l'activité}

Cette activité t'a permis de mobiliser l'ensemble des savoir-faire de la leçon dans une situation authentique de gestion communale :
\begin{itemize}
\item les \textbf{limites} ont donné naissance aux \textbf{asymptotes} $x = 0$ et $y = \num{1000}$ (Partie A), $x = -1$ et $y = 2$ (Partie B), et chacune a reçu une interprétation économique concrète (explosion du coût moyen pour un forage trop court, coût plancher de \num{1000} FCFA/m, seuil de rentabilité pour la recette) ;
\item le \textbf{tableau de variation}, complété par les asymptotes, a suffi à construire une courbe fidèle sans calculer des dizaines de points ;
\item la démonstration du \textbf{centre de symétrie} $\Omega(-1\,;\,2)$ a montré comment une propriété géométrique se prouve algébriquement (relation caractéristique ou changement de repère), et comment elle divise par deux le travail de tracé ;
\item enfin, la démarche \emph{conjecturer graphiquement puis vérifier par le calcul} illustre l'allers-retours permanent entre le graphique et l'algèbre, compétence centrale de l'épreuve de mathématiques au Baccalauréat A.
\end{itemize}

\begin{attention}[Pièges à éviter dans cette activité]
\begin{itemize}
\item Ne pas oublier le \textbf{signe du numérateur} lors d'une limite du type $\dfrac{k}{x - a}$ : c'est lui qui détermine si la limite est $+\infty$ ou $-\infty$ de chaque côté de $a$.
\item Une asymptote horizontale $y = b$ ne signifie pas que la courbe ne peut jamais couper la droite : elle décrit un comportement \emph{à l'infini}. Ici, $C(x) > \num{1000}$ pour tout $x$, mais ce n'est pas une règle générale (ex. $f(x) = \frac{\sin x}{x}$ coupe $y=0$ infiniment souvent).
\item Pour prouver un centre de symétrie, il faut vérifier la relation $f(a-h)+f(a+h)=2b$ pour \textbf{tout} $h$ admissible, pas seulement sur un exemple numérique.
\item Une conjecture graphique reste une conjecture : la résolution de l'inéquation (ou de l'équation) est la seule preuve acceptée au Baccalauréat.
\item \textbf{Asymptote oblique} : ne pas confondre « division euclidienne » (pour les fonctions rationnelles) et « calcul des limites $a = \lim f(x)/x$, $b = \lim f(x)-ax$ » (méthode générale, toujours valable). Les deux donnent le même résultat.
\end{itemize}
\end{attention}

\begin{savaistu}[L'eau, un enjeu mathématique]
Selon les programmes d'hydraulique rurale menés au Cameroun (avec l'appui de la Banque Mondiale et de l'AFD), le coût d'un forage varie fortement selon la profondeur de la nappe phréatique : de quelques dizaines de mètres dans les plaines du Nord (bassin du Lac Tchad) à plus de 100--150 mètres sur les plateaux de l'Ouest et du Sud. Les modèles de coût moyen comme celui de cette activité ($C(x) = \text{fixe}/x + \text{variable}$) aident les communes et les ONG à comparer les devis, à planifier les budgets d'adduction d'eau et à négocier les tarifs de revente. Les mathématiques de Terminale sont déjà un outil de décision publique et de développement local !
\end{savaistu}

\newpage

\coursec{Méthodes et exercices résolus}

\coursec{Courbes représentatives}

\coursec{Lecture graphique et conjectures}

\begin{definition}[Conjecture graphique]
Une \cle{conjecture} est une affirmation formulée à partir de l'observation d'une représentation graphique (courbe, nuage de points, histogramme). Elle exprime une propriété qui \emph{semble} vraie (signe, symétrie, asymptote, extremum, monotonie) mais qui nécessite une \textbf{vérification algébrique rigoureuse} pour devenir un résultat certain.
\end{definition}

\begin{methode}[Lire une courbe et vérifier une conjecture]
Face à une courbe représentative $\mathcal{C}_f$ donnée (au Bac, souvent celle d'une fonction auxiliaire ou de la fonction étudiée) :
\begin{enumerate}
\item \textbf{Observer} : positions relatives (au-dessus/en dessous d'une droite), symétries apparentes, comportement aux bornes (branches qui « collent » à une droite), extrema, intersections avec les axes.
\item \textbf{Formuler la conjecture} en langage mathématique précis : « la courbe semble admettre la droite $x=a$ pour axe de symétrie », « $f$ semble positive sur $[0\,;\,+\infty[$ », « la droite $y=2$ semble être une asymptote horizontale en $+\infty$ ».
\item \textbf{Vérifier par le calcul} : une conjecture graphique ne vaut jamais démonstration.
      \begin{itemize}
      \item Pour un \textbf{signe} : résoudre $f(x)\geq 0$ ou dresser un tableau de signes.
      \item Pour une \textbf{symétrie} : appliquer les méthodes 2 ou 3 (axe ou centre).
      \item Pour une \textbf{asymptote} : calculer la limite adéquate (finie à l'infini ou infinie en un point fini).
      \end{itemize}
\item \textbf{Conclure} explicitement : « la conjecture est validée » ou « elle est infirmée ».
\end{enumerate}
\begin{aretenir}[Réflexe Bac]
Toute lecture graphique rapporte des points si elle est \emph{formulée clairement}, mais la \cle{vérification algébrique} est exigée dès que la question dit « démontrer » ou « justifier ». Ne jamais écrire « on voit sur le graphique que... » sans calcul à l'appui.
\end{aretenir}
\end{methode}

\begin{exoresolu}[Conjecture sur le signe d'une fonction]
On donne ci-dessous la courbe $\mathcal{C}_g$ de la fonction $g$ définie sur $\mathbb{R}$ par $g(x)=x^2-2x+2$.
\begin{popfigure}
\begin{center}
\begin{tikzpicture}[scale=0.85]
\draw[->,popmuted] (-1.5,0) -- (3.5,0) node[below]{$x$};
\draw[->,popmuted] (0,-0.5) -- (0,5) node[left]{$y$};
\draw[popblue,thick,domain=-1:3,samples=100] plot (\x,{\x*\x-2*\x+2});
\draw (1,1) node[below right,popdark]{\small $S(1\,;\,1)$};
\fill[poporange] (1,1) circle (1.5pt);
\end{tikzpicture}
\end{center}
\legende{Courbe de $g(x)=x^2-2x+2$}
\end{popfigure}
\begin{enumerate}
\item Émettre une conjecture sur le signe de $g$ sur $\mathbb{R}$.
\item Démontrer cette conjecture.
\end{enumerate}
\tcblower
\textbf{1. Conjecture.} La courbe $\mathcal{C}_g$ semble située entièrement \cle{au-dessus} de l'axe des abscisses : son point le plus bas est $S(1\,;\,1)$, d'ordonnée strictement positive. On conjecture donc : $\forall x\in\mathbb{R},\ g(x)>0$.

\textbf{2. Démonstration.} La forme développée ne permet pas de conclure directement ; on utilise la \cle{forme canonique}. Pour tout réel $x$ :
\begin{align*}
g(x) &= x^2-2x+2 = (x-1)^2 - 1 + 2 = (x-1)^2 + 1.
\end{align*}
Or, pour tout réel $x$, $(x-1)^2 \geq 0$, donc $(x-1)^2 + 1 \geq 1 > 0$.
Ainsi $\forall x\in\mathbb{R},\ g(x) \geq 1 > 0$ : la fonction $g$ est strictement positive sur $\mathbb{R}$.

\textbf{Conclusion :} la conjecture émise à partir du graphique est \textbf{validée}. On retient la démarche : la forme canonique $a(x-\alpha)^2+\beta$ donne immédiatement le signe d'un trinôme lorsque $\beta>0$ et $a>0$.
\end{exoresolu}

\begin{savaistu}[Modélisation et conjecture : le coût moyen à la CAMTEL]
Dans la gestion d'un réseau de télécommunications (type CAMTEL ou MTN), le coût total $C(x)$ pour $x$ abonnés comporte souvent une part fixe (infrastructures) et une part variable. Le coût moyen $\bar{C}(x)=\frac{C(x)}{x}$ admet souvent une asymptote horizontale égale au coût variable unitaire. Observer la courbe de $\bar{C}$ permet de conjecturer ce coût limite, puis de le vérifier par le calcul de $\lim_{x\to+\infty}\bar{C}(x)$. C'est une application concrète de la démarche « conjecturer $\to$ vérifier ».
\end{savaistu}

\coursec{Axe de symétrie d'une courbe}

\begin{definition}[Axe de symétrie]
Soit $f$ une fonction définie sur $D_f$ et $a\in\mathbb{R}$. La droite $\Delta$ d'équation $x=a$ est un \cle{axe de symétrie} de la courbe $\mathcal{C}_f$ si et seulement si :
\begin{enumerate}
\item $D_f$ est symétrique par rapport à $a$ : $\forall x\in D_f,\ 2a-x\in D_f$ (équivalent : si $a+h\in D_f$ alors $a-h\in D_f$).
\item $\forall x\in D_f,\ f(2a-x)=f(x)$ (équivalent : $\forall h,\ f(a+h)=f(a-h)$).
\end{enumerate}
Géométriquement, pour tout point $M$ de la courbe, son symétrique $M'$ par rapport à $\Delta$ appartient aussi à la courbe.
\end{definition}

\begin{propriete}[Axe de symétrie d'une parabole]
Pour toute fonction polynôme du second degré $f(x)=ax^2+bx+c$ ($a\neq 0$), la droite d'équation $x=-\dfrac{b}{2a}$ (abscisse du sommet) est un axe de symétrie de $\mathcal{C}_f$. C'est le premier réflexe à avoir face à un trinôme.
\end{propriete}

\begin{methode}[Démontrer qu'une droite $x=a$ est axe de symétrie]
\begin{enumerate}
\item Vérifier que l'ensemble de définition $D_f$ est \textbf{symétrique par rapport à $a$} : si $a+h\in D_f$ alors $a-h\in D_f$.
\item Poser le changement de variable $x = a + h$ (le point d'abscisse $a-h$ est le symétrique de $a+h$ par rapport à $a$).
\item Calculer et montrer que pour tout $h$ tel que $a+h\in D_f$ :
\[ f(a+h) = f(a-h). \]
\item Conclure : les points $M(a+h\,;\,f(a+h))$ et $M'(a-h\,;\,f(a-h))$ ont même ordonnée et des abscisses symétriques par rapport à $a$ : $\Delta$ est axe de symétrie de $\mathcal{C}_f$.
\end{enumerate}
\end{methode}

\begin{exoresolu}[Axe de symétrie d'une parabole]
Soit $f$ la fonction définie sur $\mathbb{R}$ par $f(x)=x^2-4x+1$ et $\mathcal{C}_f$ sa courbe.
Démontrer que la droite $\Delta$ d'équation $x=2$ est axe de symétrie de $\mathcal{C}_f$.
\tcblower
\textbf{Étape 1.} $D_f=\mathbb{R}$ est symétrique par rapport à $2$ : si $2+h\in\mathbb{R}$ alors $2-h\in\mathbb{R}$.

\textbf{Étape 2 et 3.} Soit $h$ un réel quelconque. Calculons :
\begin{align*}
f(2+h) &= (2+h)^2 - 4(2+h) + 1 = 4 + 4h + h^2 - 8 - 4h + 1 = h^2 - 3,\\
f(2-h) &= (2-h)^2 - 4(2-h) + 1 = 4 - 4h + h^2 - 8 + 4h + 1 = h^2 - 3.
\end{align*}
Donc pour tout réel $h$ : $f(2+h)=f(2-h)=h^2-3$.

\textbf{Étape 4. Conclusion.} Les points de $\mathcal{C}_f$ d'abscisses $2+h$ et $2-h$ ont la même ordonnée : ils sont symétriques par rapport à la droite $\Delta$ d'équation $x=2$. La droite $\Delta$ est donc \textbf{axe de symétrie} de $\mathcal{C}_f$.

\emph{Remarque :} on retrouve bien $x=-\dfrac{b}{2a}=-\dfrac{-4}{2}=2$, abscisse du sommet $S(2\,;\,-3)$.
\end{exoresolu}

\coursec{Centre de symétrie d'une courbe}

\begin{definition}[Centre de symétrie]
Soit $f$ une fonction définie sur $D_f$ et $\Omega(a\,;\,b)$ un point du plan. $\Omega$ est un \cle{centre de symétrie} de la courbe $\mathcal{C}_f$ si et seulement si :
\begin{enumerate}
\item $D_f$ est symétrique par rapport à $a$ : $\forall x\in D_f,\ 2a-x\in D_f$.
\item $\forall x\in D_f,\ f(2a-x)+f(x)=2b$ (équivalent : $\forall h,\ f(a+h)+f(a-h)=2b$).
\end{enumerate}
Géométriquement, $\Omega$ est le milieu de tout segment $[MM']$ où $M$ et $M'$ sont deux points de la courbe symétriques par rapport à $\Omega$.
\end{definition}

\begin{propriete}[Centre de symétrie d'une fonction homographique]
Toute fonction homographique $f(x)=\dfrac{ax+b}{cx+d}$ ($ad-bc\neq 0$, $c\neq 0$) admet un centre de symétrie. En effectuant la division (ou identification) pour écrire $f(x)=\alpha+\dfrac{\beta}{x-x_0}$, on trouve le centre $\Omega(x_0\,;\,\alpha)$. La courbe s'obtient par translation de l'hyperbole $y=\frac{\beta}{x}$ (centre $(0;0)$) par le vecteur $\vec{u}(x_0\,;\,\alpha)$.
\end{propriete}

\begin{methode}[Démontrer qu'un point $\Omega(a\,;\,b)$ est centre de symétrie]
\begin{enumerate}
\item Vérifier que $D_f$ est symétrique par rapport à $a$.
\item Pour tout $h$ tel que $a+h\in D_f$, calculer $f(a+h)$ et $f(a-h)$, puis établir l'une des deux égalités équivalentes :
\[ f(a+h)+f(a-h)=2b \qquad\text{ou}\qquad f(a+h)-b = -\bigl(f(a-h)-b\bigr). \]
\item Conclure : $\Omega$ est le milieu de tout couple de points symétriques $M(a+h\,;\,f(a+h))$ et $M'(a-h\,;\,f(a-h))$, donc $\Omega$ est centre de symétrie de $\mathcal{C}_f$.
\end{enumerate}
\end{methode}

\begin{attention}[Ne pas confondre les deux égalités]
Axe $x=a$ : $f(a+h)=f(a-h)$ (égalité des ordonnées). Centre $(a\,;\,b)$ : $f(a+h)+f(a-h)=2b$ (la \textbf{somme} vaut $2b$). Écrire $f(a+h)=-f(a-h)$ ne convient que si $b=0$ (centre à l'origine).
\end{attention}

\begin{exoresolu}[Centre de symétrie d'une courbe]
Soit $f$ la fonction définie sur $\mathbb{R}\setminus\{1\}$ par $f(x)=\dfrac{2x-1}{x-1}$.
Démontrer que le point $\Omega(1\,;\,2)$ est centre de symétrie de la courbe $\mathcal{C}_f$.
\tcblower
\textbf{Étape 1.} $D_f=\mathbb{R}\setminus\{1\}$. Si $1+h\in D_f$ (c'est-à-dire $h\neq 0$), alors $1-h\in D_f$ car $1-h\neq 1$. Le domaine est symétrique par rapport à $1$.

\textbf{Étape 2.} Soit $h\neq 0$. Calculons :
\begin{align*}
f(1+h) &= \frac{2(1+h)-1}{(1+h)-1} = \frac{1+2h}{h} = 2 + \frac{1}{h},\\[2pt]
f(1-h) &= \frac{2(1-h)-1}{(1-h)-1} = \frac{1-2h}{-h} = 2 - \frac{1}{h}.
\end{align*}
Donc :
\[ f(1+h)+f(1-h) = \left(2+\frac{1}{h}\right)+\left(2-\frac{1}{h}\right) = 4 = 2\times 2. \]

\textbf{Étape 3. Conclusion.} Pour tout $h\neq 0$, $f(1+h)+f(1-h)=2\times 2$ : le point $\Omega(1\,;\,2)$ est le milieu des points $M(1+h\,;\,f(1+h))$ et $M'(1-h\,;\,f(1-h))$. Donc $\Omega(1\,;\,2)$ est \textbf{centre de symétrie} de $\mathcal{C}_f$.

\emph{Remarque :} on pouvait aussi réécrire $f(x)=2+\dfrac{1}{x-1}$ (division ou identification) : la courbe se déduit de l'hyperbole $y=\dfrac1x$ par la translation de vecteur $\vec{u}(1\,;\,2)$, qui envoie le centre de symétrie $(0\,;\,0)$ sur $\Omega(1\,;\,2)$.
\end{exoresolu}

\begin{experience}[Découverte : centre de symétrie d'une fonction cubique]
Soit $f(x)=x^3-3x^2+2x$. On cherche un centre de symétrie.
\begin{enumerate}
\item Calculer $f(x)+f(2-x)$ et simplifier.
\item En déduire les coordonnées du centre de symétrie $\Omega$.
\item Vérifier que $f'(x)$ admet $x=1$ pour axe de symétrie (lien avec le point d'inflexion).
\end{enumerate}
\tcblower
\textbf{Solution guidée.}
1. $f(2-x) = (2-x)^3 - 3(2-x)^2 + 2(2-x) = (8-12x+6x^2-x^3) - 3(4-4x+x^2) + 4-2x = -x^3+3x^2-2x$.
   $f(x)+f(2-x) = (x^3-3x^2+2x) + (-x^3+3x^2-2x) = 0 = 2\times 0$.
2. On a $f(1+h)+f(1-h)=0=2\times 0$. Le centre est $\Omega(1\,;\,0)$.
3. $f'(x)=3x^2-6x+2$. $f'(1+h)=3(1+h)^2-6(1+h)+2=3h^2-1$. $f'(1-h)=3h^2-1$. $f'$ est pair par rapport à $1$ : $x=1$ est axe de symétrie de $\mathcal{C}_{f'}$, ce qui correspond au point d'inflexion de $\mathcal{C}_f$ d'abscisse $1$.
\end{experience}

\coursec{Asymptotes verticales et horizontales}

\begin{definition}[Asymptotes verticale et horizontale]
Soit $f$ une fonction et $\mathcal{C}_f$ sa courbe.
\begin{itemize}
\item La droite $\Delta$ d'équation $x=a$ est une \cle{asymptote verticale} si $\displaystyle\lim_{x\to a}f(x)=+\infty$ ou $-\infty$ (ou limites latérales infinies).
\item La droite $\Delta$ d'équation $y=b$ est une \cle{asymptote horizontale} en $+\infty$ (resp. $-\infty$) si $\displaystyle\lim_{x\to+\infty}f(x)=b$ (resp. $\lim_{x\to-\infty}f(x)=b$) avec $b\in\mathbb{R}$ fini.
\end{itemize}
\end{definition}

\begin{propriete}[Asymptotes des fonctions rationnelles]
Pour $f(x)=\frac{P(x)}{Q(x)}$ avec $P,Q$ polynômes :
\begin{itemize}
\item \textbf{Verticale :} aux zéros de $Q$ qui ne sont pas des zéros de $P$ (ou d'ordre inférieur).
\item \textbf{Horizontale :} si $\deg P < \deg Q$, $y=0$ ; si $\deg P = \deg Q$, $y=\frac{a_n}{b_n}$ (quotient des coefficients dominants) ; si $\deg P > \deg Q$, pas d'asymptote horizontale (asymptote oblique possible).
\end{itemize}
\end{propriete}

\begin{methode}[Asymptotes verticales et horizontales par les limites]
\begin{enumerate}
\item \textbf{Asymptote verticale} : repérer les valeurs $a$ exclues de $D_f$ (dénominateur nul, racine d'un radicand pair, etc.). Calculer $\displaystyle\lim_{x\to a} f(x)$ (éventuellement limites à gauche et à droite). Si cette limite est $+\infty$ ou $-\infty$, alors la droite d'équation \cle{$x=a$} est asymptote verticale à $\mathcal{C}_f$.
\item \textbf{Asymptote horizontale} : calculer $\displaystyle\lim_{x\to+\infty}f(x)$ et $\displaystyle\lim_{x\to-\infty}f(x)$. Si l'une vaut un réel $b$ \textbf{fini}, alors la droite d'équation \cle{$y=b$} est asymptote horizontale à $\mathcal{C}_f$ en $+\infty$ (ou $-\infty$).
\item \textbf{Position relative} (souvent demandée) : étudier le signe de $f(x)-b$. Si $f(x)-b>0$, la courbe est au-dessus de l'asymptote ; si $f(x)-b<0$, elle est en dessous.
\end{enumerate}
\begin{aretenir}[Formulation exigée au Bac]
Rédiger toujours : « $\displaystyle\lim_{x\to a}f(x)=+\infty$, donc la droite d'équation $x=a$ est asymptote verticale à $\mathcal{C}_f$ ». Une limite infinie \emph{en un point fini} donne une asymptote \emph{verticale} ; une limite \emph{finie à l'infini} donne une asymptote \emph{horizontale}.
\end{aretenir}
\end{methode}

\begin{exoresolu}[Asymptotes d'une fonction homographique]
Soit $f$ définie sur $\mathbb{R}\setminus\{-2\}$ par $f(x)=\dfrac{3x+1}{x+2}$, et $\mathcal{C}_f$ sa courbe.
\begin{enumerate}
\item Montrer que $\mathcal{C}_f$ admet une asymptote verticale dont on donnera l'équation.
\item Montrer que $\mathcal{C}_f$ admet une asymptote horizontale en $+\infty$ et en $-\infty$.
\item Préciser la position de $\mathcal{C}_f$ par rapport à son asymptote horizontale.
\end{enumerate}
\tcblower
\textbf{1. Asymptote verticale.} La valeur interdite est $x=-2$ (elle annule le dénominateur). Étudions la limite en $-2$. Le numérateur tend vers $3\times(-2)+1=-5\neq 0$.

Pour $x>-2$, $x+2>0$, donc $\displaystyle\lim_{\substack{x\to -2\\x>-2}}f(x)=\frac{-5}{0^+}=-\infty$.

Pour $x<-2$, $x+2<0$, donc $\displaystyle\lim_{\substack{x\to -2\\x<-2}}f(x)=\frac{-5}{0^-}=+\infty$.

Les limites étant infinies, la droite d'équation \textbf{$x=-2$ est asymptote verticale} à $\mathcal{C}_f$.

\textbf{2. Asymptote horizontale.} Pour $x\neq 0$, factorisons par les termes dominants :
\[ f(x)=\frac{3x+1}{x+2}=\frac{x\left(3+\frac1x\right)}{x\left(1+\frac2x\right)}=\frac{3+\frac1x}{1+\frac2x}. \]
Comme $\displaystyle\lim_{x\to\pm\infty}\frac1x=0$ et $\displaystyle\lim_{x\to\pm\infty}\frac2x=0$, on obtient :
\[ \lim_{x\to+\infty}f(x)=3 \qquad\text{et}\qquad \lim_{x\to-\infty}f(x)=3. \]
Ces limites sont finies : la droite d'équation \textbf{$y=3$ est asymptote horizontale} à $\mathcal{C}_f$ en $+\infty$ et en $-\infty$.

\textbf{3. Position relative.} Écrivons $f(x)=3+\dfrac{?}{x+2}$ : par identification, $f(x)=\dfrac{3(x+2)-5}{x+2}=3-\dfrac{5}{x+2}$. Donc :
\[ f(x)-3 = -\frac{5}{x+2}. \]
Si $x>-2$ : $f(x)-3<0$, la courbe est \textbf{en dessous} de l'asymptote. Si $x<-2$ : $f(x)-3>0$, la courbe est \textbf{au-dessus} de l'asymptote.

\textbf{Conclusion :} $\mathcal{C}_f$ admet les asymptotes $x=-2$ et $y=3$ ; elle est au-dessus de $y=3$ sur $]-\infty\,;\,-2[$ et en dessous sur $]-2\,;\,+\infty[$.
\end{exoresolu}

\coursec{Asymptotes obliques}

\begin{definition}[Asymptote oblique]
Soit $f$ une fonction et $\mathcal{C}_f$ sa courbe. La droite $\Delta$ d'équation $y=ax+b$ (avec $a\neq 0$) est une \cle{asymptote oblique} à $\mathcal{C}_f$ en $+\infty$ (resp. $-\infty$) si et seulement si :
\[ \lim_{x\to+\infty}\bigl[f(x)-(ax+b)\bigr]=0 \qquad\bigl(\text{resp. } \lim_{x\to-\infty}\bigl[f(x)-(ax+b)\bigr]=0\bigr). \]
\end{definition}

\begin{propriete}[Recherche des coefficients $a$ et $b$]
Si $\mathcal{C}_f$ admet une asymptote oblique $y=ax+b$ en $+\infty$, alors nécessairement :
\[ a = \lim_{x\to+\infty}\frac{f(x)}{x} \quad\text{(fini et non nul)},\qquad b = \lim_{x\to+\infty}\bigl[f(x)-ax\bigr] \quad\text{(fini)}. \]
Même procédure en $-\infty$. Pour une fonction rationnelle $\frac{P}{Q}$, il y a asymptote oblique en $\pm\infty$ si $\deg P = \deg Q + 1$ ; $a$ et $b$ sont le quotient et le reste de la division euclidienne de $P$ par $Q$.
\end{propriete}

\begin{methode}[Déterminer une asymptote oblique $y=ax+b$]
\begin{enumerate}
\item \textbf{Si $a$ et $b$ sont donnés} : former la différence $f(x)-(ax+b)$ et calculer sa limite. Si elle vaut $0$, $\Delta$ est asymptote oblique.
\item \textbf{Si $a$ et $b$ sont à trouver} :
      \begin{enumerate}
      \item Calculer $a=\displaystyle\lim_{x\to+\infty}\frac{f(x)}{x}$. Si la limite n'existe pas ou est $0$ ou infinie : pas d'asymptote oblique.
      \item Calculer $b=\displaystyle\lim_{x\to+\infty}\bigl[f(x)-ax\bigr]$. Si la limite n'existe pas ou est infinie : pas d'asymptote oblique.
      \item Vérifier que $\lim_{x\to+\infty}[f(x)-(ax+b)]=0$ (souvent automatique si les deux limites précédentes sont finies).
      \end{enumerate}
\item \textbf{Position relative} : étudier le signe de $f(x)-(ax+b)$ au voisinage de l'infini.
\end{enumerate}
\end{methode}

\begin{exoresolu}[Asymptote oblique]
Soit $f$ définie sur $]0\,;\,+\infty[$ par $f(x)=2x-1+\dfrac{3}{x}$, de courbe $\mathcal{C}_f$.
\begin{enumerate}
\item Démontrer que la droite $\Delta$ d'équation $y=2x-1$ est asymptote à $\mathcal{C}_f$ en $+\infty$.
\item Préciser la position de $\mathcal{C}_f$ par rapport à $\Delta$.
\end{enumerate}
\tcblower
\textbf{1.} Formons la différence entre la fonction et la droite :
\[ f(x)-(2x-1) = 2x-1+\frac{3}{x}-2x+1 = \frac{3}{x}. \]
Or $\displaystyle\lim_{x\to+\infty}\frac{3}{x}=0$, donc :
\[ \lim_{x\to+\infty}\bigl[f(x)-(2x-1)\bigr]=0. \]
Par définition, la droite $\Delta$ d'équation \textbf{$y=2x-1$ est asymptote oblique} à $\mathcal{C}_f$ en $+\infty$.

\textbf{2. Position relative.} Pour tout $x>0$, $f(x)-(2x-1)=\dfrac{3}{x}>0$.

Donc sur tout son ensemble de définition $]0\,;\,+\infty[$, la courbe $\mathcal{C}_f$ est située \textbf{strictement au-dessus} de son asymptote $\Delta$, tout en s'en rapprochant indéfiniment quand $x$ tend vers $+\infty$.
\end{exoresolu}

\begin{experience}[Recherche guidée d'une asymptote oblique]
Soit $f(x)=\dfrac{x^2+1}{x-1}$ sur $\mathbb{R}\setminus\{1\}$. On cherche l'asymptote oblique en $+\infty$.
\begin{enumerate}
\item Effectuer la division euclidienne de $x^2+1$ par $x-1$ pour écrire $f(x)=ax+b+\frac{r}{x-1}$.
\item En déduire $a$ et $b$ et vérifier la définition.
\item Étudier la position relative.
\end{enumerate}
\tcblower
\textbf{Solution.}
1. Division : $x^2+1 = (x-1)(x+1) + 2$. Donc $f(x)=x+1+\frac{2}{x-1}$. On conjecture $a=1, b=1$.
2. $f(x)-(x+1)=\frac{2}{x-1} \xrightarrow[x\to+\infty]{} 0$. Asymptote : $y=x+1$.
3. Pour $x>1$, $f(x)-(x+1)=\frac{2}{x-1}>0$ : courbe au-dessus. Pour $x<1$, $f(x)-(x+1)<0$ : courbe en dessous.
\end{experience}

\coursec{Synthèse : construire une courbe complète}

\begin{methode}[Tracer une courbe à partir du tableau de variations et des asymptotes]
\begin{enumerate}
\item \textbf{Rassembler les informations} : ensemble de définition, limites aux bornes, tableau de variations, asymptotes (verticales, horizontales, obliques), éléments de symétrie (axe ou centre), points remarquables (intersections avec les axes, extrema).
\item \textbf{Tracer le repère et les éléments fixes} : asymptotes (en pointillés de couleur distincte), axes ou centre de symétrie, tangentes horizontales aux extrema.
\item \textbf{Placer les points clés} : extrema (coordonnées exactes), intersections avec les axes (si simples), points d'inflexion éventuels.
\item \textbf{Esquisser les branches} en respectant scrupuleusement :
      \begin{itemize}
      \item le sens de variation entre deux points clés (croissant/décroissant) ;
      \item le « collage » aux asymptotes (la courbe s'en approche sans les toucher, sauf position relative étudiée) ;
      \item la symétrie éventuelle (tracer une moitié puis compléter par symétrie axiale ou centrale).
      \end{itemize}
\item \textbf{Vérifier la cohérence} : chaque branche doit être compatible avec les limites du tableau de variations (valeurs aux bornes des intervalles de monotonie).
\end{enumerate}
\end{methode}

\begin{exoresolu}[Construction complète d'une courbe]
Soit $f$ définie sur $\mathbb{R}\setminus\{1\}$ par $f(x)=\dfrac{x^2-3x+4}{x-1}$. On admet que :
\begin{itemize}
\item $\displaystyle\lim_{x\to\pm\infty}f(x)=\pm\infty$ et la droite $\Delta : y=x-2$ est asymptote oblique en $+\infty$ et $-\infty$ ;
\item la droite $x=1$ est asymptote verticale, avec $\displaystyle\lim_{\substack{x\to1\\x<1}}f(x)=-\infty$ et $\displaystyle\lim_{\substack{x\to1\\x>1}}f(x)=+\infty$ ;
\item $f$ est croissante sur $]-\infty\,;\,-1]$ et sur $[3\,;\,+\infty[$, décroissante sur $[-1\,;\,1[$ et sur $]1\,;\,3]$ ; $f(-1)=-4$ et $f(3)=2$.
\end{itemize}
Tracer l'allure de $\mathcal{C}_f$ avec ses asymptotes.
\tcblower
\textbf{Étape 1 — Bilan.} Deux asymptotes : verticale $x=1$, oblique $y=x-2$. Deux extrema : maximum local $A(-1\,;\,-4)$, minimum local $B(3\,;\,2)$. Quatre branches à construire, séparées par l'asymptote verticale.

\textbf{Étape 2 — Éléments fixes.} On trace en pointillés les droites $x=1$ et $y=x-2$.

\textbf{Étape 3 — Points clés.} On place $A(-1\,;\,-4)$ (tangente horizontale) et $B(3\,;\,2)$ (tangente horizontale).

\textbf{Étape 4 — Branches.}
\begin{itemize}
\item Sur $]-\infty\,;\,-1]$ : $f$ croît de $-\infty$ (en suivant $\Delta$) jusqu'à $-4$ en $A$.
\item Sur $[-1\,;\,1[$ : $f$ décroît de $-4$ vers $-\infty$ en « collant » à la droite $x=1$ par la gauche.
\item Sur $]1\,;\,3]$ : $f$ décroît de $+\infty$ (le long de $x=1$ par la droite) jusqu'à $2$ en $B$.
\item Sur $[3\,;\,+\infty[$ : $f$ croît de $2$ vers $+\infty$ en se rapprochant de $\Delta$.
\end{itemize}

\begin{popfigure}
\begin{center}
\begin{tikzpicture}[scale=0.62]
\draw[->,popmuted] (-6,0) -- (7,0) node[below]{$x$};
\draw[->,popmuted] (0,-7) -- (0,6) node[left]{$y$};
\draw[poporange,dashed,thick] (1,-7) -- (1,6) node[above left]{\small $x=1$};
\draw[popgreen,dashed,thick,domain=-5.5:6.5] plot (\x,{\x-2}) node[below right,pos=0.95]{\small $y=x-2$};
\draw[popblue,thick,domain=-5.5:-1.05,samples=80] plot (\x,{(\x*\x-3*\x+4)/(\x-1)});
\draw[popblue,thick,domain=-0.98:0.82,samples=80] plot (\x,{(\x*\x-3*\x+4)/(\x-1)});
\draw[popblue,thick,domain=1.18:2.95,samples=80] plot (\x,{(\x*\x-3*\x+4)/(\x-1)});
\draw[popblue,thick,domain=3.05:6.4,samples=80] plot (\x,{(\x*\x-3*\x+4)/(\x-1)});
\fill[popdark] (-1,-4) circle (2pt) node[below left]{\small $A(-1;-4)$};
\fill[popdark] (3,2) circle (2pt) node[above right]{\small $B(3;2)$};
\end{tikzpicture}
\end{center}
\legende{Courbe de $f$ avec ses deux asymptotes}
\end{popfigure}

\textbf{Étape 5 — Vérification.} Chaque branche respecte le sens de variation annoncé et les limites : la courbe rejoint bien $-\infty$ puis $+\infty$ de part et d'autre de $x=1$, et se rapproche de $\Delta$ aux deux infinis. Le tracé est cohérent.
\end{exoresolu}

\begin{attention}[Les 5 erreurs qui coûtent des points au Bac]
\begin{enumerate}
\item \textbf{Confondre les formules de symétrie.} Axe $x=a$ : $f(a+h)=f(a-h)$. Centre $(a\,;\,b)$ : $f(a+h)+f(a-h)=2b$. Écrire la mauvaise égalité (ou oublier le facteur $2$ devant $b$) fausse toute la démonstration.
\item \textbf{Oublier de vérifier le domaine de définition.} Avant toute démonstration de symétrie, il faut écrire que $D_f$ est symétrique par rapport à $a$ ; sans cette phrase, la démonstration est incomplète et pénalisée.
\item \textbf{Affirmer une asymptote sans calcul de limite.} « La droite $x=2$ est asymptote car la fonction n'est pas définie en $2$ » est faux en général : seule une \textbf{limite infinie} en $2$ justifie l'asymptote verticale, et seule une \textbf{limite finie} à l'infini justifie une asymptote horizontale. Toujours écrire la limite avant de conclure.
\item \textbf{Conclure sur une forme indéterminée sans la lever.} Pour $\displaystyle\lim_{x\to+\infty}\frac{3x+1}{x+2}$, écrire « $\frac{\infty}{\infty}$ » ne vaut rien : il faut factoriser par $x$ (ou mettre le terme dominant en facteur) puis conclure.
\item \textbf{Tracer une courbe qui traverse son asymptote verticale ou qui contredit le tableau de variations.} Une branche ne franchit jamais une asymptote verticale (valeur interdite !), et chaque portion tracée doit respecter la monotonie et les limites annoncées. Relire le tableau avant de tracer.
\end{enumerate}
\end{attention}

\coursec{Exercices d'entraînement}

\begin{exercice}[Symétrie et asymptotes]
Soit $f$ la fonction définie sur $\mathbb{R}\setminus\{-1\}$ par $f(x)=\dfrac{x^2+2x+3}{x+1}$.
\begin{enumerate}
\item Déterminer les asymptotes de $\mathcal{C}_f$ (verticale et oblique).
\item Étudier la position de $\mathcal{C}_f$ par rapport à son asymptote oblique.
\item Démontrer que $\mathcal{C}_f$ admet un centre de symétrie et donner ses coordonnées.
\item Tracer $\mathcal{C}_f$ et ses asymptotes.
\end{enumerate}
\end{exercice}

\begin{exercice}[Conjecture et vérification]
On donne la courbe $\mathcal{C}_h$ d'une fonction $h$ définie sur $\mathbb{R}^*_+$ par $h(x)=x+\dfrac{4}{x}$.
\begin{enumerate}
\item À l'aide de la courbe (fournie à l'examen), conjecturer l'existence d'un minimum et sa valeur approchée.
\item Démontrer par le calcul que $h$ admet un minimum sur $\mathbb{R}^*_+$ et le déterminer exactement.
\item $\mathcal{C}_h$ admet-elle un axe de symétrie ? Un centre de symétrie ? Justifier.
\end{enumerate}
\end{exercice}

\begin{exercice}[Asymptotes et variations]
Soit $f(x)=\dfrac{2x^2-3x+5}{x-2}$ sur $\mathbb{R}\setminus\{2\}$.
\begin{enumerate}
\item Déterminer les asymptotes de $\mathcal{C}_f$.
\item Calculer $f'(x)$ et dresser le tableau de variations de $f$.
\item Tracer $\mathcal{C}_f$ et ses asymptotes.
\end{enumerate}
\end{exercice}

\begin{exercice}[Application concrète : Coût moyen]
Une entreprise de transformation agricole à Bafoussam a un coût total de production (en milliers de FCFA) modélisé par $C(x)=2x^2+50x+2000$ pour $x$ tonnes de produit ($x>0$).
\begin{enumerate}
\item Exprimer le coût moyen $\bar{C}(x)$.
\item Déterminer les asymptotes de la courbe de $\bar{C}$.
\item Étudier les variations de $\bar{C}$ et interpréter économiquement le minimum du coût moyen.
\end{enumerate}
\end{exercice}

\newpage

\coursec{Exercices}

\coursec{Courbes représentatives}

\courssub{Je vérifie mes connaissances}

\begin{aretenir}[Rappels essentiels : asymptotes et symétries]
\textbf{Asymptote verticale :} La droite $x=a$ est asymptote verticale à $\mathcal{C}_f$ si $\lim\limits_{x\to a^-}f(x)=\pm\infty$ ou $\lim\limits_{x\to a^+}f(x)=\pm\infty$.

\textbf{Asymptote horizontale :} La droite $y=b$ est asymptote horizontale à $\mathcal{C}_f$ en $+\infty$ (resp. $-\infty$) si $\lim\limits_{x\to +\infty}f(x)=b$ (resp. $\lim\limits_{x\to -\infty}f(x)=b$).

\textbf{Asymptote oblique :} La droite $y=ax+b$ ($a\neq 0$) est asymptote oblique à $\mathcal{C}_f$ en $+\infty$ (resp. $-\infty$) si $\lim\limits_{x\to +\infty}\bigl[f(x)-(ax+b)\bigr]=0$ (resp. $\lim\limits_{x\to -\infty}\bigl[f(x)-(ax+b)\bigr]=0$).

\textbf{Axe de symétrie :} La droite $x=a$ est axe de symétrie de $\mathcal{C}_f$ si $\forall h$, $f(a+h)=f(a-h)$ (fonction paire par rapport à $a$).

\textbf{Centre de symétrie :} Le point $\Omega(a\,;b)$ est centre de symétrie de $\mathcal{C}_f$ si $\forall h$, $f(a+h)+f(a-h)=2b$ (fonction impaire par rapport à $\Omega$).
\end{aretenir}

\begin{exercice}[QCM : une seule bonne réponse]
Pour chaque question, choisir la bonne réponse (aucune justification n'est demandée ici, mais sois capable de la donner à l'oral).

\textbf{1.} La courbe de la fonction $f$ définie par $f(x)=\dfrac{3x+1}{x-2}$ admet pour asymptote verticale la droite d'équation :

\quad a) $x=-2$ \qquad b) $x=2$ \qquad c) $y=2$ \qquad d) $y=3$

\textbf{2.} Si $\displaystyle\lim_{x\to +\infty}f(x)=5$, alors la courbe de $f$ admet :

\quad a) une asymptote verticale d'équation $x=5$ \qquad b) une asymptote horizontale d'équation $y=5$

\quad c) un centre de symétrie de coordonnées $(5\,;0)$ \qquad d) une asymptote oblique

\textbf{3.} Le point $\Omega(a\,;b)$ est centre de symétrie de la courbe de $f$ si et seulement si, pour tout $x$ tel que $a+x$ et $a-x$ appartiennent au domaine :

\quad a) $f(a+x)=f(a-x)$ \qquad b) $f(a+x)+f(a-x)=2b$

\quad c) $f(a+x)-f(a-x)=2b$ \qquad d) $f(2a-x)=f(x)$

\textbf{4.} La droite d'équation $y=2x+1$ est asymptote oblique à la courbe de $f$ en $+\infty$ si :

\quad a) $\displaystyle\lim_{x\to+\infty}f(x)=+\infty$ \qquad b) $\displaystyle\lim_{x\to+\infty}[f(x)-2x]=1$

\quad c) $\displaystyle\lim_{x\to+\infty}[f(x)-(2x+1)]=0$ \qquad d) $f(x)=2x+1$ pour $x$ assez grand
\end{exercice}

\begin{exercice}[Vrai ou faux, en justifiant]
Pour chaque affirmation, dire si elle est vraie ou fausse, en justifiant par une démonstration courte ou un contre-exemple.

\textbf{1.} « Si $a$ est une valeur interdite de $f$, alors la droite d'équation $x=a$ est asymptote verticale à la courbe de $f$. »

\textbf{2.} « Une courbe ne peut jamais couper son asymptote horizontale. »

\textbf{3.} « Si la courbe de $f$ admet le point $\Omega(1\,;2)$ pour centre de symétrie, alors $f(1+h)+f(1-h)=4$ pour tout réel $h$ tel que $1+h$ et $1-h$ soient dans le domaine de $f$. »

\textbf{4.} « Une courbe peut admettre à la fois une asymptote horizontale en $+\infty$ et une asymptote oblique en $-\infty$. »
\end{exercice}

\begin{exercice}[Compléter les définitions]
Recopier et compléter les phrases suivantes avec les mots ou expressions qui conviennent.

\textbf{1.} La droite d'équation $x=a$ est \cle{asymptote verticale} à la courbe de $f$ si $\displaystyle\lim_{x\to a}f(x)=+\infty$ ou $-\infty$ (ou l'une des limites unilatérales est infinie).

\textbf{2.} La droite d'équation $y=b$ est \cle{asymptote horizontale} à la courbe de $f$ en $+\infty$ si $\displaystyle\lim_{x\to+\infty}f(x)=b$.

\textbf{3.} La droite d'équation $x=a$ est \cle{axe de symétrie} de la courbe de $f$ si, pour tout $h$ convenable : $f(a+h)=f(a-h)$.

\textbf{4.} Pour montrer que la droite d'équation $y=ax+b$ est asymptote oblique à la courbe de $f$ en $+\infty$, on calcule la limite de $f(x)-(ax+b)$ et on vérifie qu'elle vaut $0$.
\end{exercice}

\begin{exercice}[Calculs directs de limites]
Dans chaque cas, calculer la limite demandée et en déduire l'équation de l'asymptote correspondante à la courbe de $f$, en précisant sa nature (verticale, horizontale ou oblique).

\textbf{1.} $f(x)=\dfrac{2x-3}{x+1}$ ; calculer $\displaystyle\lim_{x\to -1}f(x)$ et $\displaystyle\lim_{x\to +\infty}f(x)$.

\textbf{2.} $f(x)=x+1+\dfrac{2}{x}$ ; calculer $\displaystyle\lim_{x\to +\infty}[f(x)-(x+1)]$.

\textbf{3.} $f(x)=\dfrac{x^2+1}{x}$ ; calculer $\displaystyle\lim_{x\to 0}f(x)$ (on distinguera $x\to 0^+$ et $x\to 0^-$).
\end{exercice}

\courssub{Je m'entraîne}

\begin{exercice}[Axe de symétrie d'une parabole]
Soit $f$ la fonction définie sur $\mathbb{R}$ par $f(x)=x^2+2x-3$ et $\mathcal{C}_f$ sa courbe représentative dans un repère orthonormé.

\textbf{1.} Montrer que pour tout réel $h$ : $f(-1+h)=f(-1-h)$.

\textbf{2.} Qu'en déduit-on pour la courbe $\mathcal{C}_f$ ?

\textbf{3.} Dresser le tableau de variations de $f$, puis tracer $\mathcal{C}_f$ et son axe de symétrie (unité : $\SI{1}{cm}$).
\end{exercice}

\begin{exercice}[Centre de symétrie et réduction du domaine d'étude]
Soit $f$ la fonction définie par $f(x)=\dfrac{x-1}{x-2}$ et $\mathcal{C}_f$ sa courbe représentative.

\textbf{1.} Déterminer le domaine de définition $D_f$ de $f$.

\textbf{2.} Montrer que pour tout réel $h\neq 0$ : $f(2+h)+f(2-h)=2$.

\textbf{3.} En déduire que le point $\Omega(2\,;1)$ est centre de symétrie de $\mathcal{C}_f$.

\textbf{4.} Expliquer pourquoi il suffit d'étudier $f$ sur $]2\,;+\infty[$ pour obtenir toute la courbe. Préciser la transformation qui permet de compléter le tracé.
\end{exercice}

\begin{exercice}[Asymptotes d'une fonction homographique]
Une station-service de la ville de Bafoussam modélise le coût unitaire de stockage du carburant (en milliers de FCFA par mètre cube) par la fonction $f$ définie sur $]0\,;+\infty[$ par $f(x)=\dfrac{4x+8}{x}$, où $x$ est la quantité stockée en $\mathrm{m}^3$.

\textbf{1.} Calculer $\displaystyle\lim_{x\to 0^+}f(x)$. Interpréter graphiquement le résultat.

\textbf{2.} Calculer $\displaystyle\lim_{x\to +\infty}f(x)$. En déduire une asymptote à la courbe de $f$ et interpréter ce résultat en termes de coût unitaire.

\textbf{3.} Montrer que le point $\Omega(0\,;4)$ n'appartient pas à la courbe. Déterminer le centre de symétrie de $\mathcal{C}_f$ (on prolonge $f$ sur $\mathbb{R}\setminus\{0\}$ par la même formule).
\end{exercice}

\begin{exercice}[Lecture graphique et conjectures]
On donne ci-dessous la courbe représentative $\mathcal{C}$ d'une fonction $g$ définie sur $\mathbb{R}$.

\begin{popfigure}
\begin{center}
\begin{tikzpicture}[scale=0.9]
\draw[popline] (-4.5,0) -- (4.5,0) node[below left]{$x$};
\draw[popline] (0,-0.8) -- (0,3.2) node[below left]{$y$};
\draw[popline] (0,0) node[below left]{$O$};
\draw[popline] (1,0.06) -- (1,-0.06) node[below]{\small $1$};
\draw[popline] (-1,0.06) -- (-1,-0.06) node[below]{\small $-1$};
\draw[popline] (0.06,1) -- (-0.06,1) node[left]{\small $1$};
\draw[popline] (0.06,2) -- (-0.06,2) node[left]{\small $2$};
\draw[popline] (0.06,3) -- (-0.06,3) node[left]{\small $3$};
\draw[dashed,poporange,thick] (-4.5,2) -- (4.5,2);
\begin{scope}
\clip (-4.5,-0.8) rectangle (4.5,3.2);
\draw[popblue,very thick,smooth,domain=-4.5:4.5,samples=200] plot (\x,{(2*\x*\x+3)/(\x*\x+1)});
\end{scope}
\draw[popblue] (3.4,2.25) node[above]{$\mathcal{C}$};
\end{tikzpicture}
\end{center}
\legende{Courbe représentative de la fonction $g$}
\end{popfigure}

\textbf{1.} Conjectures graphiques :
\begin{itemize}
\item Quelle semble être l'asymptote horizontale de $\mathcal{C}$ en $+\infty$ et en $-\infty$ ?
\item Quel semble être l'axe de symétrie de $\mathcal{C}$ ?
\item Quelle semble être la position de $\mathcal{C}$ par rapport à son asymptote ?
\end{itemize}

\textbf{2.} Vérification par le calcul : on sait que $g(x)=\dfrac{2x^2+3}{x^2+1}$.
\begin{itemize}
\item[a)] Calculer $\displaystyle\lim_{x\to +\infty}g(x)$ et valider ou invalider la première conjecture.
\item[b)] Montrer que $g$ est paire et conclure.
\item[c)] Étudier le signe de $g(x)-2$ et conclure sur la position relative.
\end{itemize}
\end{exercice}

\courssub{Je me prépare au Bac}

\begin{exercice}[Type Bac : étude complète avec asymptote oblique]
\textbf{(D'après Baccalauréat A, thème classique.)} Soit $f$ la fonction définie par $f(x)=\dfrac{x^2-3x+4}{x-1}$ et $\mathcal{C}_f$ sa courbe représentative dans un repère orthonormé $(O\,;\vec{i},\vec{j})$ (unité : $\SI{1}{cm}$).

\textbf{Partie A : domaine et limites}

\textbf{1.} Déterminer le domaine de définition $D_f$ de $f$.

\textbf{2.} Calculer les limites de $f$ en $1$ (à gauche et à droite). En déduire une asymptote $\mathcal{D}_1$ à $\mathcal{C}_f$.

\textbf{3.} Calculer les limites de $f$ en $+\infty$ et en $-\infty$.

\textbf{Partie B : asymptote oblique}

\textbf{4.} Montrer que pour tout $x\neq 1$ : $f(x)=x-2+\dfrac{2}{x-1}$.

\textbf{5.} En déduire que la droite $\mathcal{D}_2$ d'équation $y=x-2$ est asymptote oblique à $\mathcal{C}_f$ en $+\infty$ et en $-\infty$.

\textbf{6.} Étudier, suivant les valeurs de $x$, la position relative de $\mathcal{C}_f$ et de $\mathcal{D}_2$.

\textbf{Partie C : variations et tracé}

\textbf{7.} Montrer que pour tout $x\neq 1$ : $f'(x)=\dfrac{x^2-2x-1}{(x-1)^2}$.

\textbf{8.} Étudier le signe de $f'(x)$ et dresser le tableau de variations complet de $f$ (on donnera les valeurs exactes des extrema, atteints en $1-\sqrt{2}$ et $1+\sqrt{2}$).

\textbf{9.} Tracer soigneusement $\mathcal{D}_1$, $\mathcal{D}_2$ et $\mathcal{C}_f$.

\textbf{10.} Déterminer graphiquement, puis algébriquement, le nombre de solutions de l'équation $f(x)=3$.
\end{exercice}

\begin{exercice}[Type Bac : parité, asymptotes et équation]
\textbf{(Contexte : réseau de distribution.)} Une coopérative d'électrification rurale de la région de l'Adamaoua modélise une grandeur physique par la fonction $f$ définie sur $\mathbb{R}\setminus\{0\}$ par :
\[ f(x)=x+\dfrac{4}{x}. \]
On note $\mathcal{C}_f$ sa courbe représentative dans un repère orthonormé $(O\,;\vec{i},\vec{j})$ (unité : $\SI{1}{cm}$).

\textbf{Partie A : éléments de symétrie}

\textbf{1.} Montrer que $f$ est impaire. Que peut-on en déduire pour $\mathcal{C}_f$ ? Justifier qu'on peut restreindre l'étude à $]0\,;+\infty[$.

\textbf{Partie B : limites et asymptotes}

\textbf{2.} Calculer $\displaystyle\lim_{x\to 0^+}f(x)$. En déduire une asymptote à $\mathcal{C}_f$.

\textbf{3.} Calculer $\displaystyle\lim_{x\to +\infty}f(x)$.

\textbf{4.} Montrer que la droite $\mathcal{D}$ d'équation $y=x$ est asymptote oblique à $\mathcal{C}_f$ en $+\infty$.

\textbf{5.} Étudier la position relative de $\mathcal{C}_f$ et de $\mathcal{D}$ sur $]0\,;+\infty[$.

\textbf{Partie C : variations et tracé}

\textbf{6.} Calculer $f'(x)$ et montrer que $f'(x)=\dfrac{x^2-4}{x^2}$.

\textbf{7.} Dresser le tableau de variations de $f$ sur $]0\,;+\infty[$, puis sur $\mathbb{R}\setminus\{0\}$ en utilisant la parité.

\textbf{8.} Tracer $\mathcal{D}$, l'asymptote verticale et $\mathcal{C}_f$ en justifiant le tracé de la partie sur $]-\infty\,;0[$ par symétrie.

\textbf{Partie D : résolution d'une équation}

\textbf{9.} Résoudre graphiquement l'équation $f(x)=5$.

\textbf{10.} Retrouver algébriquement les solutions exactes de $f(x)=5$ (on se ramènera à une équation du second degré).

\textbf{11.} La coopérative estime que la grandeur modélisée est rentable lorsque $f(x)\leq 5$ pour $x>0$. Déterminer l'intervalle des valeurs de $x$ correspondantes.
\end{exercice}

\begin{exercice}[Type Bac : fonction rationnelle et centre de symétrie]
\textbf{(Contexte : transport urbain.)} Une société de transport de Yaoundé étudie le temps moyen de parcours d'une ligne de bus. Une modélisation conduit à la fonction $f$ définie par :
\[ f(x)=\dfrac{2x^2-5x+6}{x-2} \]
où $x$ représente le nombre de bus en circulation ($x>2$) et $f(x)$ le temps moyen en minutes. On note $\mathcal{C}_f$ la courbe de $f$ dans un repère orthogonal $(O\,;\vec{i},\vec{j})$ (unités : $\SI{2}{cm}$ en abscisse, $\SI{0,5}{cm}$ en ordonnée).

\textbf{Partie A : limites et asymptotes}

\textbf{1.} Préciser le domaine de définition de $f$ et calculer les limites de $f$ en $2$ (à gauche et à droite). Interpréter graphiquement.

\textbf{2.} Montrer que pour tout $x\neq 2$ : $f(x)=2x-1+\dfrac{4}{x-2}$.

\textbf{3.} En déduire que la droite $\Delta$ d'équation $y=2x-1$ est asymptote oblique à $\mathcal{C}_f$ en $+\infty$ et en $-\infty$, et étudier la position relative de $\mathcal{C}_f$ par rapport à $\Delta$.

\textbf{Partie B : centre de symétrie}

\textbf{4.} Soit $\Omega$ le point d'intersection des deux asymptotes. Déterminer les coordonnées de $\Omega$.

\textbf{5.} Montrer que pour tout réel $h\neq 0$ : $f(2+h)+f(2-h)=6$.

\textbf{6.} En déduire que $\Omega$ est centre de symétrie de $\mathcal{C}_f$. Comment cela permet-il de réduire le domaine d'étude ?

\textbf{Partie C : variations, tracé et exploitation}

\textbf{7.} Montrer que $f'(x)=\dfrac{2x^2-8x+4}{(x-2)^2}$, puis dresser le tableau de variations de $f$ sur $]2\,;+\infty[$ (les extrema sont atteints en $2+\sqrt{2}$ ; on donnera la valeur exacte du minimum).

\textbf{8.} Tracer les asymptotes, le centre de symétrie $\Omega$ et la courbe $\mathcal{C}_f$.

\textbf{9.} La société souhaite que le temps moyen de parcours reste inférieur à $15$ minutes. Déterminer graphiquement, puis par le calcul (on résoudra une inéquation du second degré), les valeurs entières de $x$ convenables.
\end{exercice}

\newpage

\coursec{Corrigés des exercices}

\coursec{Exercices corrigés : Courbes représentatives}

\courssub{Exercices d'application directe : définitions et premiers calculs}

\begin{corrige}[Exercice 1 — QCM : une seule bonne réponse]
\textbf{1. Réponse \cle{b)} $x=2$.} La valeur interdite est $x=2$ ; or $\lim\limits_{x\to 2}(3x+1)=7\neq 0$, donc $\lim\limits_{x\to 2^-}f(x)=-\infty$ et $\lim\limits_{x\to 2^+}f(x)=+\infty$ : la droite $x=2$ est \cle{asymptote verticale}. (La droite $y=3$ est, elle, \cle{asymptote horizontale}.)

\textbf{2. Réponse \cle{b)}.} Par définition, si $\lim\limits_{x\to+\infty}f(x)=5$ (limite finie en l'infini), la droite d'équation $y=5$ est \cle{asymptote horizontale} à $\mathcal{C}_f$ en $+\infty$.

\textbf{3. Réponse \cle{b)}.} C'est la caractérisation du \cle{centre de symétrie} : $\Omega(a\,;b)$ est centre de symétrie si et seulement si $f(a+x)+f(a-x)=2b$ pour tout $x$ convenable. La condition a) correspond à un \cle{axe de symétrie} $x=a$.

\textbf{4. Réponse \cle{c)}.} Par définition, $y=2x+1$ est \cle{asymptote oblique} en $+\infty$ si et seulement si $\lim\limits_{x\to+\infty}\bigl[f(x)-(2x+1)\bigr]=0$. La condition b) $\bigl(\lim[f(x)-2x]=1\bigr)$ en est une reformulation \emph{équivalente} (car $f(x)-(2x+1)=[f(x)-2x]-1$), mais c) est la formulation \emph{canonique} du cours, celle qu'il faut retenir et utiliser pour justifier.
\end{corrige}

\begin{corrige}[Exercice 2 — Vrai ou faux, en justifiant]
\textbf{1. \cle{FAUX}.} Une valeur interdite ne donne pas nécessairement une limite infinie. Contre-exemple : $f(x)=\dfrac{x^2-1}{x-1}$ définie sur $\mathbb{R}\setminus\{1\}$. Pour $x\neq 1$, $f(x)=x+1$, donc $\lim\limits_{x\to 1}f(x)=2$ (finie) : la droite $x=1$ n'est pas asymptote, la courbe présente seulement un « trou ». Il faut \emph{calculer la limite} pour conclure.

\textbf{2. \cle{FAUX}.} Une courbe \emph{peut} couper son asymptote horizontale. Contre-exemple classique : $f(x)=\dfrac{\sin x}{x}$ admet $y=0$ pour asymptote horizontale en $+\infty$, et sa courbe coupe cette asymptote une infinité de fois (en $x=k\pi$, $k\in\mathbb{Z}^*$). Exemple de fonction rationnelle : $f(x)=\dfrac{x^2-x}{x^2+1}$ a pour asymptote horizontale $y=1$ (car $\lim_{x\to\pm\infty}f(x)=1$) et la coupe en $x=-1$ car $f(-1)=1$.

\textbf{3. \cle{VRAI}.} $\Omega(1\,;2)$ centre de symétrie signifie, avec $a=1$ et $b=2$ : $f(1+h)+f(1-h)=2b=4$ pour tout $h$ tel que $1+h$ et $1-h$ appartiennent au domaine. C'est exactement la définition.

\textbf{4. \cle{VRAI}.} Les comportements en $+\infty$ et en $-\infty$ sont indépendants. Exemple : $f(x)=\dfrac{x+|x|}{2}+\dfrac{1}{x^2+1}$. Pour $x>0$, $f(x)=x+\dfrac{1}{x^2+1}$ (asymptote oblique $y=x$ en $+\infty$) ; pour $x<0$, $f(x)=\dfrac{1}{x^2+1}\to 0$ (asymptote horizontale $y=0$ en $-\infty$). C'est donc possible.
\end{corrige}

\begin{corrige}[Exercice 3 — Compléter les définitions]
\textbf{1.} La droite d'équation $x=a$ est \cle{asymptote verticale} à la courbe de $f$ si $\lim\limits_{x\to a}f(x)=+\infty$ ou $-\infty$ (ou l'une des limites unilatérales est infinie).

\textbf{2.} La droite d'équation $y=b$ est \cle{asymptote horizontale} à la courbe de $f$ en $+\infty$ si $\lim\limits_{x\to+\infty}f(x)=b$.

\textbf{3.} La droite d'équation $x=a$ est \cle{axe de symétrie} de la courbe de $f$ si, pour tout $h$ convenable : $f(a+h)=f(a-h)$.

\textbf{4.} Pour montrer que la droite d'équation $y=ax+b$ est \cle{asymptote oblique} à la courbe de $f$ en $+\infty$, on calcule la limite de $f(x)-(ax+b)$ et on vérifie qu'elle vaut $\mathbf{0}$.
\end{corrige}

\begin{corrige}[Exercice 4 — Calculs directs de limites]
\textbf{1.} $f(x)=\dfrac{2x-3}{x+1}$.
\begin{itemize}
\item En $-1$ : $\lim\limits_{x\to -1}(2x-3)=-5\neq 0$ et $x+1\to 0$. Pour $x\to -1^+$, $x+1>0$, donc $\lim\limits_{x\to -1^+}f(x)=-\infty$ ; pour $x\to -1^-$, $x+1<0$, donc $\lim\limits_{x\to -1^-}f(x)=+\infty$. La droite $\boxed{x=-1}$ est \cle{asymptote verticale}.
\item En $+\infty$ : $f(x)=\dfrac{x(2-\frac3x)}{x(1+\frac1x)}=\dfrac{2-\frac3x}{1+\frac1x}\to 2$. La droite $\boxed{y=2}$ est \cle{asymptote horizontale} en $+\infty$ (et de même en $-\infty$).
\end{itemize}

\textbf{2.} $f(x)-(x+1)=\dfrac{2}{x}$, donc $\lim\limits_{x\to+\infty}\bigl[f(x)-(x+1)\bigr]=\lim\limits_{x\to+\infty}\dfrac{2}{x}=0$. La droite $\boxed{y=x+1}$ est \cle{asymptote oblique} en $+\infty$ (et aussi en $-\infty$).

\textbf{3.} $f(x)=\dfrac{x^2+1}{x}$. Numérateur : $x^2+1\to 1>0$. Pour $x\to 0^+$, $x>0$, donc $\lim\limits_{x\to 0^+}f(x)=+\infty$ ; pour $x\to 0^-$, $x<0$, donc $\lim\limits_{x\to 0^-}f(x)=-\infty$. La droite $\boxed{x=0}$ (l'axe des ordonnées) est \cle{asymptote verticale}.
\end{corrige}

\courssub{Axe et centre de symétrie}

\begin{corrige}[Exercice 5 — Axe de symétrie d'une parabole]
\textbf{1.} Calculons les deux membres pour tout réel $h$ :
\begin{align*}
f(-1+h)&=(-1+h)^2+2(-1+h)-3=1-2h+h^2-2+2h-3=h^2-4\\
f(-1-h)&=(-1-h)^2+2(-1-h)-3=1+2h+h^2-2-2h-3=h^2-4
\end{align*}
Donc pour tout réel $h$ : $f(-1+h)=f(-1-h)$.

\textbf{2.} On en déduit que la droite d'équation $\boxed{x=-1}$ est \cle{axe de symétrie} de $\mathcal{C}_f$.

\textbf{3.} $f'(x)=2x+2=2(x+1)$ : $f'(x)<0$ pour $x<-1$, $f'(x)>0$ pour $x>-1$. Minimum : $f(-1)=1-2-3=-4$. Limites : $\lim\limits_{\pm\infty}f=+\infty$.

\begin{center}
\begin{tabular}{|c|ccccc|}\hline
$x$ & $-\infty$ & & $-1$ & & $+\infty$ \\ \hline
$f'(x)$ & & $-$ & $0$ & $+$ & \\ \hline
$f$ & $+\infty$ & $\searrow$ & $-4$ & $\nearrow$ & $+\infty$ \\ \hline
\end{tabular}
\end{center}

Points de tracé : $f(0)=-3$, $f(1)=0$, $f(-3)=0$, $f(2)=5$, $f(-4)=5$. On trace la parabole de sommet $S(-1\,;-4)$, passant par $(1\,;0)$ et $(-3\,;0)$, et l'axe $x=-1$ en pointillés. Les points $(1\,;0)$ et $(-3\,;0)$, $(2\,;5)$ et $(-4\,;5)$ sont symétriques par rapport à cet axe, ce qui confirme le résultat.
\end{corrige}

\begin{corrige}[Exercice 6 — Centre de symétrie et réduction du domaine d'étude]
\textbf{1.} $f$ est définie si et seulement si $x-2\neq 0$, donc $D_f=\mathbb{R}\setminus\{2\}=]-\infty\,;2[\cup]2\,;+\infty[$. Ce domaine est symétrique par rapport à $2$ : si $2+h\in D_f$ alors $h\neq 0$ et $2-h\in D_f$.

\textbf{2.} Pour $h\neq 0$ :
\[ f(2+h)=\frac{(2+h)-1}{(2+h)-2}=\frac{1+h}{h}, \qquad f(2-h)=\frac{(2-h)-1}{(2-h)-2}=\frac{1-h}{-h}=\frac{h-1}{h}. \]
Donc :
\[ f(2+h)+f(2-h)=\frac{1+h}{h}+\frac{h-1}{h}=\frac{2h}{h}=2. \]

\textbf{3.} On a montré $f(2+h)+f(2-h)=2=2\times 1$ pour tout $h\neq 0$. Avec $a=2$ et $b=1$, c'est la caractérisation du \cle{centre de symétrie} : le point $\boxed{\Omega(2\,;1)}$ est centre de symétrie de $\mathcal{C}_f$.

\textbf{4.} Si $M(2+h\,;f(2+h))$ est un point de la courbe avec $h>0$, alors son symétrique par rapport à $\Omega$, qui est $M'(2-h\,;2-f(2+h))=M'(2-h\,;f(2-h))$, appartient aussi à la courbe. La partie de la courbe sur $]-\infty\,;2[$ est donc l'image de celle sur $]2\,;+\infty[$ par la \cle{symétrie centrale de centre $\Omega$}. Il suffit d'étudier $f$ sur $]2\,;+\infty[$, puis de compléter le tracé par cette symétrie : le domaine d'étude est réduit de moitié.
\end{corrige}

\courssub{Asymptotes verticales, horizontales et obliques}

\begin{corrige}[Exercice 7 — Asymptotes d'une fonction homographique]
\textbf{1.} Pour $x\to 0^+$ : $4x+8\to 8>0$ et $x\to 0^+$, donc $\lim\limits_{x\to 0^+}f(x)=+\infty$. Interprétation graphique : la droite $\boxed{x=0}$ (axe des ordonnées) est \cle{asymptote verticale} à la courbe. Concrètement : quand la quantité stockée devient très faible, le coût unitaire explose.

\textbf{2.} $f(x)=\dfrac{4x+8}{x}=4+\dfrac{8}{x}$, donc $\lim\limits_{x\to+\infty}f(x)=4$. La droite $\boxed{y=4}$ est \cle{asymptote horizontale} en $+\infty$. Interprétation : quand la quantité stockée devient très grande, le coût unitaire se rapproche de $4$ milliers de FCFA par $\mathrm{m}^3$ (coût plancher), sans jamais l'atteindre.

\textbf{3.} $\Omega(0\,;4)$ n'appartient pas à la courbe car $0\notin D_f$ (la fonction n'est pas définie en $0$).

Pour le centre de symétrie, prolongeons $f$ sur $\mathbb{R}\setminus\{0\}$. Pour tout $h\neq 0$ :
\[ f(h)+f(-h)=\left(4+\frac{8}{h}\right)+\left(4-\frac{8}{h}\right)=8=2\times 4. \]
Donc $f(0+h)+f(0-h)=2\times 4$ : le point $\boxed{\Omega(0\,;4)}$ est \cle{centre de symétrie} de $\mathcal{C}_f$. Remarque : $\Omega$ est l'intersection des deux asymptotes $x=0$ et $y=4$, ce qui est général pour les fonctions homographiques.
\end{corrige}

\courssub{Lecture graphique et conjectures}

\begin{corrige}[Exercice 8 — Lecture graphique et conjectures]
\textbf{1.} Conjectures à partir du graphique :
\begin{itemize}
\item La courbe semble se rapprocher de la droite $\boxed{y=2}$ en $+\infty$ et en $-\infty$ : \cle{asymptote horizontale} d'équation $y=2$.
\item La courbe semble symétrique par rapport à l'axe des ordonnées : \cle{axe de symétrie} $\boxed{x=0}$.
\item La courbe semble située \textbf{au-dessus} de son asymptote $y=2$ (elle s'en rapproche par valeurs supérieures).
\end{itemize}

\textbf{2.a)} $g(x)=\dfrac{2x^2+3}{x^2+1}=\dfrac{x^2\left(2+\frac{3}{x^2}\right)}{x^2\left(1+\frac{1}{x^2}\right)}=\dfrac{2+\frac{3}{x^2}}{1+\frac{1}{x^2}}$, donc $\lim\limits_{x\to+\infty}g(x)=2$. La conjecture est \textbf{validée} : $y=2$ est asymptote horizontale en $+\infty$ (et le même calcul donne la limite $2$ en $-\infty$).

\textbf{2.b)} Le domaine $\mathbb{R}$ est symétrique par rapport à $0$ et, pour tout réel $x$ :
\[ g(-x)=\frac{2(-x)^2+3}{(-x)^2+1}=\frac{2x^2+3}{x^2+1}=g(x). \]
Donc $g$ est \cle{paire} : sa courbe admet l'axe des ordonnées ($x=0$) pour axe de symétrie. Conjecture validée.

\textbf{2.c)} $g(x)-2=\dfrac{2x^2+3}{x^2+1}-2=\dfrac{2x^2+3-2x^2-2}{x^2+1}=\dfrac{1}{x^2+1}$. Or $x^2+1>0$ pour tout $x$, donc $g(x)-2>0$, c'est-à-dire $g(x)>2$ pour tout réel $x$ : la courbe est \textbf{strictement au-dessus} de son asymptote $y=2$. Conjecture validée. (De plus $g(x)-2\to 0$, ce qui redémontre l'asymptote.)
\end{corrige}

\courssub{Synthèse : études complètes type Bac}

\begin{corrige}[Exercice 9 — Type Bac : étude complète avec asymptote oblique]
\textbf{Barème indicatif (sur 10 pts) :} Partie A : 2 pts ; Partie B : 3 pts ; Partie C : 5 pts (dont 1,5 pt pour le tracé).

\textbf{Partie A.}

\textbf{1.} $f$ est définie ssi $x-1\neq 0$ : $D_f=\mathbb{R}\setminus\{1\}$.

\textbf{2.} Numérateur en $1$ : $1-3+4=2>0$. Dénominateur : $x-1\to 0$, positif si $x>1$, négatif si $x<1$. Donc :
\[ \lim_{x\to 1^+}f(x)=+\infty \qquad \text{et} \qquad \lim_{x\to 1^-}f(x)=-\infty. \]
La droite $\mathcal{D}_1$ d'équation $\boxed{x=1}$ est \cle{asymptote verticale} à $\mathcal{C}_f$.

\textbf{3.} $f(x)=\dfrac{x^2\left(1-\frac3x+\frac{4}{x^2}\right)}{x\left(1-\frac1x\right)}=x\cdot\dfrac{1-\frac3x+\frac{4}{x^2}}{1-\frac1x}$. La fraction tend vers $1$, donc $\lim\limits_{x\to+\infty}f(x)=+\infty$ et $\lim\limits_{x\to-\infty}f(x)=-\infty$ (car $x\to-\infty$ multiplié par un facteur tendant vers $1$).

\textbf{Partie B.}

\textbf{4.} Pour $x\neq 1$ :
\[ x-2+\frac{2}{x-1}=\frac{(x-2)(x-1)+2}{x-1}=\frac{x^2-3x+2+2}{x-1}=\frac{x^2-3x+4}{x-1}=f(x). \]

\textbf{5.} $f(x)-(x-2)=\dfrac{2}{x-1}$, donc $\lim\limits_{x\to+\infty}\bigl[f(x)-(x-2)\bigr]=0$ et de même en $-\infty$. La droite $\mathcal{D}_2$ : $\boxed{y=x-2}$ est \cle{asymptote oblique} à $\mathcal{C}_f$ en $+\infty$ et en $-\infty$.

\textbf{6.} Le signe de $f(x)-(x-2)=\dfrac{2}{x-1}$ est celui de $x-1$ :
\begin{itemize}
\item si $x>1$ : $f(x)-(x-2)>0$, $\mathcal{C}_f$ est \textbf{au-dessus} de $\mathcal{D}_2$ ;
\item si $x<1$ : $f(x)-(x-2)<0$, $\mathcal{C}_f$ est \textbf{en dessous} de $\mathcal{D}_2$.
\end{itemize}

\textbf{Partie C.}

\textbf{7.} $f=\dfrac{u}{v}$ avec $u=x^2-3x+4$, $v=x-1$, $u'=2x-3$, $v'=1$ :
\[ f'(x)=\frac{(2x-3)(x-1)-(x^2-3x+4)}{(x-1)^2}=\frac{2x^2-5x+3-x^2+3x-4}{(x-1)^2}=\frac{x^2-2x-1}{(x-1)^2}. \]

\textbf{8.} $(x-1)^2>0$ sur $D_f$, donc $f'(x)$ a le signe de $x^2-2x-1$. Discriminant : $\Delta=4+4=8$, racines $x_1=\dfrac{2-\sqrt{8}}{2}=1-\sqrt{2}$ et $x_2=1+\sqrt{2}$. Le trinôme est positif à l'extérieur des racines, négatif entre elles.

Valeurs exactes : $f(1-\sqrt{2})$ : avec $x-1=-\sqrt{2}$, $f(x)=x-2+\dfrac{2}{x-1}=-1-\sqrt{2}-\dfrac{2}{\sqrt{2}}=-1-\sqrt{2}-\sqrt{2}=-1-2\sqrt{2}$. De même $f(1+\sqrt{2})=-1+\sqrt{2}+\sqrt{2}=-1+2\sqrt{2}$.

\emph{Tableau de variations sur $]-\infty\,;1[$ :}
\begin{center}
\begin{tabular}{|c|ccccc|}\hline
$x$ & $-\infty$ & & $1-\sqrt{2}$ & & $1$ \\ \hline
$f'(x)$ & & $+$ & $0$ & $-$ & \\ \hline
$f$ & $-\infty$ & $\nearrow$ & $-1-2\sqrt{2}$ & $\searrow$ & $-\infty$ \\ \hline
\end{tabular}
\end{center}

\emph{Tableau de variations sur $]1\,;+\infty[$ :}
\begin{center}
\begin{tabular}{|c|ccccc|}\hline
$x$ & $1$ & & $1+\sqrt{2}$ & & $+\infty$ \\ \hline
$f'(x)$ & & $-$ & $0$ & $+$ & \\ \hline
$f$ & $+\infty$ & $\searrow$ & $-1+2\sqrt{2}$ & $\nearrow$ & $+\infty$ \\ \hline
\end{tabular}
\end{center}

\textbf{9.} Tracé : asymptote verticale $x=1$, asymptote oblique $y=x-2$ (passant par $(0\,;-2)$ et $(2\,;0)$). Branche de gauche : au-dessous de $\mathcal{D}_2$, maximum en $(1-\sqrt{2}\,;-1-2\sqrt{2})\approx(-0{,}41\,;-3{,}83)$, plongeant vers $-\infty$ le long de $x=1$. Branche de droite : part de $+\infty$ le long de $x=1$, minimum en $(1+\sqrt{2}\,;-1+2\sqrt{2})\approx(2{,}41\,;1{,}83)$, remonte en restant au-dessus de $\mathcal{D}_2$.

\textbf{10.} \emph{Graphiquement} : la droite $y=3$ coupe la branche de droite en deux points (car $3>-1+2\sqrt{2}\approx 1{,}83$, valeur du minimum) et ne coupe pas la branche de gauche (dont le maximum vaut $\approx -3{,}83<3$) : \textbf{2 solutions}.

\emph{Algébriquement} : $f(x)=3\iff x^2-3x+4=3(x-1)\iff x^2-6x+7=0$ (avec $x\neq 1$). $\Delta=36-28=8$, solutions $x=\dfrac{6\pm 2\sqrt{2}}{2}=3\pm\sqrt{2}$, toutes deux différentes de $1$. Donc $S=\{3-\sqrt{2}\,;\,3+\sqrt{2}\}$ : bien 2 solutions.

\textbf{Points de vigilance :} citer la condition $x\neq 1$ avant tout calcul ; justifier l'asymptote oblique par la limite de la \emph{différence} et non par « $f(x)\approx x-2$ » ; pour la position relative, étudier le \emph{signe} de la différence ; dans le tableau, interdire à $f$ d'être définie en $1$ (double barre ou deux tableaux séparés).
\end{corrige}

\begin{corrige}[Exercice 10 — Type Bac : parité, asymptotes et équation]
\textbf{Barème indicatif (sur 10 pts) :} Partie A : 1 pt ; Partie B : 2,5 pts ; Partie C : 3,5 pts ; Partie D : 3 pts.

\textbf{Partie A.}

\textbf{1.} Le domaine $\mathbb{R}\setminus\{0\}$ est symétrique par rapport à $0$. Pour tout $x\neq 0$ :
\[ f(-x)=-x+\frac{4}{-x}=-x-\frac{4}{x}=-\left(x+\frac{4}{x}\right)=-f(x). \]
Donc $f$ est \cle{impaire} : $\mathcal{C}_f$ admet l'origine $O$ pour \cle{centre de symétrie}. La partie de la courbe sur $]-\infty\,;0[$ se déduit de celle sur $]0\,;+\infty[$ par la symétrie de centre $O$ : on peut restreindre l'étude à $]0\,;+\infty[$.

\textbf{Partie B.}

\textbf{2.} Pour $x\to 0^+$ : $x\to 0$ et $\dfrac{4}{x}\to +\infty$, donc $\lim\limits_{x\to 0^+}f(x)=+\infty$. La droite $\boxed{x=0}$ (axe des ordonnées) est \cle{asymptote verticale}.

\textbf{3.} $\lim\limits_{x\to+\infty}x=+\infty$ et $\lim\limits_{x\to+\infty}\dfrac4x=0$, donc $\lim\limits_{x\to+\infty}f(x)=+\infty$.

\textbf{4.} $f(x)-x=\dfrac{4}{x}$, donc $\lim\limits_{x\to+\infty}\bigl[f(x)-x\bigr]=0$ : la droite $\mathcal{D}$ d'équation $\boxed{y=x}$ est \cle{asymptote oblique} en $+\infty$ (et en $-\infty$ par imparité).

\textbf{5.} Sur $]0\,;+\infty[$ : $f(x)-x=\dfrac{4}{x}>0$, donc $\mathcal{C}_f$ est strictement \textbf{au-dessus} de $\mathcal{D}$.

\textbf{Partie C.}

\textbf{6.} $f'(x)=1-\dfrac{4}{x^2}=\dfrac{x^2-4}{x^2}$.

\textbf{7.} Sur $]0\,;+\infty[$ : $x^2>0$, donc $f'(x)$ a le signe de $x^2-4=(x-2)(x+2)$, c'est-à-dire négatif sur $]0\,;2[$ et positif sur $]2\,;+\infty[$. Minimum : $f(2)=2+2=4$.

\begin{center}
\begin{tabular}{|c|ccccc|}\hline
$x$ & $0$ & & $2$ & & $+\infty$ \\ \hline
$f'(x)$ & & $-$ & $0$ & $+$ & \\ \hline
$f$ & $+\infty$ & $\searrow$ & $4$ & $\nearrow$ & $+\infty$ \\ \hline
\end{tabular}
\end{center}

Par imparité (symétrie de centre $O$), sur $]-\infty\,;0[$ : $f$ est croissante sur $]-\infty\,;-2]$ de $-\infty$ à $f(-2)=-4$, puis décroissante sur $[-2\,;0[$ de $-4$ à $-\infty$.

\textbf{8.} Tracé : asymptote verticale $x=0$, asymptote oblique $y=x$. Sur $]0\,;+\infty[$ : la courbe descend de $+\infty$ (le long de l'axe des ordonnées) jusqu'au minimum $(2\,;4)$, puis remonte en restant au-dessus de $y=x$ dont elle se rapproche. On obtient la partie sur $]-\infty\,;0[$ par symétrie centrale de centre $O$ : maximum en $(-2\,;-4)$, courbe en dessous de $y=x$.

\textbf{Partie D.}

\textbf{9.} Graphiquement, la droite $y=5$ coupe la branche de droite en deux points (car $5>4$, le minimum) et ne coupe pas la branche de gauche (située sous $y=x<0$ pour $x<0$... plus précisément $f(x)<0$ pour $x<0$ donc $f(x)=5$ n'a pas de solution négative) : \textbf{2 solutions}, lues approximativement $x\approx 1$ et $x\approx 4$.

\textbf{10.} $f(x)=5\iff x+\dfrac4x=5\iff x^2-5x+4=0$ (en multipliant par $x\neq 0$). $\Delta=25-16=9$, solutions $x=\dfrac{5\pm 3}{2}$, soit $x=1$ ou $x=4$. Donc $S=\{1\,;\,4\}$, ce qui confirme la lecture graphique.

\textbf{11.} Pour $x>0$ : $f(x)\leq 5\iff x^2-5x+4\leq 0$ (car $x>0$, multiplier par $x$ conserve le sens). Le trinôme s'annule en $1$ et $4$ et est négatif entre ses racines : $x\in[1\,;4]$. La grandeur est rentable pour une valeur de $x$ comprise entre $1$ et $4$.

\textbf{Points de vigilance :} pour l'imparité, vérifier d'abord la symétrie du domaine ; la multiplication par $x$ dans une inéquation exige de discuter le signe de $x$ (ici $x>0$, donc le sens est conservé) ; ne pas oublier que la symétrie de centre $O$ transforme un minimum en maximum.
\end{corrige}

\begin{corrige}[Exercice 11 — Type Bac : fonction rationnelle et centre de symétrie]
\textbf{Barème indicatif (sur 10 pts) :} Partie A : 3 pts ; Partie B : 2,5 pts ; Partie C : 4,5 pts.

\textbf{Partie A.}

\textbf{1.} $f$ est définie ssi $x\neq 2$ : $D_f=\mathbb{R}\setminus\{2\}$. Numérateur en $2$ : $2(4)-5(2)+6=8-10+6=4>0$. Dénominateur $x-2\to 0$, positif à droite, négatif à gauche :
\[ \lim_{x\to 2^+}f(x)=+\infty, \qquad \lim_{x\to 2^-}f(x)=-\infty. \]
Interprétation graphique : la droite $\boxed{x=2}$ est \cle{asymptote verticale} à $\mathcal{C}_f$.

\textbf{2.} Pour $x\neq 2$ :
\[ 2x-1+\frac{4}{x-2}=\frac{(2x-1)(x-2)+4}{x-2}=\frac{2x^2-4x-x+2+4}{x-2}=\frac{2x^2-5x+6}{x-2}=f(x). \]

\textbf{3.} $f(x)-(2x-1)=\dfrac{4}{x-2}$, donc $\lim\limits_{x\to+\infty}\bigl[f(x)-(2x-1)\bigr]=0$ et de même $\lim\limits_{x\to-\infty}\bigl[f(x)-(2x-1)\bigr]=0$. La droite $\Delta$ : $\boxed{y=2x-1}$ est \cle{asymptote oblique} en $+\infty$ et en $-\infty$.

Position relative : le signe de $f(x)-(2x-1)=\dfrac{4}{x-2}$ est celui de $x-2$ :
\begin{itemize}
\item si $x>2$ : $\mathcal{C}_f$ est \textbf{au-dessus} de $\Delta$ ;
\item si $x<2$ : $\mathcal{C}_f$ est \textbf{en dessous} de $\Delta$.
\end{itemize}

\textbf{Partie B.}

\textbf{4.} $\Omega$ est l'intersection de $x=2$ et $y=2x-1$ : $y=2(2)-1=3$, donc $\boxed{\Omega(2\,;3)}$.

\textbf{5.} Pour $h\neq 0$, en utilisant la forme de la question 2 :
\begin{align*}
f(2+h)&=2(2+h)-1+\frac{4}{h}=3+2h+\frac{4}{h}\\
f(2-h)&=2(2-h)-1+\frac{4}{-h}=3-2h-\frac{4}{h}
\end{align*}
Donc $f(2+h)+f(2-h)=6$.

\textbf{6.} On a $f(2+h)+f(2-h)=6=2\times 3$ pour tout $h\neq 0$ : c'est la caractérisation du \cle{centre de symétrie} avec $a=2$, $b=3$. Donc $\boxed{\Omega(2\,;3)}$ est centre de symétrie de $\mathcal{C}_f$. Conséquence : il suffit d'étudier $f$ sur $]2\,;+\infty[$ (qui est d'ailleurs le domaine concret du modèle) ; la partie sur $]-\infty\,;2[$ s'obtient par la symétrie de centre $\Omega$.

\textbf{Partie C.}

\textbf{7.} Dérivons $f(x)=2x-1+\dfrac{4}{x-2}$ : $f'(x)=2-\dfrac{4}{(x-2)^2}=\dfrac{2(x-2)^2-4}{(x-2)^2}=\dfrac{2(x^2-4x+4)-4}{(x-2)^2}=\dfrac{2x^2-8x+4}{(x-2)^2}$.

Sur $]2\,;+\infty[$, $(x-2)^2>0$, donc $f'(x)$ a le signe de $2x^2-8x+4=2(x^2-4x+2)$. Racines : $\Delta=16-8=8$, $x=\dfrac{4\pm 2\sqrt{2}}{2}=2\pm\sqrt{2}$. Seule $2+\sqrt{2}$ appartient à $]2\,;+\infty[$. Le trinôme est négatif sur $]2\,;2+\sqrt{2}[$ et positif ensuite.

Minimum : $f(2+\sqrt{2})=3+2\sqrt{2}+\dfrac{4}{\sqrt{2}}=3+2\sqrt{2}+2\sqrt{2}=\boxed{3+4\sqrt{2}}\approx 8{,}66$.

\begin{center}
\begin{tabular}{|c|ccccc|}\hline
$x$ & $2$ & & $2+\sqrt{2}$ & & $+\infty$ \\ \hline
$f'(x)$ & & $-$ & $0$ & $+$ & \\ \hline
$f$ & $+\infty$ & $\searrow$ & $3+4\sqrt{2}$ & $\nearrow$ & $+\infty$ \\ \hline
\end{tabular}
\end{center}

\textbf{8.} Tracé : asymptotes $x=2$ (verticale) et $y=2x-1$ (oblique, passant par $(0\,;-1)$ et $(2\,;3)$), centre $\Omega(2\,;3)$ placé à leur intersection. Sur $]2\,;+\infty[$ : la courbe descend de $+\infty$ le long de $x=2$ jusqu'au minimum $(2+\sqrt{2}\,;3+4\sqrt{2})\approx(3{,}41\,;8{,}66)$, puis remonte en restant au-dessus de $\Delta$. La branche sur $]-\infty\,;2[$ se déduit par la symétrie de centre $\Omega$ (maximum en $(2-\sqrt{2}\,;3-4\sqrt{2})$).

\textbf{9.} \emph{Lecture graphique} : la droite $y=15$ coupe la branche de droite en deux points d'abscisses approximatives $x\approx 2{,}3$ et $x\approx 6{,}7$ ; le temps est inférieur à $15$ minutes entre ces deux valeurs.

\emph{Par le calcul} : pour $x>2$ :
\[ f(x)\leq 15\iff 2x^2-5x+6\leq 15(x-2)\iff 2x^2-20x+36\leq 0\iff x^2-10x+18\leq 0 \]
(la multiplication par $x-2>0$ conserve le sens). $\Delta=100-72=28$, racines $x=\dfrac{10\pm 2\sqrt{7}}{2}=5\pm\sqrt{7}$, soit $5-\sqrt{7}\approx 2{,}35$ et $5+\sqrt{7}\approx 7{,}65$. Le trinôme est négatif entre ses racines : $x\in[5-\sqrt{7}\,;\,5+\sqrt{7}]$.

Les valeurs \textbf{entières} convenables sont donc $x\in\{3\,;\,4\,;\,5\,;\,6\,;\,7\}$ : la société doit faire circuler entre 3 et 7 bus pour que le temps moyen reste inférieur à 15 minutes.

\textbf{Points de vigilance :} l'asymptote oblique se justifie par $\lim[f(x)-(2x-1)]=0$, pas par la seule écriture de $f$ ; le centre de symétrie exige l'égalité $f(a+h)+f(a-h)=2b$ démontrée pour \emph{tout} $h\neq 0$ ; dans l'inéquation, justifier le sens de l'inégalité par $x-2>0$ ; ne pas oublier de ne retenir que les solutions \emph{entières} dans la réponse finale en contexte.
\end{corrige}

\newpage

\coursec{Fiche bilan}

\coursec{Courbes représentatives}

\courssub{Lecture graphique et conjectures}

\begin{definition}[Lecture graphique]
Soit $f$ une fonction définie sur un ensemble $D_f \subset \mathbb{R}$ et $\mathcal{C}_f$ sa courbe représentative dans un repère orthonormé $(O;\vec{\imath},\vec{\jmath})$. La lecture graphique permet de \emph{conjecturer} (supposer) des propriétés de $f$ :
\begin{itemize}
\item les \cle{variations} (croissance/décroissance) par la direction de la courbe ;
\item le \cle{signe} de $f(x)$ par la position par rapport à l'axe des abscisses (au-dessus $\Leftrightarrow f(x)>0$, en-dessous $\Leftrightarrow f(x)<0$) ;
\item les \cle{limites} aux bornes de $D_f$ et en $\pm\infty$ par le comportement des branches ;
\item d'éventuelles \cle{symétries} (axe vertical $x=a$, centre $\Omega(a;b)$) ;
\item d'éventuelles \cle{asymptotes} (verticales, horizontales, obliques).
\end{itemize}
\end{definition}

\begin{attention}[Conjecture $\neq$ Preuve]
Une lecture graphique, même précise, ne constitue \textbf{jamais} une démonstration. Au Baccalauréat, toute conjecture lue sur un graphique doit être \textbf{vérifiée par le calcul algébrique} (limites, dérivées, égalités fonctionnelles). Une affirmation non justifiée vaut 0 point.
\end{attention}

\begin{methode}[Émettre et vérifier une conjecture]
\begin{enumerate}
\item \textbf{Observer} la courbe : repérer les points remarquables, le sens de parcours, les directions asymptotiques, les symétries visuelles.
\item \textbf{Formuler} la conjecture par écrit (ex: « La courbe semble admettre $x=2$ comme axe de symétrie »).
\item \textbf{Vérifier} par le calcul :
\begin{itemize}
\item Symétrie axe $x=a$ : montrer $f(a+h)=f(a-h)$ pour tout $h$ tel que $a\pm h\in D_f$.
\item Symétrie centre $\Omega(a;b)$ : montrer $f(a+h)+f(a-h)=2b$.
\item Asymptote verticale $x=a$ : calculer $\lim_{x\to a^{\pm}}f(x)=\pm\infty$.
\item Asymptote horizontale $y=b$ : calculer $\lim_{x\to\pm\infty}f(x)=b$.
\item Asymptote oblique $y=ax+b$ : déterminer $a=\lim_{x\to\pm\infty}\frac{f(x)}{x}$ et $b=\lim_{x\to\pm\infty}(f(x)-ax)$, puis vérifier $\lim_{x\to\pm\infty}(f(x)-(ax+b))=0$.
\end{itemize}
\item \textbf{Conclure} par une phrase affirmative ou négative.
\end{enumerate}
\end{methode}

\begin{exemplebox}[Conjecture sur une courbe donnée]
On donne la courbe $\mathcal{C}_f$ ci-dessous. Conjecturer puis vérifier la nature de la symétrie et l'asymptote horizontale.
\begin{popfigure}
\begin{tikzpicture}[scale=0.8, domain=-4:4]
\draw[popline, thin, ->] (-4.5,0) -- (4.5,0) node[right] {$x$};
\draw[popline, thin, ->] (0,-1.5) -- (0,3.5) node[above] {$y$};
\foreach \x in {-4,-3,-2,-1,1,2,3,4} \draw (\x,0.05) -- (\x,-0.05) node[below] {\small $\x$};
\foreach \y in {-1,1,2,3} \draw (0.05,\y) -- (-0.05,\y) node[left] {\small $\y$};
% Courbe y = 1/(x^2+1) + 1  -> symétrie x=0, asymptote y=1
\draw[popcurve, thick, samples=200, domain=-4:4] plot (\x, {1/(\x*\x+1)+1});
\draw[popmuted, dashed] (-4.5,1) -- (4.5,1) node[right] {$y=1$};
\draw[popmuted, dashed] (0,-1.5) -- (0,3.5);
\node[popink] at (3.5, 2.5) {$\mathcal{C}_f$};
\node[popmuted] at (-4, 1.3) {Asymptote ?};
\end{tikzpicture}
\legende{Courbe suggérant une symétrie par rapport à l'axe des ordonnées et une asymptote horizontale $y=1$.}
\end{popfigure}
\textbf{Conjectures :} La courbe est symétrique par rapport à l'axe des ordonnées ($x=0$). Elle admet la droite $y=1$ comme asymptote horizontale en $+\infty$ et $-\infty$.
\textbf{Vérification :} On suppose $f(x)=\frac{1}{x^2+1}+1$ (déduite de la forme).
\begin{itemize}
\item Symétrie : $f(-x)=\frac{1}{(-x)^2+1}+1=\frac{1}{x^2+1}+1=f(x)$. Donc $x=0$ est axe de symétrie.
\item Asymptote horizontale : $\lim_{x\to\pm\infty}f(x)=\lim_{x\to\pm\infty}\left(\frac{1}{x^2+1}+1\right)=0+1=1$. Donc $y=1$ est asymptote horizontale.
\end{itemize}
\end{exemplebox}

\courssub{Axe de symétrie d'une courbe}

\begin{definition}[Axe de symétrie]
Soit $f$ une fonction définie sur $D_f$. La droite d'équation $x=a$ est un \cle{axe de symétrie} de la courbe $\mathcal{C}_f$ si et seulement si :
\[
\forall h \in \mathbb{R},\quad a+h \in D_f \implies a-h \in D_f \quad\text{et}\quad f(a+h)=f(a-h).
\]
Équivalentement : $\forall x \in D_f,\quad 2a-x \in D_f \quad\text{et}\quad f(2a-x)=f(x)$.
\end{definition}

\begin{propriete}[Caractérisation par la fonction paire]
La courbe $\mathcal{C}_f$ admet $x=a$ comme axe de symétrie $\iff$ la fonction $g$ définie par $g(x)=f(x+a)$ est \textbf{paire} sur $D_g = \{x \mid x+a \in D_f\}$.
\emph{Preuve :} $g(-x)=f(-x+a)=f(a-x)=f(a+x)=g(x)$.
\end{propriete}

\begin{methode}[Démontrer qu'une droite $x=a$ est axe de symétrie]
\begin{enumerate}
\item Vérifier la condition sur le domaine : $\forall x \in D_f,\ 2a-x \in D_f$.
\item Calculer $f(a+h)$ et $f(a-h)$ (ou $f(2a-x)$ et $f(x)$).
\item Montrer l'égalité $f(a+h)=f(a-h)$.
\item Conclure : « Donc la droite d'équation $x=a$ est un axe de symétrie de $\mathcal{C}_f$. »
\end{enumerate}
\end{methode}

\begin{exemplebox}[Axe de symétrie : $f(x)=\frac{2x-1}{x-2}$]
Déterminer l'axe de symétrie de $\mathcal{C}_f$ pour $f(x)=\frac{2x-1}{x-2}$ ($D_f=\mathbb{R}^*_2$).
\textbf{Recherche :} On cherche $a$ tel que $f(a+h)=f(a-h)$.
$f(a+h)=\frac{2(a+h)-1}{a+h-2}=\frac{2a+2h-1}{a+h-2}$,
$f(a-h)=\frac{2a-2h-1}{a-h-2}$.
On veut $\frac{2a+2h-1}{a+h-2}=\frac{2a-2h-1}{a-h-2}$.
Produit en croix : $(2a+2h-1)(a-h-2)=(2a-2h-1)(a+h-2)$.
Développement et simplification conduisent à $a=2$.
\textbf{Vérification :} $D_f=\mathbb{R}\setminus\{2\}$. Si $x\neq 2$, $4-x\neq 2$. $f(4-x)=\frac{2(4-x)-1}{4-x-2}=\frac{7-2x}{2-x}=\frac{2x-7}{x-2}$.
$f(x)=\frac{2x-1}{x-2}$. Ne sont pas égales ! \textbf{Attention :} Cette fonction admet un \emph{centre} de symétrie, pas un axe. (Voir section centre).
\end{exemplebox}

\begin{exemplebox}[Axe de symétrie : $f(x)=x^2-4x+5$]
$f(x)=(x-2)^2+1$. Posons $g(x)=f(x+2)=x^2+1$ (paire). Donc $x=2$ est axe de symétrie.
Vérification : $f(2+h)=(h)^2+1$, $f(2-h)=(-h)^2+1=h^2+1$. Égalité vérifiée.
\end{exemplebox}

\courssub{Centre de symétrie d'une courbe}

\begin{definition}[Centre de symétrie]
Soit $f$ une fonction définie sur $D_f$. Le point $\Omega(a\,;\,b)$ est un \cle{centre de symétrie} de la courbe $\mathcal{C}_f$ si et seulement si :
\[
\forall h \in \mathbb{R},\quad a+h \in D_f \implies a-h \in D_f \quad\text{et}\quad f(a+h)+f(a-h)=2b.
\]
Équivalentement : $\forall x \in D_f,\quad 2a-x \in D_f \quad\text{et}\quad f(2a-x)=2b-f(x)$.
\end{definition}

\begin{propriete}[Caractérisation par la fonction impaire]
La courbe $\mathcal{C}_f$ admet $\Omega(a;b)$ comme centre de symétrie $\iff$ la fonction $g$ définie par $g(x)=f(x+a)-b$ est \textbf{impaire} sur $D_g = \{x \mid x+a \in D_f\}$.
\emph{Preuve :} $g(-x)=f(-x+a)-b=f(a-x)-b$. Comme $f(a-x)+f(a+x)=2b \iff f(a-x)-b = -(f(a+x)-b) = -g(x)$.
\end{propriete}

\begin{methode}[Démontrer qu'un point $\Omega(a;b)$ est centre de symétrie]
\begin{enumerate}
\item Vérifier la condition sur le domaine : $\forall x \in D_f,\ 2a-x \in D_f$.
\item Calculer la somme $f(a+h)+f(a-h)$ (ou $f(x)+f(2a-x)$).
\item Montrer que cette somme vaut $2b$ (constante, indépendante de $h$ ou $x$).
\item Conclure : « Donc le point $\Omega(a\,;\,b)$ est un centre de symétrie de $\mathcal{C}_f$. »
\end{enumerate}
\end{methode}

\begin{attention}[Piège fréquent : Somme vs Différence]
Pour un \textbf{axe} $x=a$ : on compare $f(a+h)$ et $f(a-h)$ $\rightarrow$ \textbf{Égalité} ($=$).
Pour un \textbf{centre} $\Omega(a;b)$ : on fait la \textbf{Somme} $f(a+h)+f(a-h)$ $\rightarrow$ Résultat $2b$.
Ne pas confondre : $f(a+h)-f(a-h)=0$ (axe) vs $f(a+h)+f(a-h)=2b$ (centre).
\end{attention}

\begin{exemplebox}[Centre de symétrie : $f(x)=\frac{2x-1}{x-2}$]
$D_f=\mathbb{R}\setminus\{2\}$. On cherche $\Omega(a;b)$.
Méthode classique : $f(x)=\frac{2x-1}{x-2}=\frac{2(x-2)+3}{x-2}=2+\frac{3}{x-2}$.
La courbe est l'image de $y=\frac{3}{x}$ (centre $O(0;0)$) par translation $\vec{v}(2;2)$.
Donc centre $\Omega(2;2)$.
\textbf{Vérification formelle :} $a=2, b=2$.
$f(2+h)=2+\frac{3}{h}$, $f(2-h)=2-\frac{3}{h}$.
Somme : $f(2+h)+f(2-h)=4=2b$. Domaine ok ($h\neq 0 \iff 2\pm h \neq 2$).
Donc $\Omega(2;2)$ est centre de symétrie.
\end{exemplebox}

\begin{exemplebox}[Centre de symétrie : Fonction polynôme degré 3]
Soit $f(x)=x^3-3x^2+2$. $f'(x)=3x^2-6x=3x(x-2)$. Point d'inflexion à $x=1$ ($f''(x)=6x-6=0 \Rightarrow x=1$).
$f(1)=0$. Conjecture : $\Omega(1;0)$.
Vérification : $f(1+h)=(1+h)^3-3(1+h)^2+2 = 1+3h+3h^2+h^3 -3(1+2h+h^2)+2 = h^3 -3h$.
$f(1-h)= -h^3 +3h$.
Somme : $f(1+h)+f(1-h)=0=2b \Rightarrow b=0$. CQFD.
\end{exemplebox}

\courssub{Asymptotes verticales et horizontales}

\begin{definition}[Asymptote verticale]
Soit $a \in \mathbb{R}$ (ou borne de $D_f$). La droite d'équation $x=a$ est une \cle{asymptote verticale} de $\mathcal{C}_f$ si et seulement si :
\[
\lim_{x\to a^+}f(x)=+\infty \text{ ou } -\infty \quad\text{OU}\quad \lim_{x\to a^-}f(x)=+\infty \text{ ou } -\infty.
\]
(Il suffit qu'une seule des limites latérales soit infinie).
\end{definition}

\begin{definition}[Asymptote horizontale]
La droite d'équation $y=b$ ($b \in \mathbb{R}$) est une \cle{asymptote horizontale} de $\mathcal{C}_f$ si et seulement si :
\[
\lim_{x\to +\infty}f(x)=b \quad\text{OU}\quad \lim_{x\to -\infty}f(x)=b.
\]
(Si les deux limites valent $b$, c'est la même asymptote pour les deux branches).
\end{definition}

\begin{methode}[Recherche des asymptotes verticales et horizontales]
\begin{enumerate}
\item Déterminer $D_f$ et identifier les valeurs interdites $a$ (zéros de dénominateur, bornes de racine, etc.).
\item Pour chaque $a$ « suspect » : calculer $\lim_{x\to a^-}f(x)$ et $\lim_{x\to a^+}f(x)$. Si l'une est infinie $\Rightarrow$ asymptote verticale $x=a$.
\item Calculer $\lim_{x\to +\infty}f(x)$ et $\lim_{x\to -\infty}f(x)$. Si limite finie $b \Rightarrow$ asymptote horizontale $y=b$.
\end{enumerate}
\end{methode}

\begin{exemplebox}[Asymptotes verticales et horizontales]
$f(x)=\frac{x^2+1}{x-1}$, $D_f=\mathbb{R}\setminus\{1\}$.
\begin{itemize}
\item \textbf{Verticale :} $x=1$. $\lim_{x\to 1^-}f(x)=-\infty$, $\lim_{x\to 1^+}f(x)=+\infty$. Asymptote verticale $x=1$.
\item \textbf{Horizontale :} $\lim_{x\to\pm\infty}f(x)=\lim_{x\to\pm\infty}\frac{x^2}{x}=\pm\infty$. Pas d'asymptote horizontale (mais oblique, voir section suivante).
\end{itemize}
\end{exemplebox}

\begin{savaistu}[Asymptote verticale et domaine]
Une asymptote verticale $x=a$ correspond \textbf{toujours} à une valeur $a$ exclue du domaine $D_f$ (ou borne de $D_f$). Réciproquement, une valeur exclue ne donne pas toujours une asymptote verticale (ex: $f(x)=\frac{\sin x}{x}$ en $0$, limite finie $1$ : trou, pas asymptote).
\end{savaistu}

\courssub{Asymptotes obliques}

\begin{definition}[Asymptote oblique]
La droite d'équation $y=ax+b$ ($a \neq 0$) est une \cle{asymptote oblique} de $\mathcal{C}_f$ en $+\infty$ (resp. $-\infty$) si et seulement si :
\[
\lim_{x\to +\infty}\big[f(x)-(ax+b)\big]=0 \quad\left(\text{resp. } \lim_{x\to -\infty}\big[f(x)-(ax+b)\big]=0\right).
\]
\end{definition}

\begin{propriete}[Détermination des coefficients $a$ et $b$]
Si $\mathcal{C}_f$ admet une asymptote oblique $y=ax+b$ en $+\infty$, alors nécessairement :
\[
a = \lim_{x\to +\infty}\frac{f(x)}{x} \quad\text{et}\quad b = \lim_{x\to +\infty}\big(f(x)-ax\big).
\]
Mêmes formules en $-\infty$. \textbf{Attention :} Il faut vérifier que $a \neq 0$ (sinon c'est une horizontale) et que les deux limites existent et sont finies.
\end{propriete}

\begin{methode}[Recherche d'une asymptote oblique (algorithme standard)]
\begin{enumerate}
\item Calculer $a = \lim_{x\to \pm\infty}\frac{f(x)}{x}$. Si $a=0$ ou n'existe pas (infini) $\rightarrow$ Pas d'asymptote oblique (chercher horizontale).
\item Calculer $b = \lim_{x\to \pm\infty}\big(f(x)-ax\big)$.
\item \textbf{Vérifier} que $\lim_{x\to \pm\infty}\big[f(x)-(ax+b)\big]=0$ (souvent automatique si étapes 1 et 2 ok, mais exigé au Bac).
\item Écrire l'équation $y=ax+b$.
\item \textbf{Étudier la position relative} : signe de $f(x)-(ax+b)$ pour savoir si la courbe est au-dessus ($>0$) ou en-dessous ($<0$) de l'asymptote.
\end{enumerate}
\end{methode}

\begin{attention}[Erreurs classiques asymptote oblique]
\begin{itemize}
\item Oublier de vérifier $a \neq 0$.
\item Calculer $b = \lim f(x)$ au lieu de $\lim (f(x)-ax)$.
\item Ne pas traiter $+\infty$ et $-\infty$ séparément (les asymptotes peuvent être différentes !).
\item Ne pas étudier la position relative (souvent demandé : « Préciser la position de $\mathcal{C}_f$ par rapport à son asymptote »).
\end{itemize}
\end{attention}

\begin{exemplebox}[Asymptote oblique : $f(x)=\frac{x^2+1}{x-1}$]
$D_f=\mathbb{R}\setminus\{1\}$. Asymptote verticale $x=1$.
\textbf{En $+\infty$ :}
$a = \lim_{x\to+\infty}\frac{x^2+1}{x(x-1)} = \lim\frac{x^2}{x^2}=1$.
$b = \lim_{x\to+\infty}\left(\frac{x^2+1}{x-1}-x\right) = \lim\frac{x^2+1-x(x-1)}{x-1} = \lim\frac{x+1}{x-1}=1$.
Asymptote : $y=x+1$.
Vérification : $f(x)-(x+1)=\frac{x^2+1-(x+1)(x-1)}{x-1}=\frac{x^2+1-(x^2-1)}{x-1}=\frac{2}{x-1}$.
$\lim_{x\to+\infty}\frac{2}{x-1}=0$. Position : $\frac{2}{x-1}>0$ pour $x>1$ $\Rightarrow$ Courbe \textbf{au-dessus} de l'asymptote.
\textbf{En $-\infty$ :} Même calcul, $a=1, b=1$. Asymptote $y=x+1$.
Position : $x<1 \Rightarrow x-1<0 \Rightarrow \frac{2}{x-1}<0$ $\Rightarrow$ Courbe \textbf{en-dessous}.
\end{exemplebox}

\begin{exemplebox}[Asymptote oblique : Racine carrée $f(x)=\sqrt{x^2+x}$]
$D_f=(-\infty,-1]\cup[0,+\infty)$.
\textbf{En $+\infty$ :} $a=\lim\frac{\sqrt{x^2+x}}{x}=\lim\sqrt{1+\frac{1}{x}}=1$.
$b=\lim(\sqrt{x^2+x}-x)=\lim\frac{x}{\sqrt{x^2+x}+x}=\lim\frac{1}{\sqrt{1+1/x}+1}=\frac{1}{2}$.
Asymptote $y=x+\frac{1}{2}$. Position : $f(x)-(x+1/2)=\frac{-1/4}{\sqrt{x^2+x}+x+1/2}<0$ $\Rightarrow$ Courbe \textbf{en-dessous}.
\textbf{En $-\infty$ :} $x<0$, $\sqrt{x^2}=-x$. $a=\lim\frac{\sqrt{x^2+x}}{x}=\lim\frac{-x\sqrt{1+1/x}}{x}=-1$.
$b=\lim(\sqrt{x^2+x}+x)=\lim\frac{x}{\sqrt{x^2+x}-x}=\lim\frac{1}{\sqrt{1+1/x}-1}$ (forme $\frac{1}{0^-}$) $\to -\infty$.
Pas d'asymptote oblique en $-\infty$ (branche parabolique).
\end{exemplebox}

\courssub{Synthèse : construire une courbe complète}

\begin{methode}[Tracer $\mathcal{C}_f$ : La check-list complète (ordre recommandé)]
\begin{enumerate}
\item \textbf{Domaine $D_f$} : Ensemble de définition. Noter les valeurs interdites.
\item \textbf{Symétries} : Chercher axe $x=a$ ou centre $\Omega(a;b)$. Si trouvés, \textbf{réduire l'étude} à une moitié du domaine (ex: $x \ge a$ pour axe, ou $x \ge a$ pour centre).
\item \textbf{Limites et Asymptotes} :
\begin{itemize}
\item Aux bornes de $D_f$ (asymptotes verticales).
\item En $\pm\infty$ (asymptotes horizontales/obliques + position relative).
\end{itemize}
\item \textbf{Dérivée $f'$ et Tableau de variations} : Calculer $f'(x)$, signe, variations. Reporter les limites aux bornes du tableau.
\item \textbf{Points remarquables} : Intersections avec axes ($f(x)=0$, $f(0)$), points d'inflexion ($f''(x)=0$ et changement de signe), extremums.
\item \textbf{Tracé} :
\begin{itemize}
\item Repère orthonormé (échelle adaptée).
\item Tracer les asymptotes en \textbf{pointillés} et les nommer ($x=a$, $y=ax+b$).
\item Placer les points remarquables.
\item Tracer la courbe en respectant variations, convexité (si étudiée), position par rapport aux asymptotes.
\item Utiliser la symétrie pour compléter l'autre moitié.
\end{itemize}
\end{enumerate}
\end{methode}

\begin{exoresolu}[Étude complète : $f(x)=\frac{x^2-2x-3}{x+1}$]
\textbf{1. Domaine :} $D_f=\mathbb{R}\setminus\{-1\}$.
\textbf{2. Simplification :} $f(x)=\frac{(x-3)(x+1)}{x+1}=x-3$ pour $x\neq -1$.
C'est une droite $y=x-3$ \textbf{avec un trou} en $x=-1$ (point $A(-1;-4)$ exclu).
\textbf{3. Symétries :} Pas de symétrie axe/centre évidente sur la forme simplifiée (droite non passant par l'origine, pas parallèle aux axes).
\textbf{4. Limites / Asymptotes :}
$\lim_{x\to -1}f(x)=-4$ (finie) $\Rightarrow$ \textbf{Pas d'asymptote verticale} (trou).
$+\infty$ : $f(x)\sim x \Rightarrow$ Asymptote oblique $y=x-3$ (c'est la droite elle-même !).
Vérification : $f(x)-(x-3)=0$ pour $x\neq -1$. Position : confondue (sauf au trou).
\textbf{5. Variations :} $f'(x)=1>0$. Strictement croissante sur $]-\infty;-1[$ et $]-1;+\infty[$.
\textbf{6. Tracé :} Droite $y=x-3$ en trait plein, point $A(-1;-4)$ en cercle vide. Asymptote oblique = la droite (trait pointillé superposé ou non tracé si confondu, mais mentionner « asymptote oblique confondue avec la courbe hors point $A$ »).
\end{exoresolu}

\begin{exoresolu}[Étude complète : $f(x)=x+\frac{4}{x}$ (Contexte : Coût moyen)]
\textit{Une entreprise de télécom (type MTN Cameroon) a un coût total $C(x)=x^2+4$ (en millions FCFA) pour produire $x$ milliers de cartes SIM. Le coût moyen est $f(x)=\frac{C(x)}{x}=x+\frac{4}{x}$ pour $x>0$. Étudier $f$ et tracer $\mathcal{C}_f$.}
\textbf{1. Domaine :} $D_f=]0;+\infty[$ (contexte physique) ou $\mathbb{R}^*$ (maths pures). On prend $\mathbb{R}^*$.
\textbf{2. Symétries :} $f(-x)=-x-\frac{4}{x}=-f(x)$. Fonction \textbf{impaire} $\Rightarrow$ Centre de symétrie $\Omega(0;0)$.
$\rightarrow$ Étudier sur $]0;+\infty[$, déduire $]-\infty;0[$ par rotation $180^\circ$.
\textbf{3. Limites / Asymptotes ($x>0$) :}
$\lim_{x\to 0^+}f(x)=+\infty \Rightarrow$ Asymptote verticale $x=0$ (axe des ordonnées).
$+\infty$ : $a=\lim\frac{f(x)}{x}=\lim(1+\frac{4}{x^2})=1$.
$b=\lim(f(x)-x)=\lim\frac{4}{x}=0$.
Asymptote oblique $y=x$.
Position : $f(x)-x=\frac{4}{x}>0$ sur $]0;+\infty[$ $\Rightarrow$ Courbe \textbf{au-dessus} de $y=x$.
\textbf{4. Variations ($x>0$) :} $f'(x)=1-\frac{4}{x^2}=\frac{x^2-4}{x^2}=\frac{(x-2)(x+2)}{x^2}$.
Signe : $f'<0$ sur $]0;2[$, $f'=0$ en $2$, $f'>0$ sur $]2;+\infty[$.
Minimum $f(2)=4$.
\textbf{5. Points :} $A(2;4)$ minimum. Pas d'intersection axe $x$ ($x^2+4=0$ impossible).
\textbf{6. Tracé :} Branche $x>0$ : part de $+\infty$ (asymptote $x=0$), descend jusqu'à $A(2;4)$, remonte vers $+\infty$ au-dessus de $y=x$. Branche $x<0$ : symétrique par rapport à $O$.
\end{exoresolu}

\begin{popfigure}
\begin{tikzpicture}[mindmap, grow cyclic,
every node/.style={concept, text=white, font=\small\bfseries},
root concept/.append style={concept color=popink, font=\bfseries},
level 1/.append style={level distance=3.6cm, sibling angle=72},
level 2/.append style={level distance=2.2cm, font=\scriptsize},
concept connection/.append style={color=popmuted}]
\node[root concept]{Courbes représentatives}
child[concept color=popblue]{ node{Lecture graphique}
  child[concept color=popblue!70]{ node{Conjecturer puis vérifier} }
  child[concept color=popblue!70]{ node{Variations, signe, limites} }
}
child[concept color=popgreen]{ node{Axe de symétrie $x=a$}
  child[concept color=popgreen!70]{ node{$f(a+h)=f(a-h)$} }
  child[concept color=popgreen!70]{ node{Domaine symétrique} }
}
child[concept color=poporange]{ node{Centre de symétrie $\Omega(a\,;\,b)$}
  child[concept color=poporange!70]{ node{$f(a+h)+f(a-h)=2b$} }
  child[concept color=poporange!70]{ node{Somme = 2b} }
}
child[concept color=poppurple]{ node{Asymptotes V \& H}
  child[concept color=poppurple!70]{ node{Vert. $x=a$: $\lim=\pm\infty$} }
  child[concept color=poppurple!70]{ node{Horiz. $y=b$: $\lim_{\pm\infty}=b$} }
}
child[concept color=popgold]{ node{Asymptote oblique $y=ax+b$}
  child[concept color=popgold!70]{ node{$a=\lim f/x$} }
  child[concept color=popgold!70]{ node{$b=\lim f-ax$} }
  child[concept color=popgold!70]{ node{Position : signe $f-(ax+b)$} }
}
child[concept color=poppink]{ node{Tracer une courbe}
  child[concept color=poppink!70]{ node{Domaine $\to$ Symétries} }
  child[concept color=poppink!70]{ node{Limites $\to$ Asymptotes} }
  child[concept color=poppink!70]{ node{Derivée $\to$ Tab. Var.} }
  child[concept color=poppink!70]{ node{Points rem. $\to$ Tracé} }
};
\end{tikzpicture}
\legende{Carte mentale de la leçon : les six compétences clés autour des courbes représentatives.}
\end{popfigure}

\begin{propriete}[Synthèse des définitions et formules clés]
\begin{center}
\begin{tabularx}{\linewidth}{|l|X|}
\hline
\rowcolor{popblueL} \textbf{Objet} & \textbf{Condition nécessaire et suffisante (à vérifier)} \\
\hline
Axe de symétrie $x=a$ & $\forall x\in D_f,\ 2a-x\in D_f$ et $f(2a-x)=f(x)$ \quad (ou $f(a+h)=f(a-h)$) \\
\hline
Centre de symétrie $\Omega(a\,;\,b)$ & $\forall x\in D_f,\ 2a-x\in D_f$ et $f(2a-x)=2b-f(x)$ \quad (ou $f(a+h)+f(a-h)=2b$) \\
\hline
Asymptote verticale $x=a$ & $\lim_{x\to a^-}f(x)=\pm\infty$ \textbf{ou} $\lim_{x\to a^+}f(x)=\pm\infty$ ($a$ borne de $D_f$ ou valeur interdite) \\
\hline
Asymptote horizontale $y=b$ & $\lim_{x\to +\infty}f(x)=b$ \textbf{ou} $\lim_{x\to -\infty}f(x)=b$ ($b\in\mathbb{R}$) \\
\hline
Asymptote oblique $y=ax+b$ ($a\neq0$) & $a=\lim_{x\to\pm\infty}\frac{f(x)}{x} \in \mathbb{R}^*$,\quad $b=\lim_{x\to\pm\infty}(f(x)-ax) \in \mathbb{R}$,\quad \textbf{et} $\lim_{x\to\pm\infty}[f(x)-(ax+b)]=0$ \\
\hline
Position courbe / Asymptote oblique & Signe de $\Delta(x)=f(x)-(ax+b)$ : $\Delta(x)>0 \Rightarrow \mathcal{C}_f$ \textbf{au-dessus} ; $\Delta(x)<0 \Rightarrow \mathcal{C}_f$ \textbf{en-dessous} \\
\hline
\end{tabularx}
\end{center}
\end{propriete}

\begin{aretenir}[Checklist avant le Bac — Points de vigilance]
Avant de rendre ta copie, vérifie chaque point :
\begin{itemize}
\item[$\square$] J'ai précisé l'\textbf{ensemble de définition $D_f$} avant tout calcul de limite.
\item[$\square$] Pour un axe $x=a$, j'ai bien vérifié que $2a-x\in D_f$ dès que $x\in D_f$ (domaine symétrique).
\item[$\square$] Pour un centre $\Omega(a\,;\,b)$, j'ai calculé la \textbf{somme} $f(a+h)+f(a-h)$ (ou $f(x)+f(2a-x)$) et non la différence.
\item[$\square$] J'ai conclu chaque symétrie par une \textbf{phrase complète} : « Donc la droite d'équation $x=a$ est axe de symétrie de $\mathcal{C}_f$ » ou « Donc le point $\Omega(a\,;\,b)$ est centre de symétrie de $\mathcal{C}_f$ ».
\item[$\square$] Une asymptote verticale n'est annoncée qu'après avoir obtenu une limite \textbf{infinie} en un point $a$ (préciser $a^+$ ou $a^-$).
\item[$\square$] Une asymptote horizontale n'existe qu'en $+\infty$ ou $-\infty$, jamais en un point fini.
\item[$\square$] Pour l'asymptote oblique, j'ai bien calculé $a=\lim\frac{f(x)}{x}$ (vérifié $a\neq 0$), puis $b=\lim(f(x)-ax)$, \textbf{puis vérifié} $\lim[f(x)-(ax+b)]=0$.
\item[$\square$] J'ai étudié le \textbf{signe de $f(x)-(ax+b)$} pour préciser la position relative (au-dessus / en-dessous / coupure).
\item[$\square$] Mon tableau de variation est cohérent avec les limites calculées (valeurs aux bornes, $\pm\infty$ aux asymptotes verticales).
\item[$\square$] J'ai utilisé les \textbf{symétries} pour limiter l'étude à une moitié du domaine (gain de temps, rigueur).
\item[$\square$] Toute conjecture lue sur un graphique a été \textbf{démontrée par le calcul} : une lecture graphique seule ne vaut pas preuve.
\item[$\square$] Mes asymptotes sont tracées en \textbf{pointillés} et \textbf{nommées} sur la figure ($x=a$, $y=ax+b$).
\item[$\square$] Les points exclus du domaine (trous) sont représentés par des \textbf{cercles vides}.
\end{itemize}
\end{aretenir}

\begin{savaistu}[Le saviez-vous ? — Asymptotes et ingénierie au Cameroun]
Dans la conception des réseaux de fibre optique (CAMTEL, MTN, Orange), les ingénieurs modélisent l'atténuation du signal par des fonctions rationnelles. Les \textbf{asymptotes verticales} correspondent à des fréquences de résonance où le signal devient instable (gain infini théorique). Les \textbf{asymptotes horizontales} modélisent le bruit de fond thermique (limite du signal à très haute fréquence). Maîtriser ces courbes, c'est assurer la qualité de votre connexion 4G/5G à Yaoundé ou Douala !
\end{savaistu}

\begin{exercice}[Entraînement Bac - Série A]
\begin{enumerate}
\item Soit $f(x)=\frac{x^2+2x+3}{x+1}$ sur $D_f=\mathbb{R}\setminus\{-1\}$.
\begin{enumerate}
\item Déterminer les asymptotes de $\mathcal{C}_f$.
\item Étudier la position de $\mathcal{C}_f$ par rapport à son asymptote oblique.
\item Montrer que $\mathcal{C}_f$ admet un centre de symétrie. Donner ses coordonnées.
\item Tracer $\mathcal{C}_f$ et ses asymptotes.
\end{enumerate}
\item Soit $g(x)=x^3-3x$ sur $\mathbb{R}$.
\begin{enumerate}
\item Étudier les symétries de $\mathcal{C}_g$.
\item Déterminer les asymptotes.
\item Tracer $\mathcal{C}_g$.
\end{enumerate}
\end{enumerate}
\end{exercice}

\begin{corrige}[Entraînement Bac — question 1]
\begin{enumerate}
\item \textbf{Asymptotes :}
\begin{itemize}
\item Verticale : $x=-1$. $\lim_{x\to -1^-}f(x)=-\infty$, $\lim_{x\to -1^+}f(x)=+\infty$. Asymptote $x=-1$.
\item Oblique : $a=\lim_{x\to\pm\infty}\frac{x^2+2x+3}{x(x+1)}=1$.
$b=\lim_{x\to\pm\infty}\left(\frac{x^2+2x+3}{x+1}-x\right)=\lim\frac{x^2+2x+3-x^2-x}{x+1}=\lim\frac{x+3}{x+1}=1$.
Asymptote $y=x+1$.
Vérification : $f(x)-(x+1)=\frac{x^2+2x+3-(x+1)^2}{x+1}=\frac{2}{x+1}\to 0$. OK.
\end{itemize}
\item \textbf{Position :} $\Delta(x)=\frac{2}{x+1}$.
$x>-1 \Rightarrow \Delta>0$ (au-dessus). $x<-1 \Rightarrow \Delta<0$ (en-dessous).
\item \textbf{Centre de symétrie :} Intersection asymptotes : $x=-1 \Rightarrow y=0$. $\Omega(-1;0)$.
Vérification : $f(-1+h)=\frac{(-1+h)^2+2(-1+h)+3}{h}=\frac{h^2+2}{h}$ et $f(-1-h)=\frac{h^2+2}{-h}$, donc $f(-1+h)+f(-1-h)=0=2\times 0$ : $\Omega(-1;0)$ est bien centre de symétrie.
\item \textbf{Tracé :} tracer d'abord en pointillés les asymptotes $x=-1$ et $y=x+1$, placer $\Omega(-1;0)$, puis les deux branches : au-dessus de l'asymptote oblique pour $x>-1$, en dessous pour $x<-1$, symétriques par rapport à $\Omega$.
\end{enumerate}
\end{corrige}

\begin{corrige}[Entraînement Bac — question 2]
\begin{enumerate}
\item $g(-x)=-g(x)$ : fonction impaire $\Rightarrow$ Centre de symétrie $O(0;0)$.
\item $\lim_{x\to\pm\infty}g(x)=\pm\infty$. $a=\lim g/x = \lim x^2-3 = +\infty$. Pas d'asymptote oblique ni horizontale.
\item Variations : $g'(x)=3x^2-3=3(x-1)(x+1)$.
Max $g(-1)=2$, Min $g(1)=-2$. Tracé cubique standard symétrique par rapport à $O$.
\end{enumerate}
\end{corrige}

\findecours

\end{document}
